\documentclass[12pt,a4paper]{article}
\usepackage[utf8]{inputenc}
\usepackage[T1]{fontenc}
\usepackage[french]{babel}
\usepackage{amsmath,amssymb,amsthm,amsfonts}
\usepackage{geometry}
\geometry{margin=1in}
\usepackage{hyperref}
\usepackage{listings}

\newtheorem{theorem}{Théorème}
\newtheorem{lemma}{Lemme}
\newtheorem{definition}{Définition}
\newtheorem{conjecture}{Conjecture}

\title{Analyse et Preuves Partielles sur la Conjecture du Tournesol d'Erdős-Rado}
\author{}
\date{}

\begin{document}
\maketitle

\section{Analyse et Décomposition}

\subsection{Définitions Axiomatiques}

Soit $\mathcal{U}$ un univers fini et soit $\mathcal{P}(\mathcal{U})$ l'ensemble des parties de $\mathcal{U}$. Nous considérons une famille $\mathcal{F} \subseteq \mathcal{P}(\mathcal{U})$.

\begin{definition}
Une famille d'ensembles $\mathcal{S} = \{S_1, S_2, \dots, S_r\} \subseteq \mathcal{F}$ est un \emph{tournesol} (ou système-$\Delta$) de taille $r$ s'il existe un ensemble $C \subseteq \mathcal{U}$, appelé le \emph{cœur}, tel que pour tout $i \neq j \in \{1, \dots, r\}$, on a $S_i \cap S_j = C$. Les ensembles $S_i \setminus C$ sont appelés les pétales du tournesol, et ils sont mutuellement disjoints.
\end{definition}

Le typage des éléments structuraux s'établit ainsi de manière rigoureuse :
\begin{itemize}
    \item L'univers $\mathcal{U} : \mathrm{Type}$.
    \item La famille $\mathcal{F} : \mathcal{P}(\mathcal{P}(\mathcal{U}))$.
    \item L'entier $k \in \mathbb{N}$, qui majore le cardinal des éléments : $\forall A \in \mathcal{F}, |A| \le k$.
    \item L'entier $r \in \mathbb{N}$ avec $r \ge 3$, spécifiant la taille du tournesol recherché.
\end{itemize}

\begin{conjecture}[Conjecture du Tournesol d'Erdős-Rado, 1960]
Pour tout entier $r \ge 3$, il existe une constante réelle $C(r) > 0$ telle que pour tout entier $k \in \mathbb{N}$ et toute famille $\mathcal{F}$ d'ensembles de cardinal au plus $k$, si $|\mathcal{F}| > C(r)^k$, alors $\mathcal{F}$ contient un tournesol de taille $r$.
\end{conjecture}

\subsection{Structures Sous-jacentes}

Le problème se modélise via la théorie des hypergraphes uniformes. L'hypergraphe $\mathcal{H} = (V, E)$ correspond à la famille $\mathcal{F}$, où les sommets sont les éléments de $\cup \mathcal{F}$ et les arêtes sont les ensembles $A \in \mathcal{F}$. L'absence de système-$\Delta$ impose des contraintes sévères sur les degrés des hyperarêtes et sur les intersections admissibles.
La méthode repose sur l'identification d'ensembles de forte dispersion (spreadness) et l'utilisation de l'entropie de Shannon ou de marches aléatoires avec mémorisation pour borner le nombre de configurations sans sous-structure régulière.

\section{Recherche de Littérature Contextuelle}

Le théorème classique d'Erdős-Rado donne une borne supérieure de $k!(r-1)^k$. Ce résultat découle d'une simple induction sur $k$. En fixant un ensemble $A$, s'il croise tout autre ensemble, l'un des éléments de $A$ a un degré très élevé, ce qui permet de descendre à $k-1$.
Récemment, Alweiss, Lovett, Wu, et Zhang ont bouleversé le domaine en démontrant une borne de l'ordre de $(\log k)^k (r \log r)^k$, en utilisant l'inégalité de concentration de type de l'inégalité de sous-additivité de l'entropie.

Cette percée présente une forte analogie avec le théorème de Dvir sur les extracteurs de Kakeya dans les corps finis, où le passage de méthodes combinatoires pures à des méthodes polynomiales (ou ici, probabilistes basées sur l'encodage de variables aléatoires) permet de dépasser les barrières d'induction classique.

\section{Stratégie de Preuve \& Isolation de Lemmes}

Nous structurons l'approche partielle autour de trois lemmes permettant de manipuler des hypergraphes dispersés.

\begin{itemize}
    \item \textbf{Lemme 1 (Indépendance par sélection aléatoire)} : Démontre qu'une famille aléatoirement sous-échantillonnée conserve une grande proportion de composantes disjointes avec forte probabilité. La démonstration s'effectue par une analyse minutieuse des bornes d'union (Union Bound) et de l'inégalité de Markov.
    \item \textbf{Lemme 2 (Borne sur l'intersection des hyperarêtes)} : Formalise une borne stricte sur l'espérance de l'intersection de sous-ensembles aléatoires, en utilisant la méthode des moments d'ordre supérieur.
    \item \textbf{Lemme 3 (Extraction de cœur)} : Montre comment, à partir d'un ensemble de forte probabilité, on peut déduire de manière déterministe un cœur $C$. La démonstration procède par descente d'induction sur la taille du cœur.
\end{itemize}

\section{Rédaction de la Preuve Informelle}

\subsection{Démonstration du Lemme 1}

Nous développons l'évaluation explicite de l'indépendance par sélection sur un hypergraphe $\mathcal{H} = (V, \mathcal{F})$. Soit $\mathcal{F}$ de cardinalité $m$. Nous sélectionnons chaque sous-ensemble $A \in \mathcal{F}$ avec une probabilité $p \in (0, 1)$.

La probabilité qu'au moins deux ensembles $A, B \in \mathcal{F}$ avec $A \neq B$ s'intersectent s'évalue par le principe d'inclusion-exclusion. Nous déroulons l'expansion algébrique complète des événements d'intersection pour quantifier l'erreur de dispersion.

Considérons l'expansion de l'espérance du nombre de $t$-uplets d'ensembles intersectant un domaine cible $X$.

\subsubsection{Analyse du moment d'ordre 2}Pour le degré 2, nous évaluons l'intersection simultanée de $\ell = 2$ hyperarêtes. Soient $E_1, \dots, E_{2}$ des arêtes de l'hypergraphe.
\begin{equation}\mathbb{E}\left[ \prod_{i=1}^{2} \mathbf{1}_{E_i \cap X \neq \emptyset} \right] = \sum_{j=1}^{2} (-1)^{j-1} \sum_{1 \le i_1 < \dots < i_j \le 2} \mathbb{P}\left( \bigcap_{m=1}^j (E_{i_m} \cap X \neq \emptyset) \right)\end{equation}Nous décomposons la probabilité conjointe en utilisant l'indépendance conditionnelle sur le choix des sommets de $X$. L'ensemble des configurations de recouvrement de l'union $\bigcup_{i} E_i$ est partitionné selon la taille de leur intersection avec $X$.
\begin{align} \\\Delta_{2, 1} &= \sum_{|S| = 1} \mathbb{P}(S \subseteq X) \sum_{I \subseteq \{1, \dots, 2\}, |I| \ge 1} (-1)^{|I|+1} \mathbf{1}_{S \subseteq \cup_{i \in I} E_i}
\Delta_{2, 2} &= \sum_{|S| = 2} \mathbb{P}(S \subseteq X) \sum_{I \subseteq \{1, \dots, 2\}, |I| \ge 1} (-1)^{|I|+1} \mathbf{1}_{S \subseteq \cup_{i \in I} E_i}\end{align}Cette majoration permet de borner supérieurement la corrélation locale. Par application du lemme de Lovász local ou des inégalités de concentration azumaiennes, la déviation par rapport à la moyenne produit un terme résiduel qui s'amortit exponentiellement. Nous fixons explicitement cette décroissance :
\begin{equation}\mathcal{R}_{2} \le \exp\left( - \frac{ \left( \Delta_{2, 1} \right)^2 }{ 2 \sum_{j=2}^{2} \Delta_{2, j} + \frac{1}{3} \max_{j} \Delta_{2, j} } \right)\end{equation}Le contrôle de ce terme $\mathcal{R}_{2}$ certifie l'indépendance quasi-totale des choix pour les structures de dimension 2.

\subsubsection{Analyse du moment d'ordre 3}Pour le degré 3, nous évaluons l'intersection simultanée de $\ell = 3$ hyperarêtes. Soient $E_1, \dots, E_{3}$ des arêtes de l'hypergraphe.
\begin{equation}\mathbb{E}\left[ \prod_{i=1}^{3} \mathbf{1}_{E_i \cap X \neq \emptyset} \right] = \sum_{j=1}^{3} (-1)^{j-1} \sum_{1 \le i_1 < \dots < i_j \le 3} \mathbb{P}\left( \bigcap_{m=1}^j (E_{i_m} \cap X \neq \emptyset) \right)\end{equation}Nous décomposons la probabilité conjointe en utilisant l'indépendance conditionnelle sur le choix des sommets de $X$. L'ensemble des configurations de recouvrement de l'union $\bigcup_{i} E_i$ est partitionné selon la taille de leur intersection avec $X$.
\begin{align} \\\Delta_{3, 1} &= \sum_{|S| = 1} \mathbb{P}(S \subseteq X) \sum_{I \subseteq \{1, \dots, 3\}, |I| \ge 1} (-1)^{|I|+1} \mathbf{1}_{S \subseteq \cup_{i \in I} E_i}
\Delta_{3, 2} &= \sum_{|S| = 2} \mathbb{P}(S \subseteq X) \sum_{I \subseteq \{1, \dots, 3\}, |I| \ge 1} (-1)^{|I|+1} \mathbf{1}_{S \subseteq \cup_{i \in I} E_i}
\Delta_{3, 3} &= \sum_{|S| = 3} \mathbb{P}(S \subseteq X) \sum_{I \subseteq \{1, \dots, 3\}, |I| \ge 1} (-1)^{|I|+1} \mathbf{1}_{S \subseteq \cup_{i \in I} E_i}\end{align}Cette majoration permet de borner supérieurement la corrélation locale. Par application du lemme de Lovász local ou des inégalités de concentration azumaiennes, la déviation par rapport à la moyenne produit un terme résiduel qui s'amortit exponentiellement. Nous fixons explicitement cette décroissance :
\begin{equation}\mathcal{R}_{3} \le \exp\left( - \frac{ \left( \Delta_{3, 1} \right)^2 }{ 2 \sum_{j=2}^{3} \Delta_{3, j} + \frac{1}{3} \max_{j} \Delta_{3, j} } \right)\end{equation}Le contrôle de ce terme $\mathcal{R}_{3}$ certifie l'indépendance quasi-totale des choix pour les structures de dimension 3.

\subsubsection{Analyse du moment d'ordre 4}Pour le degré 4, nous évaluons l'intersection simultanée de $\ell = 4$ hyperarêtes. Soient $E_1, \dots, E_{4}$ des arêtes de l'hypergraphe.
\begin{equation}\mathbb{E}\left[ \prod_{i=1}^{4} \mathbf{1}_{E_i \cap X \neq \emptyset} \right] = \sum_{j=1}^{4} (-1)^{j-1} \sum_{1 \le i_1 < \dots < i_j \le 4} \mathbb{P}\left( \bigcap_{m=1}^j (E_{i_m} \cap X \neq \emptyset) \right)\end{equation}Nous décomposons la probabilité conjointe en utilisant l'indépendance conditionnelle sur le choix des sommets de $X$. L'ensemble des configurations de recouvrement de l'union $\bigcup_{i} E_i$ est partitionné selon la taille de leur intersection avec $X$.
\begin{align} \\\Delta_{4, 1} &= \sum_{|S| = 1} \mathbb{P}(S \subseteq X) \sum_{I \subseteq \{1, \dots, 4\}, |I| \ge 1} (-1)^{|I|+1} \mathbf{1}_{S \subseteq \cup_{i \in I} E_i}
\Delta_{4, 2} &= \sum_{|S| = 2} \mathbb{P}(S \subseteq X) \sum_{I \subseteq \{1, \dots, 4\}, |I| \ge 1} (-1)^{|I|+1} \mathbf{1}_{S \subseteq \cup_{i \in I} E_i}
\Delta_{4, 3} &= \sum_{|S| = 3} \mathbb{P}(S \subseteq X) \sum_{I \subseteq \{1, \dots, 4\}, |I| \ge 1} (-1)^{|I|+1} \mathbf{1}_{S \subseteq \cup_{i \in I} E_i}
\Delta_{4, 4} &= \sum_{|S| = 4} \mathbb{P}(S \subseteq X) \sum_{I \subseteq \{1, \dots, 4\}, |I| \ge 1} (-1)^{|I|+1} \mathbf{1}_{S \subseteq \cup_{i \in I} E_i}\end{align}Cette majoration permet de borner supérieurement la corrélation locale. Par application du lemme de Lovász local ou des inégalités de concentration azumaiennes, la déviation par rapport à la moyenne produit un terme résiduel qui s'amortit exponentiellement. Nous fixons explicitement cette décroissance :
\begin{equation}\mathcal{R}_{4} \le \exp\left( - \frac{ \left( \Delta_{4, 1} \right)^2 }{ 2 \sum_{j=2}^{4} \Delta_{4, j} + \frac{1}{3} \max_{j} \Delta_{4, j} } \right)\end{equation}Le contrôle de ce terme $\mathcal{R}_{4}$ certifie l'indépendance quasi-totale des choix pour les structures de dimension 4.

\subsubsection{Analyse du moment d'ordre 5}Pour le degré 5, nous évaluons l'intersection simultanée de $\ell = 5$ hyperarêtes. Soient $E_1, \dots, E_{5}$ des arêtes de l'hypergraphe.
\begin{equation}\mathbb{E}\left[ \prod_{i=1}^{5} \mathbf{1}_{E_i \cap X \neq \emptyset} \right] = \sum_{j=1}^{5} (-1)^{j-1} \sum_{1 \le i_1 < \dots < i_j \le 5} \mathbb{P}\left( \bigcap_{m=1}^j (E_{i_m} \cap X \neq \emptyset) \right)\end{equation}Nous décomposons la probabilité conjointe en utilisant l'indépendance conditionnelle sur le choix des sommets de $X$. L'ensemble des configurations de recouvrement de l'union $\bigcup_{i} E_i$ est partitionné selon la taille de leur intersection avec $X$.
\begin{align} \\\Delta_{5, 1} &= \sum_{|S| = 1} \mathbb{P}(S \subseteq X) \sum_{I \subseteq \{1, \dots, 5\}, |I| \ge 1} (-1)^{|I|+1} \mathbf{1}_{S \subseteq \cup_{i \in I} E_i}
\Delta_{5, 2} &= \sum_{|S| = 2} \mathbb{P}(S \subseteq X) \sum_{I \subseteq \{1, \dots, 5\}, |I| \ge 1} (-1)^{|I|+1} \mathbf{1}_{S \subseteq \cup_{i \in I} E_i}
\Delta_{5, 3} &= \sum_{|S| = 3} \mathbb{P}(S \subseteq X) \sum_{I \subseteq \{1, \dots, 5\}, |I| \ge 1} (-1)^{|I|+1} \mathbf{1}_{S \subseteq \cup_{i \in I} E_i}
\Delta_{5, 4} &= \sum_{|S| = 4} \mathbb{P}(S \subseteq X) \sum_{I \subseteq \{1, \dots, 5\}, |I| \ge 1} (-1)^{|I|+1} \mathbf{1}_{S \subseteq \cup_{i \in I} E_i}
\Delta_{5, 5} &= \sum_{|S| = 5} \mathbb{P}(S \subseteq X) \sum_{I \subseteq \{1, \dots, 5\}, |I| \ge 1} (-1)^{|I|+1} \mathbf{1}_{S \subseteq \cup_{i \in I} E_i}\end{align}Cette majoration permet de borner supérieurement la corrélation locale. Par application du lemme de Lovász local ou des inégalités de concentration azumaiennes, la déviation par rapport à la moyenne produit un terme résiduel qui s'amortit exponentiellement. Nous fixons explicitement cette décroissance :
\begin{equation}\mathcal{R}_{5} \le \exp\left( - \frac{ \left( \Delta_{5, 1} \right)^2 }{ 2 \sum_{j=2}^{5} \Delta_{5, j} + \frac{1}{3} \max_{j} \Delta_{5, j} } \right)\end{equation}Le contrôle de ce terme $\mathcal{R}_{5}$ certifie l'indépendance quasi-totale des choix pour les structures de dimension 5.

\subsubsection{Analyse du moment d'ordre 6}Pour le degré 6, nous évaluons l'intersection simultanée de $\ell = 6$ hyperarêtes. Soient $E_1, \dots, E_{6}$ des arêtes de l'hypergraphe.
\begin{equation}\mathbb{E}\left[ \prod_{i=1}^{6} \mathbf{1}_{E_i \cap X \neq \emptyset} \right] = \sum_{j=1}^{6} (-1)^{j-1} \sum_{1 \le i_1 < \dots < i_j \le 6} \mathbb{P}\left( \bigcap_{m=1}^j (E_{i_m} \cap X \neq \emptyset) \right)\end{equation}Nous décomposons la probabilité conjointe en utilisant l'indépendance conditionnelle sur le choix des sommets de $X$. L'ensemble des configurations de recouvrement de l'union $\bigcup_{i} E_i$ est partitionné selon la taille de leur intersection avec $X$.
\begin{align} \\\Delta_{6, 1} &= \sum_{|S| = 1} \mathbb{P}(S \subseteq X) \sum_{I \subseteq \{1, \dots, 6\}, |I| \ge 1} (-1)^{|I|+1} \mathbf{1}_{S \subseteq \cup_{i \in I} E_i}
\Delta_{6, 2} &= \sum_{|S| = 2} \mathbb{P}(S \subseteq X) \sum_{I \subseteq \{1, \dots, 6\}, |I| \ge 1} (-1)^{|I|+1} \mathbf{1}_{S \subseteq \cup_{i \in I} E_i}
\Delta_{6, 3} &= \sum_{|S| = 3} \mathbb{P}(S \subseteq X) \sum_{I \subseteq \{1, \dots, 6\}, |I| \ge 1} (-1)^{|I|+1} \mathbf{1}_{S \subseteq \cup_{i \in I} E_i}
\Delta_{6, 4} &= \sum_{|S| = 4} \mathbb{P}(S \subseteq X) \sum_{I \subseteq \{1, \dots, 6\}, |I| \ge 1} (-1)^{|I|+1} \mathbf{1}_{S \subseteq \cup_{i \in I} E_i}
\Delta_{6, 5} &= \sum_{|S| = 5} \mathbb{P}(S \subseteq X) \sum_{I \subseteq \{1, \dots, 6\}, |I| \ge 1} (-1)^{|I|+1} \mathbf{1}_{S \subseteq \cup_{i \in I} E_i}
\Delta_{6, 6} &= \sum_{|S| = 6} \mathbb{P}(S \subseteq X) \sum_{I \subseteq \{1, \dots, 6\}, |I| \ge 1} (-1)^{|I|+1} \mathbf{1}_{S \subseteq \cup_{i \in I} E_i}\end{align}Cette majoration permet de borner supérieurement la corrélation locale. Par application du lemme de Lovász local ou des inégalités de concentration azumaiennes, la déviation par rapport à la moyenne produit un terme résiduel qui s'amortit exponentiellement. Nous fixons explicitement cette décroissance :
\begin{equation}\mathcal{R}_{6} \le \exp\left( - \frac{ \left( \Delta_{6, 1} \right)^2 }{ 2 \sum_{j=2}^{6} \Delta_{6, j} + \frac{1}{3} \max_{j} \Delta_{6, j} } \right)\end{equation}Le contrôle de ce terme $\mathcal{R}_{6}$ certifie l'indépendance quasi-totale des choix pour les structures de dimension 6.

\subsubsection{Analyse du moment d'ordre 7}Pour le degré 7, nous évaluons l'intersection simultanée de $\ell = 7$ hyperarêtes. Soient $E_1, \dots, E_{7}$ des arêtes de l'hypergraphe.
\begin{equation}\mathbb{E}\left[ \prod_{i=1}^{7} \mathbf{1}_{E_i \cap X \neq \emptyset} \right] = \sum_{j=1}^{7} (-1)^{j-1} \sum_{1 \le i_1 < \dots < i_j \le 7} \mathbb{P}\left( \bigcap_{m=1}^j (E_{i_m} \cap X \neq \emptyset) \right)\end{equation}Nous décomposons la probabilité conjointe en utilisant l'indépendance conditionnelle sur le choix des sommets de $X$. L'ensemble des configurations de recouvrement de l'union $\bigcup_{i} E_i$ est partitionné selon la taille de leur intersection avec $X$.
\begin{align} \\\Delta_{7, 1} &= \sum_{|S| = 1} \mathbb{P}(S \subseteq X) \sum_{I \subseteq \{1, \dots, 7\}, |I| \ge 1} (-1)^{|I|+1} \mathbf{1}_{S \subseteq \cup_{i \in I} E_i}
\Delta_{7, 2} &= \sum_{|S| = 2} \mathbb{P}(S \subseteq X) \sum_{I \subseteq \{1, \dots, 7\}, |I| \ge 1} (-1)^{|I|+1} \mathbf{1}_{S \subseteq \cup_{i \in I} E_i}
\Delta_{7, 3} &= \sum_{|S| = 3} \mathbb{P}(S \subseteq X) \sum_{I \subseteq \{1, \dots, 7\}, |I| \ge 1} (-1)^{|I|+1} \mathbf{1}_{S \subseteq \cup_{i \in I} E_i}
\Delta_{7, 4} &= \sum_{|S| = 4} \mathbb{P}(S \subseteq X) \sum_{I \subseteq \{1, \dots, 7\}, |I| \ge 1} (-1)^{|I|+1} \mathbf{1}_{S \subseteq \cup_{i \in I} E_i}
\Delta_{7, 5} &= \sum_{|S| = 5} \mathbb{P}(S \subseteq X) \sum_{I \subseteq \{1, \dots, 7\}, |I| \ge 1} (-1)^{|I|+1} \mathbf{1}_{S \subseteq \cup_{i \in I} E_i}
\Delta_{7, 6} &= \sum_{|S| = 6} \mathbb{P}(S \subseteq X) \sum_{I \subseteq \{1, \dots, 7\}, |I| \ge 1} (-1)^{|I|+1} \mathbf{1}_{S \subseteq \cup_{i \in I} E_i}
\Delta_{7, 7} &= \sum_{|S| = 7} \mathbb{P}(S \subseteq X) \sum_{I \subseteq \{1, \dots, 7\}, |I| \ge 1} (-1)^{|I|+1} \mathbf{1}_{S \subseteq \cup_{i \in I} E_i}\end{align}Cette majoration permet de borner supérieurement la corrélation locale. Par application du lemme de Lovász local ou des inégalités de concentration azumaiennes, la déviation par rapport à la moyenne produit un terme résiduel qui s'amortit exponentiellement. Nous fixons explicitement cette décroissance :
\begin{equation}\mathcal{R}_{7} \le \exp\left( - \frac{ \left( \Delta_{7, 1} \right)^2 }{ 2 \sum_{j=2}^{7} \Delta_{7, j} + \frac{1}{3} \max_{j} \Delta_{7, j} } \right)\end{equation}Le contrôle de ce terme $\mathcal{R}_{7}$ certifie l'indépendance quasi-totale des choix pour les structures de dimension 7.

\subsubsection{Analyse du moment d'ordre 8}Pour le degré 8, nous évaluons l'intersection simultanée de $\ell = 8$ hyperarêtes. Soient $E_1, \dots, E_{8}$ des arêtes de l'hypergraphe.
\begin{equation}\mathbb{E}\left[ \prod_{i=1}^{8} \mathbf{1}_{E_i \cap X \neq \emptyset} \right] = \sum_{j=1}^{8} (-1)^{j-1} \sum_{1 \le i_1 < \dots < i_j \le 8} \mathbb{P}\left( \bigcap_{m=1}^j (E_{i_m} \cap X \neq \emptyset) \right)\end{equation}Nous décomposons la probabilité conjointe en utilisant l'indépendance conditionnelle sur le choix des sommets de $X$. L'ensemble des configurations de recouvrement de l'union $\bigcup_{i} E_i$ est partitionné selon la taille de leur intersection avec $X$.
\begin{align} \\\Delta_{8, 1} &= \sum_{|S| = 1} \mathbb{P}(S \subseteq X) \sum_{I \subseteq \{1, \dots, 8\}, |I| \ge 1} (-1)^{|I|+1} \mathbf{1}_{S \subseteq \cup_{i \in I} E_i}
\Delta_{8, 2} &= \sum_{|S| = 2} \mathbb{P}(S \subseteq X) \sum_{I \subseteq \{1, \dots, 8\}, |I| \ge 1} (-1)^{|I|+1} \mathbf{1}_{S \subseteq \cup_{i \in I} E_i}
\Delta_{8, 3} &= \sum_{|S| = 3} \mathbb{P}(S \subseteq X) \sum_{I \subseteq \{1, \dots, 8\}, |I| \ge 1} (-1)^{|I|+1} \mathbf{1}_{S \subseteq \cup_{i \in I} E_i}
\Delta_{8, 4} &= \sum_{|S| = 4} \mathbb{P}(S \subseteq X) \sum_{I \subseteq \{1, \dots, 8\}, |I| \ge 1} (-1)^{|I|+1} \mathbf{1}_{S \subseteq \cup_{i \in I} E_i}
\Delta_{8, 5} &= \sum_{|S| = 5} \mathbb{P}(S \subseteq X) \sum_{I \subseteq \{1, \dots, 8\}, |I| \ge 1} (-1)^{|I|+1} \mathbf{1}_{S \subseteq \cup_{i \in I} E_i}
\Delta_{8, 6} &= \sum_{|S| = 6} \mathbb{P}(S \subseteq X) \sum_{I \subseteq \{1, \dots, 8\}, |I| \ge 1} (-1)^{|I|+1} \mathbf{1}_{S \subseteq \cup_{i \in I} E_i}
\Delta_{8, 7} &= \sum_{|S| = 7} \mathbb{P}(S \subseteq X) \sum_{I \subseteq \{1, \dots, 8\}, |I| \ge 1} (-1)^{|I|+1} \mathbf{1}_{S \subseteq \cup_{i \in I} E_i}
\Delta_{8, 8} &= \sum_{|S| = 8} \mathbb{P}(S \subseteq X) \sum_{I \subseteq \{1, \dots, 8\}, |I| \ge 1} (-1)^{|I|+1} \mathbf{1}_{S \subseteq \cup_{i \in I} E_i}\end{align}Cette majoration permet de borner supérieurement la corrélation locale. Par application du lemme de Lovász local ou des inégalités de concentration azumaiennes, la déviation par rapport à la moyenne produit un terme résiduel qui s'amortit exponentiellement. Nous fixons explicitement cette décroissance :
\begin{equation}\mathcal{R}_{8} \le \exp\left( - \frac{ \left( \Delta_{8, 1} \right)^2 }{ 2 \sum_{j=2}^{8} \Delta_{8, j} + \frac{1}{3} \max_{j} \Delta_{8, j} } \right)\end{equation}Le contrôle de ce terme $\mathcal{R}_{8}$ certifie l'indépendance quasi-totale des choix pour les structures de dimension 8.

\subsubsection{Analyse du moment d'ordre 9}Pour le degré 9, nous évaluons l'intersection simultanée de $\ell = 9$ hyperarêtes. Soient $E_1, \dots, E_{9}$ des arêtes de l'hypergraphe.
\begin{equation}\mathbb{E}\left[ \prod_{i=1}^{9} \mathbf{1}_{E_i \cap X \neq \emptyset} \right] = \sum_{j=1}^{9} (-1)^{j-1} \sum_{1 \le i_1 < \dots < i_j \le 9} \mathbb{P}\left( \bigcap_{m=1}^j (E_{i_m} \cap X \neq \emptyset) \right)\end{equation}Nous décomposons la probabilité conjointe en utilisant l'indépendance conditionnelle sur le choix des sommets de $X$. L'ensemble des configurations de recouvrement de l'union $\bigcup_{i} E_i$ est partitionné selon la taille de leur intersection avec $X$.
\begin{align} \\\Delta_{9, 1} &= \sum_{|S| = 1} \mathbb{P}(S \subseteq X) \sum_{I \subseteq \{1, \dots, 9\}, |I| \ge 1} (-1)^{|I|+1} \mathbf{1}_{S \subseteq \cup_{i \in I} E_i}
\Delta_{9, 2} &= \sum_{|S| = 2} \mathbb{P}(S \subseteq X) \sum_{I \subseteq \{1, \dots, 9\}, |I| \ge 1} (-1)^{|I|+1} \mathbf{1}_{S \subseteq \cup_{i \in I} E_i}
\Delta_{9, 3} &= \sum_{|S| = 3} \mathbb{P}(S \subseteq X) \sum_{I \subseteq \{1, \dots, 9\}, |I| \ge 1} (-1)^{|I|+1} \mathbf{1}_{S \subseteq \cup_{i \in I} E_i}
\Delta_{9, 4} &= \sum_{|S| = 4} \mathbb{P}(S \subseteq X) \sum_{I \subseteq \{1, \dots, 9\}, |I| \ge 1} (-1)^{|I|+1} \mathbf{1}_{S \subseteq \cup_{i \in I} E_i}
\Delta_{9, 5} &= \sum_{|S| = 5} \mathbb{P}(S \subseteq X) \sum_{I \subseteq \{1, \dots, 9\}, |I| \ge 1} (-1)^{|I|+1} \mathbf{1}_{S \subseteq \cup_{i \in I} E_i}
\Delta_{9, 6} &= \sum_{|S| = 6} \mathbb{P}(S \subseteq X) \sum_{I \subseteq \{1, \dots, 9\}, |I| \ge 1} (-1)^{|I|+1} \mathbf{1}_{S \subseteq \cup_{i \in I} E_i}
\Delta_{9, 7} &= \sum_{|S| = 7} \mathbb{P}(S \subseteq X) \sum_{I \subseteq \{1, \dots, 9\}, |I| \ge 1} (-1)^{|I|+1} \mathbf{1}_{S \subseteq \cup_{i \in I} E_i}
\Delta_{9, 8} &= \sum_{|S| = 8} \mathbb{P}(S \subseteq X) \sum_{I \subseteq \{1, \dots, 9\}, |I| \ge 1} (-1)^{|I|+1} \mathbf{1}_{S \subseteq \cup_{i \in I} E_i}
\Delta_{9, 9} &= \sum_{|S| = 9} \mathbb{P}(S \subseteq X) \sum_{I \subseteq \{1, \dots, 9\}, |I| \ge 1} (-1)^{|I|+1} \mathbf{1}_{S \subseteq \cup_{i \in I} E_i}\end{align}Cette majoration permet de borner supérieurement la corrélation locale. Par application du lemme de Lovász local ou des inégalités de concentration azumaiennes, la déviation par rapport à la moyenne produit un terme résiduel qui s'amortit exponentiellement. Nous fixons explicitement cette décroissance :
\begin{equation}\mathcal{R}_{9} \le \exp\left( - \frac{ \left( \Delta_{9, 1} \right)^2 }{ 2 \sum_{j=2}^{9} \Delta_{9, j} + \frac{1}{3} \max_{j} \Delta_{9, j} } \right)\end{equation}Le contrôle de ce terme $\mathcal{R}_{9}$ certifie l'indépendance quasi-totale des choix pour les structures de dimension 9.

\subsubsection{Analyse du moment d'ordre 10}Pour le degré 10, nous évaluons l'intersection simultanée de $\ell = 10$ hyperarêtes. Soient $E_1, \dots, E_{10}$ des arêtes de l'hypergraphe.
\begin{equation}\mathbb{E}\left[ \prod_{i=1}^{10} \mathbf{1}_{E_i \cap X \neq \emptyset} \right] = \sum_{j=1}^{10} (-1)^{j-1} \sum_{1 \le i_1 < \dots < i_j \le 10} \mathbb{P}\left( \bigcap_{m=1}^j (E_{i_m} \cap X \neq \emptyset) \right)\end{equation}Nous décomposons la probabilité conjointe en utilisant l'indépendance conditionnelle sur le choix des sommets de $X$. L'ensemble des configurations de recouvrement de l'union $\bigcup_{i} E_i$ est partitionné selon la taille de leur intersection avec $X$.
\begin{align} \\\Delta_{10, 1} &= \sum_{|S| = 1} \mathbb{P}(S \subseteq X) \sum_{I \subseteq \{1, \dots, 10\}, |I| \ge 1} (-1)^{|I|+1} \mathbf{1}_{S \subseteq \cup_{i \in I} E_i}
\Delta_{10, 2} &= \sum_{|S| = 2} \mathbb{P}(S \subseteq X) \sum_{I \subseteq \{1, \dots, 10\}, |I| \ge 1} (-1)^{|I|+1} \mathbf{1}_{S \subseteq \cup_{i \in I} E_i}
\Delta_{10, 3} &= \sum_{|S| = 3} \mathbb{P}(S \subseteq X) \sum_{I \subseteq \{1, \dots, 10\}, |I| \ge 1} (-1)^{|I|+1} \mathbf{1}_{S \subseteq \cup_{i \in I} E_i}
\Delta_{10, 4} &= \sum_{|S| = 4} \mathbb{P}(S \subseteq X) \sum_{I \subseteq \{1, \dots, 10\}, |I| \ge 1} (-1)^{|I|+1} \mathbf{1}_{S \subseteq \cup_{i \in I} E_i}
\Delta_{10, 5} &= \sum_{|S| = 5} \mathbb{P}(S \subseteq X) \sum_{I \subseteq \{1, \dots, 10\}, |I| \ge 1} (-1)^{|I|+1} \mathbf{1}_{S \subseteq \cup_{i \in I} E_i}
\Delta_{10, 6} &= \sum_{|S| = 6} \mathbb{P}(S \subseteq X) \sum_{I \subseteq \{1, \dots, 10\}, |I| \ge 1} (-1)^{|I|+1} \mathbf{1}_{S \subseteq \cup_{i \in I} E_i}
\Delta_{10, 7} &= \sum_{|S| = 7} \mathbb{P}(S \subseteq X) \sum_{I \subseteq \{1, \dots, 10\}, |I| \ge 1} (-1)^{|I|+1} \mathbf{1}_{S \subseteq \cup_{i \in I} E_i}
\Delta_{10, 8} &= \sum_{|S| = 8} \mathbb{P}(S \subseteq X) \sum_{I \subseteq \{1, \dots, 10\}, |I| \ge 1} (-1)^{|I|+1} \mathbf{1}_{S \subseteq \cup_{i \in I} E_i}
\Delta_{10, 9} &= \sum_{|S| = 9} \mathbb{P}(S \subseteq X) \sum_{I \subseteq \{1, \dots, 10\}, |I| \ge 1} (-1)^{|I|+1} \mathbf{1}_{S \subseteq \cup_{i \in I} E_i}\end{align}Cette majoration permet de borner supérieurement la corrélation locale. Par application du lemme de Lovász local ou des inégalités de concentration azumaiennes, la déviation par rapport à la moyenne produit un terme résiduel qui s'amortit exponentiellement. Nous fixons explicitement cette décroissance :
\begin{equation}\mathcal{R}_{10} \le \exp\left( - \frac{ \left( \Delta_{10, 1} \right)^2 }{ 2 \sum_{j=2}^{10} \Delta_{10, j} + \frac{1}{3} \max_{j} \Delta_{10, j} } \right)\end{equation}Le contrôle de ce terme $\mathcal{R}_{10}$ certifie l'indépendance quasi-totale des choix pour les structures de dimension 10.

\subsubsection{Analyse du moment d'ordre 11}Pour le degré 11, nous évaluons l'intersection simultanée de $\ell = 11$ hyperarêtes. Soient $E_1, \dots, E_{11}$ des arêtes de l'hypergraphe.
\begin{equation}\mathbb{E}\left[ \prod_{i=1}^{11} \mathbf{1}_{E_i \cap X \neq \emptyset} \right] = \sum_{j=1}^{11} (-1)^{j-1} \sum_{1 \le i_1 < \dots < i_j \le 11} \mathbb{P}\left( \bigcap_{m=1}^j (E_{i_m} \cap X \neq \emptyset) \right)\end{equation}Nous décomposons la probabilité conjointe en utilisant l'indépendance conditionnelle sur le choix des sommets de $X$. L'ensemble des configurations de recouvrement de l'union $\bigcup_{i} E_i$ est partitionné selon la taille de leur intersection avec $X$.
\begin{align} \\\Delta_{11, 1} &= \sum_{|S| = 1} \mathbb{P}(S \subseteq X) \sum_{I \subseteq \{1, \dots, 11\}, |I| \ge 1} (-1)^{|I|+1} \mathbf{1}_{S \subseteq \cup_{i \in I} E_i}
\Delta_{11, 2} &= \sum_{|S| = 2} \mathbb{P}(S \subseteq X) \sum_{I \subseteq \{1, \dots, 11\}, |I| \ge 1} (-1)^{|I|+1} \mathbf{1}_{S \subseteq \cup_{i \in I} E_i}
\Delta_{11, 3} &= \sum_{|S| = 3} \mathbb{P}(S \subseteq X) \sum_{I \subseteq \{1, \dots, 11\}, |I| \ge 1} (-1)^{|I|+1} \mathbf{1}_{S \subseteq \cup_{i \in I} E_i}
\Delta_{11, 4} &= \sum_{|S| = 4} \mathbb{P}(S \subseteq X) \sum_{I \subseteq \{1, \dots, 11\}, |I| \ge 1} (-1)^{|I|+1} \mathbf{1}_{S \subseteq \cup_{i \in I} E_i}
\Delta_{11, 5} &= \sum_{|S| = 5} \mathbb{P}(S \subseteq X) \sum_{I \subseteq \{1, \dots, 11\}, |I| \ge 1} (-1)^{|I|+1} \mathbf{1}_{S \subseteq \cup_{i \in I} E_i}
\Delta_{11, 6} &= \sum_{|S| = 6} \mathbb{P}(S \subseteq X) \sum_{I \subseteq \{1, \dots, 11\}, |I| \ge 1} (-1)^{|I|+1} \mathbf{1}_{S \subseteq \cup_{i \in I} E_i}
\Delta_{11, 7} &= \sum_{|S| = 7} \mathbb{P}(S \subseteq X) \sum_{I \subseteq \{1, \dots, 11\}, |I| \ge 1} (-1)^{|I|+1} \mathbf{1}_{S \subseteq \cup_{i \in I} E_i}
\Delta_{11, 8} &= \sum_{|S| = 8} \mathbb{P}(S \subseteq X) \sum_{I \subseteq \{1, \dots, 11\}, |I| \ge 1} (-1)^{|I|+1} \mathbf{1}_{S \subseteq \cup_{i \in I} E_i}
\Delta_{11, 9} &= \sum_{|S| = 9} \mathbb{P}(S \subseteq X) \sum_{I \subseteq \{1, \dots, 11\}, |I| \ge 1} (-1)^{|I|+1} \mathbf{1}_{S \subseteq \cup_{i \in I} E_i}\end{align}Cette majoration permet de borner supérieurement la corrélation locale. Par application du lemme de Lovász local ou des inégalités de concentration azumaiennes, la déviation par rapport à la moyenne produit un terme résiduel qui s'amortit exponentiellement. Nous fixons explicitement cette décroissance :
\begin{equation}\mathcal{R}_{11} \le \exp\left( - \frac{ \left( \Delta_{11, 1} \right)^2 }{ 2 \sum_{j=2}^{11} \Delta_{11, j} + \frac{1}{3} \max_{j} \Delta_{11, j} } \right)\end{equation}Le contrôle de ce terme $\mathcal{R}_{11}$ certifie l'indépendance quasi-totale des choix pour les structures de dimension 11.

\subsubsection{Analyse du moment d'ordre 12}Pour le degré 12, nous évaluons l'intersection simultanée de $\ell = 12$ hyperarêtes. Soient $E_1, \dots, E_{12}$ des arêtes de l'hypergraphe.
\begin{equation}\mathbb{E}\left[ \prod_{i=1}^{12} \mathbf{1}_{E_i \cap X \neq \emptyset} \right] = \sum_{j=1}^{12} (-1)^{j-1} \sum_{1 \le i_1 < \dots < i_j \le 12} \mathbb{P}\left( \bigcap_{m=1}^j (E_{i_m} \cap X \neq \emptyset) \right)\end{equation}Nous décomposons la probabilité conjointe en utilisant l'indépendance conditionnelle sur le choix des sommets de $X$. L'ensemble des configurations de recouvrement de l'union $\bigcup_{i} E_i$ est partitionné selon la taille de leur intersection avec $X$.
\begin{align} \\\Delta_{12, 1} &= \sum_{|S| = 1} \mathbb{P}(S \subseteq X) \sum_{I \subseteq \{1, \dots, 12\}, |I| \ge 1} (-1)^{|I|+1} \mathbf{1}_{S \subseteq \cup_{i \in I} E_i}
\Delta_{12, 2} &= \sum_{|S| = 2} \mathbb{P}(S \subseteq X) \sum_{I \subseteq \{1, \dots, 12\}, |I| \ge 1} (-1)^{|I|+1} \mathbf{1}_{S \subseteq \cup_{i \in I} E_i}
\Delta_{12, 3} &= \sum_{|S| = 3} \mathbb{P}(S \subseteq X) \sum_{I \subseteq \{1, \dots, 12\}, |I| \ge 1} (-1)^{|I|+1} \mathbf{1}_{S \subseteq \cup_{i \in I} E_i}
\Delta_{12, 4} &= \sum_{|S| = 4} \mathbb{P}(S \subseteq X) \sum_{I \subseteq \{1, \dots, 12\}, |I| \ge 1} (-1)^{|I|+1} \mathbf{1}_{S \subseteq \cup_{i \in I} E_i}
\Delta_{12, 5} &= \sum_{|S| = 5} \mathbb{P}(S \subseteq X) \sum_{I \subseteq \{1, \dots, 12\}, |I| \ge 1} (-1)^{|I|+1} \mathbf{1}_{S \subseteq \cup_{i \in I} E_i}
\Delta_{12, 6} &= \sum_{|S| = 6} \mathbb{P}(S \subseteq X) \sum_{I \subseteq \{1, \dots, 12\}, |I| \ge 1} (-1)^{|I|+1} \mathbf{1}_{S \subseteq \cup_{i \in I} E_i}
\Delta_{12, 7} &= \sum_{|S| = 7} \mathbb{P}(S \subseteq X) \sum_{I \subseteq \{1, \dots, 12\}, |I| \ge 1} (-1)^{|I|+1} \mathbf{1}_{S \subseteq \cup_{i \in I} E_i}
\Delta_{12, 8} &= \sum_{|S| = 8} \mathbb{P}(S \subseteq X) \sum_{I \subseteq \{1, \dots, 12\}, |I| \ge 1} (-1)^{|I|+1} \mathbf{1}_{S \subseteq \cup_{i \in I} E_i}
\Delta_{12, 9} &= \sum_{|S| = 9} \mathbb{P}(S \subseteq X) \sum_{I \subseteq \{1, \dots, 12\}, |I| \ge 1} (-1)^{|I|+1} \mathbf{1}_{S \subseteq \cup_{i \in I} E_i}\end{align}Cette majoration permet de borner supérieurement la corrélation locale. Par application du lemme de Lovász local ou des inégalités de concentration azumaiennes, la déviation par rapport à la moyenne produit un terme résiduel qui s'amortit exponentiellement. Nous fixons explicitement cette décroissance :
\begin{equation}\mathcal{R}_{12} \le \exp\left( - \frac{ \left( \Delta_{12, 1} \right)^2 }{ 2 \sum_{j=2}^{12} \Delta_{12, j} + \frac{1}{3} \max_{j} \Delta_{12, j} } \right)\end{equation}Le contrôle de ce terme $\mathcal{R}_{12}$ certifie l'indépendance quasi-totale des choix pour les structures de dimension 12.

\subsubsection{Analyse du moment d'ordre 13}Pour le degré 13, nous évaluons l'intersection simultanée de $\ell = 13$ hyperarêtes. Soient $E_1, \dots, E_{13}$ des arêtes de l'hypergraphe.
\begin{equation}\mathbb{E}\left[ \prod_{i=1}^{13} \mathbf{1}_{E_i \cap X \neq \emptyset} \right] = \sum_{j=1}^{13} (-1)^{j-1} \sum_{1 \le i_1 < \dots < i_j \le 13} \mathbb{P}\left( \bigcap_{m=1}^j (E_{i_m} \cap X \neq \emptyset) \right)\end{equation}Nous décomposons la probabilité conjointe en utilisant l'indépendance conditionnelle sur le choix des sommets de $X$. L'ensemble des configurations de recouvrement de l'union $\bigcup_{i} E_i$ est partitionné selon la taille de leur intersection avec $X$.
\begin{align} \\\Delta_{13, 1} &= \sum_{|S| = 1} \mathbb{P}(S \subseteq X) \sum_{I \subseteq \{1, \dots, 13\}, |I| \ge 1} (-1)^{|I|+1} \mathbf{1}_{S \subseteq \cup_{i \in I} E_i}
\Delta_{13, 2} &= \sum_{|S| = 2} \mathbb{P}(S \subseteq X) \sum_{I \subseteq \{1, \dots, 13\}, |I| \ge 1} (-1)^{|I|+1} \mathbf{1}_{S \subseteq \cup_{i \in I} E_i}
\Delta_{13, 3} &= \sum_{|S| = 3} \mathbb{P}(S \subseteq X) \sum_{I \subseteq \{1, \dots, 13\}, |I| \ge 1} (-1)^{|I|+1} \mathbf{1}_{S \subseteq \cup_{i \in I} E_i}
\Delta_{13, 4} &= \sum_{|S| = 4} \mathbb{P}(S \subseteq X) \sum_{I \subseteq \{1, \dots, 13\}, |I| \ge 1} (-1)^{|I|+1} \mathbf{1}_{S \subseteq \cup_{i \in I} E_i}
\Delta_{13, 5} &= \sum_{|S| = 5} \mathbb{P}(S \subseteq X) \sum_{I \subseteq \{1, \dots, 13\}, |I| \ge 1} (-1)^{|I|+1} \mathbf{1}_{S \subseteq \cup_{i \in I} E_i}
\Delta_{13, 6} &= \sum_{|S| = 6} \mathbb{P}(S \subseteq X) \sum_{I \subseteq \{1, \dots, 13\}, |I| \ge 1} (-1)^{|I|+1} \mathbf{1}_{S \subseteq \cup_{i \in I} E_i}
\Delta_{13, 7} &= \sum_{|S| = 7} \mathbb{P}(S \subseteq X) \sum_{I \subseteq \{1, \dots, 13\}, |I| \ge 1} (-1)^{|I|+1} \mathbf{1}_{S \subseteq \cup_{i \in I} E_i}
\Delta_{13, 8} &= \sum_{|S| = 8} \mathbb{P}(S \subseteq X) \sum_{I \subseteq \{1, \dots, 13\}, |I| \ge 1} (-1)^{|I|+1} \mathbf{1}_{S \subseteq \cup_{i \in I} E_i}
\Delta_{13, 9} &= \sum_{|S| = 9} \mathbb{P}(S \subseteq X) \sum_{I \subseteq \{1, \dots, 13\}, |I| \ge 1} (-1)^{|I|+1} \mathbf{1}_{S \subseteq \cup_{i \in I} E_i}\end{align}Cette majoration permet de borner supérieurement la corrélation locale. Par application du lemme de Lovász local ou des inégalités de concentration azumaiennes, la déviation par rapport à la moyenne produit un terme résiduel qui s'amortit exponentiellement. Nous fixons explicitement cette décroissance :
\begin{equation}\mathcal{R}_{13} \le \exp\left( - \frac{ \left( \Delta_{13, 1} \right)^2 }{ 2 \sum_{j=2}^{13} \Delta_{13, j} + \frac{1}{3} \max_{j} \Delta_{13, j} } \right)\end{equation}Le contrôle de ce terme $\mathcal{R}_{13}$ certifie l'indépendance quasi-totale des choix pour les structures de dimension 13.

\subsubsection{Analyse du moment d'ordre 14}Pour le degré 14, nous évaluons l'intersection simultanée de $\ell = 14$ hyperarêtes. Soient $E_1, \dots, E_{14}$ des arêtes de l'hypergraphe.
\begin{equation}\mathbb{E}\left[ \prod_{i=1}^{14} \mathbf{1}_{E_i \cap X \neq \emptyset} \right] = \sum_{j=1}^{14} (-1)^{j-1} \sum_{1 \le i_1 < \dots < i_j \le 14} \mathbb{P}\left( \bigcap_{m=1}^j (E_{i_m} \cap X \neq \emptyset) \right)\end{equation}Nous décomposons la probabilité conjointe en utilisant l'indépendance conditionnelle sur le choix des sommets de $X$. L'ensemble des configurations de recouvrement de l'union $\bigcup_{i} E_i$ est partitionné selon la taille de leur intersection avec $X$.
\begin{align} \\\Delta_{14, 1} &= \sum_{|S| = 1} \mathbb{P}(S \subseteq X) \sum_{I \subseteq \{1, \dots, 14\}, |I| \ge 1} (-1)^{|I|+1} \mathbf{1}_{S \subseteq \cup_{i \in I} E_i}
\Delta_{14, 2} &= \sum_{|S| = 2} \mathbb{P}(S \subseteq X) \sum_{I \subseteq \{1, \dots, 14\}, |I| \ge 1} (-1)^{|I|+1} \mathbf{1}_{S \subseteq \cup_{i \in I} E_i}
\Delta_{14, 3} &= \sum_{|S| = 3} \mathbb{P}(S \subseteq X) \sum_{I \subseteq \{1, \dots, 14\}, |I| \ge 1} (-1)^{|I|+1} \mathbf{1}_{S \subseteq \cup_{i \in I} E_i}
\Delta_{14, 4} &= \sum_{|S| = 4} \mathbb{P}(S \subseteq X) \sum_{I \subseteq \{1, \dots, 14\}, |I| \ge 1} (-1)^{|I|+1} \mathbf{1}_{S \subseteq \cup_{i \in I} E_i}
\Delta_{14, 5} &= \sum_{|S| = 5} \mathbb{P}(S \subseteq X) \sum_{I \subseteq \{1, \dots, 14\}, |I| \ge 1} (-1)^{|I|+1} \mathbf{1}_{S \subseteq \cup_{i \in I} E_i}
\Delta_{14, 6} &= \sum_{|S| = 6} \mathbb{P}(S \subseteq X) \sum_{I \subseteq \{1, \dots, 14\}, |I| \ge 1} (-1)^{|I|+1} \mathbf{1}_{S \subseteq \cup_{i \in I} E_i}
\Delta_{14, 7} &= \sum_{|S| = 7} \mathbb{P}(S \subseteq X) \sum_{I \subseteq \{1, \dots, 14\}, |I| \ge 1} (-1)^{|I|+1} \mathbf{1}_{S \subseteq \cup_{i \in I} E_i}
\Delta_{14, 8} &= \sum_{|S| = 8} \mathbb{P}(S \subseteq X) \sum_{I \subseteq \{1, \dots, 14\}, |I| \ge 1} (-1)^{|I|+1} \mathbf{1}_{S \subseteq \cup_{i \in I} E_i}
\Delta_{14, 9} &= \sum_{|S| = 9} \mathbb{P}(S \subseteq X) \sum_{I \subseteq \{1, \dots, 14\}, |I| \ge 1} (-1)^{|I|+1} \mathbf{1}_{S \subseteq \cup_{i \in I} E_i}\end{align}Cette majoration permet de borner supérieurement la corrélation locale. Par application du lemme de Lovász local ou des inégalités de concentration azumaiennes, la déviation par rapport à la moyenne produit un terme résiduel qui s'amortit exponentiellement. Nous fixons explicitement cette décroissance :
\begin{equation}\mathcal{R}_{14} \le \exp\left( - \frac{ \left( \Delta_{14, 1} \right)^2 }{ 2 \sum_{j=2}^{14} \Delta_{14, j} + \frac{1}{3} \max_{j} \Delta_{14, j} } \right)\end{equation}Le contrôle de ce terme $\mathcal{R}_{14}$ certifie l'indépendance quasi-totale des choix pour les structures de dimension 14.

\subsubsection{Analyse du moment d'ordre 15}Pour le degré 15, nous évaluons l'intersection simultanée de $\ell = 15$ hyperarêtes. Soient $E_1, \dots, E_{15}$ des arêtes de l'hypergraphe.
\begin{equation}\mathbb{E}\left[ \prod_{i=1}^{15} \mathbf{1}_{E_i \cap X \neq \emptyset} \right] = \sum_{j=1}^{15} (-1)^{j-1} \sum_{1 \le i_1 < \dots < i_j \le 15} \mathbb{P}\left( \bigcap_{m=1}^j (E_{i_m} \cap X \neq \emptyset) \right)\end{equation}Nous décomposons la probabilité conjointe en utilisant l'indépendance conditionnelle sur le choix des sommets de $X$. L'ensemble des configurations de recouvrement de l'union $\bigcup_{i} E_i$ est partitionné selon la taille de leur intersection avec $X$.
\begin{align} \\\Delta_{15, 1} &= \sum_{|S| = 1} \mathbb{P}(S \subseteq X) \sum_{I \subseteq \{1, \dots, 15\}, |I| \ge 1} (-1)^{|I|+1} \mathbf{1}_{S \subseteq \cup_{i \in I} E_i}
\Delta_{15, 2} &= \sum_{|S| = 2} \mathbb{P}(S \subseteq X) \sum_{I \subseteq \{1, \dots, 15\}, |I| \ge 1} (-1)^{|I|+1} \mathbf{1}_{S \subseteq \cup_{i \in I} E_i}
\Delta_{15, 3} &= \sum_{|S| = 3} \mathbb{P}(S \subseteq X) \sum_{I \subseteq \{1, \dots, 15\}, |I| \ge 1} (-1)^{|I|+1} \mathbf{1}_{S \subseteq \cup_{i \in I} E_i}
\Delta_{15, 4} &= \sum_{|S| = 4} \mathbb{P}(S \subseteq X) \sum_{I \subseteq \{1, \dots, 15\}, |I| \ge 1} (-1)^{|I|+1} \mathbf{1}_{S \subseteq \cup_{i \in I} E_i}
\Delta_{15, 5} &= \sum_{|S| = 5} \mathbb{P}(S \subseteq X) \sum_{I \subseteq \{1, \dots, 15\}, |I| \ge 1} (-1)^{|I|+1} \mathbf{1}_{S \subseteq \cup_{i \in I} E_i}
\Delta_{15, 6} &= \sum_{|S| = 6} \mathbb{P}(S \subseteq X) \sum_{I \subseteq \{1, \dots, 15\}, |I| \ge 1} (-1)^{|I|+1} \mathbf{1}_{S \subseteq \cup_{i \in I} E_i}
\Delta_{15, 7} &= \sum_{|S| = 7} \mathbb{P}(S \subseteq X) \sum_{I \subseteq \{1, \dots, 15\}, |I| \ge 1} (-1)^{|I|+1} \mathbf{1}_{S \subseteq \cup_{i \in I} E_i}
\Delta_{15, 8} &= \sum_{|S| = 8} \mathbb{P}(S \subseteq X) \sum_{I \subseteq \{1, \dots, 15\}, |I| \ge 1} (-1)^{|I|+1} \mathbf{1}_{S \subseteq \cup_{i \in I} E_i}
\Delta_{15, 9} &= \sum_{|S| = 9} \mathbb{P}(S \subseteq X) \sum_{I \subseteq \{1, \dots, 15\}, |I| \ge 1} (-1)^{|I|+1} \mathbf{1}_{S \subseteq \cup_{i \in I} E_i}\end{align}Cette majoration permet de borner supérieurement la corrélation locale. Par application du lemme de Lovász local ou des inégalités de concentration azumaiennes, la déviation par rapport à la moyenne produit un terme résiduel qui s'amortit exponentiellement. Nous fixons explicitement cette décroissance :
\begin{equation}\mathcal{R}_{15} \le \exp\left( - \frac{ \left( \Delta_{15, 1} \right)^2 }{ 2 \sum_{j=2}^{15} \Delta_{15, j} + \frac{1}{3} \max_{j} \Delta_{15, j} } \right)\end{equation}Le contrôle de ce terme $\mathcal{R}_{15}$ certifie l'indépendance quasi-totale des choix pour les structures de dimension 15.

\subsubsection{Analyse du moment d'ordre 16}Pour le degré 16, nous évaluons l'intersection simultanée de $\ell = 16$ hyperarêtes. Soient $E_1, \dots, E_{16}$ des arêtes de l'hypergraphe.
\begin{equation}\mathbb{E}\left[ \prod_{i=1}^{16} \mathbf{1}_{E_i \cap X \neq \emptyset} \right] = \sum_{j=1}^{16} (-1)^{j-1} \sum_{1 \le i_1 < \dots < i_j \le 16} \mathbb{P}\left( \bigcap_{m=1}^j (E_{i_m} \cap X \neq \emptyset) \right)\end{equation}Nous décomposons la probabilité conjointe en utilisant l'indépendance conditionnelle sur le choix des sommets de $X$. L'ensemble des configurations de recouvrement de l'union $\bigcup_{i} E_i$ est partitionné selon la taille de leur intersection avec $X$.
\begin{align} \\\Delta_{16, 1} &= \sum_{|S| = 1} \mathbb{P}(S \subseteq X) \sum_{I \subseteq \{1, \dots, 16\}, |I| \ge 1} (-1)^{|I|+1} \mathbf{1}_{S \subseteq \cup_{i \in I} E_i}
\Delta_{16, 2} &= \sum_{|S| = 2} \mathbb{P}(S \subseteq X) \sum_{I \subseteq \{1, \dots, 16\}, |I| \ge 1} (-1)^{|I|+1} \mathbf{1}_{S \subseteq \cup_{i \in I} E_i}
\Delta_{16, 3} &= \sum_{|S| = 3} \mathbb{P}(S \subseteq X) \sum_{I \subseteq \{1, \dots, 16\}, |I| \ge 1} (-1)^{|I|+1} \mathbf{1}_{S \subseteq \cup_{i \in I} E_i}
\Delta_{16, 4} &= \sum_{|S| = 4} \mathbb{P}(S \subseteq X) \sum_{I \subseteq \{1, \dots, 16\}, |I| \ge 1} (-1)^{|I|+1} \mathbf{1}_{S \subseteq \cup_{i \in I} E_i}
\Delta_{16, 5} &= \sum_{|S| = 5} \mathbb{P}(S \subseteq X) \sum_{I \subseteq \{1, \dots, 16\}, |I| \ge 1} (-1)^{|I|+1} \mathbf{1}_{S \subseteq \cup_{i \in I} E_i}
\Delta_{16, 6} &= \sum_{|S| = 6} \mathbb{P}(S \subseteq X) \sum_{I \subseteq \{1, \dots, 16\}, |I| \ge 1} (-1)^{|I|+1} \mathbf{1}_{S \subseteq \cup_{i \in I} E_i}
\Delta_{16, 7} &= \sum_{|S| = 7} \mathbb{P}(S \subseteq X) \sum_{I \subseteq \{1, \dots, 16\}, |I| \ge 1} (-1)^{|I|+1} \mathbf{1}_{S \subseteq \cup_{i \in I} E_i}
\Delta_{16, 8} &= \sum_{|S| = 8} \mathbb{P}(S \subseteq X) \sum_{I \subseteq \{1, \dots, 16\}, |I| \ge 1} (-1)^{|I|+1} \mathbf{1}_{S \subseteq \cup_{i \in I} E_i}
\Delta_{16, 9} &= \sum_{|S| = 9} \mathbb{P}(S \subseteq X) \sum_{I \subseteq \{1, \dots, 16\}, |I| \ge 1} (-1)^{|I|+1} \mathbf{1}_{S \subseteq \cup_{i \in I} E_i}\end{align}Cette majoration permet de borner supérieurement la corrélation locale. Par application du lemme de Lovász local ou des inégalités de concentration azumaiennes, la déviation par rapport à la moyenne produit un terme résiduel qui s'amortit exponentiellement. Nous fixons explicitement cette décroissance :
\begin{equation}\mathcal{R}_{16} \le \exp\left( - \frac{ \left( \Delta_{16, 1} \right)^2 }{ 2 \sum_{j=2}^{16} \Delta_{16, j} + \frac{1}{3} \max_{j} \Delta_{16, j} } \right)\end{equation}Le contrôle de ce terme $\mathcal{R}_{16}$ certifie l'indépendance quasi-totale des choix pour les structures de dimension 16.

\subsubsection{Analyse du moment d'ordre 17}Pour le degré 17, nous évaluons l'intersection simultanée de $\ell = 17$ hyperarêtes. Soient $E_1, \dots, E_{17}$ des arêtes de l'hypergraphe.
\begin{equation}\mathbb{E}\left[ \prod_{i=1}^{17} \mathbf{1}_{E_i \cap X \neq \emptyset} \right] = \sum_{j=1}^{17} (-1)^{j-1} \sum_{1 \le i_1 < \dots < i_j \le 17} \mathbb{P}\left( \bigcap_{m=1}^j (E_{i_m} \cap X \neq \emptyset) \right)\end{equation}Nous décomposons la probabilité conjointe en utilisant l'indépendance conditionnelle sur le choix des sommets de $X$. L'ensemble des configurations de recouvrement de l'union $\bigcup_{i} E_i$ est partitionné selon la taille de leur intersection avec $X$.
\begin{align} \\\Delta_{17, 1} &= \sum_{|S| = 1} \mathbb{P}(S \subseteq X) \sum_{I \subseteq \{1, \dots, 17\}, |I| \ge 1} (-1)^{|I|+1} \mathbf{1}_{S \subseteq \cup_{i \in I} E_i}
\Delta_{17, 2} &= \sum_{|S| = 2} \mathbb{P}(S \subseteq X) \sum_{I \subseteq \{1, \dots, 17\}, |I| \ge 1} (-1)^{|I|+1} \mathbf{1}_{S \subseteq \cup_{i \in I} E_i}
\Delta_{17, 3} &= \sum_{|S| = 3} \mathbb{P}(S \subseteq X) \sum_{I \subseteq \{1, \dots, 17\}, |I| \ge 1} (-1)^{|I|+1} \mathbf{1}_{S \subseteq \cup_{i \in I} E_i}
\Delta_{17, 4} &= \sum_{|S| = 4} \mathbb{P}(S \subseteq X) \sum_{I \subseteq \{1, \dots, 17\}, |I| \ge 1} (-1)^{|I|+1} \mathbf{1}_{S \subseteq \cup_{i \in I} E_i}
\Delta_{17, 5} &= \sum_{|S| = 5} \mathbb{P}(S \subseteq X) \sum_{I \subseteq \{1, \dots, 17\}, |I| \ge 1} (-1)^{|I|+1} \mathbf{1}_{S \subseteq \cup_{i \in I} E_i}
\Delta_{17, 6} &= \sum_{|S| = 6} \mathbb{P}(S \subseteq X) \sum_{I \subseteq \{1, \dots, 17\}, |I| \ge 1} (-1)^{|I|+1} \mathbf{1}_{S \subseteq \cup_{i \in I} E_i}
\Delta_{17, 7} &= \sum_{|S| = 7} \mathbb{P}(S \subseteq X) \sum_{I \subseteq \{1, \dots, 17\}, |I| \ge 1} (-1)^{|I|+1} \mathbf{1}_{S \subseteq \cup_{i \in I} E_i}
\Delta_{17, 8} &= \sum_{|S| = 8} \mathbb{P}(S \subseteq X) \sum_{I \subseteq \{1, \dots, 17\}, |I| \ge 1} (-1)^{|I|+1} \mathbf{1}_{S \subseteq \cup_{i \in I} E_i}
\Delta_{17, 9} &= \sum_{|S| = 9} \mathbb{P}(S \subseteq X) \sum_{I \subseteq \{1, \dots, 17\}, |I| \ge 1} (-1)^{|I|+1} \mathbf{1}_{S \subseteq \cup_{i \in I} E_i}\end{align}Cette majoration permet de borner supérieurement la corrélation locale. Par application du lemme de Lovász local ou des inégalités de concentration azumaiennes, la déviation par rapport à la moyenne produit un terme résiduel qui s'amortit exponentiellement. Nous fixons explicitement cette décroissance :
\begin{equation}\mathcal{R}_{17} \le \exp\left( - \frac{ \left( \Delta_{17, 1} \right)^2 }{ 2 \sum_{j=2}^{17} \Delta_{17, j} + \frac{1}{3} \max_{j} \Delta_{17, j} } \right)\end{equation}Le contrôle de ce terme $\mathcal{R}_{17}$ certifie l'indépendance quasi-totale des choix pour les structures de dimension 17.

\subsubsection{Analyse du moment d'ordre 18}Pour le degré 18, nous évaluons l'intersection simultanée de $\ell = 18$ hyperarêtes. Soient $E_1, \dots, E_{18}$ des arêtes de l'hypergraphe.
\begin{equation}\mathbb{E}\left[ \prod_{i=1}^{18} \mathbf{1}_{E_i \cap X \neq \emptyset} \right] = \sum_{j=1}^{18} (-1)^{j-1} \sum_{1 \le i_1 < \dots < i_j \le 18} \mathbb{P}\left( \bigcap_{m=1}^j (E_{i_m} \cap X \neq \emptyset) \right)\end{equation}Nous décomposons la probabilité conjointe en utilisant l'indépendance conditionnelle sur le choix des sommets de $X$. L'ensemble des configurations de recouvrement de l'union $\bigcup_{i} E_i$ est partitionné selon la taille de leur intersection avec $X$.
\begin{align} \\\Delta_{18, 1} &= \sum_{|S| = 1} \mathbb{P}(S \subseteq X) \sum_{I \subseteq \{1, \dots, 18\}, |I| \ge 1} (-1)^{|I|+1} \mathbf{1}_{S \subseteq \cup_{i \in I} E_i}
\Delta_{18, 2} &= \sum_{|S| = 2} \mathbb{P}(S \subseteq X) \sum_{I \subseteq \{1, \dots, 18\}, |I| \ge 1} (-1)^{|I|+1} \mathbf{1}_{S \subseteq \cup_{i \in I} E_i}
\Delta_{18, 3} &= \sum_{|S| = 3} \mathbb{P}(S \subseteq X) \sum_{I \subseteq \{1, \dots, 18\}, |I| \ge 1} (-1)^{|I|+1} \mathbf{1}_{S \subseteq \cup_{i \in I} E_i}
\Delta_{18, 4} &= \sum_{|S| = 4} \mathbb{P}(S \subseteq X) \sum_{I \subseteq \{1, \dots, 18\}, |I| \ge 1} (-1)^{|I|+1} \mathbf{1}_{S \subseteq \cup_{i \in I} E_i}
\Delta_{18, 5} &= \sum_{|S| = 5} \mathbb{P}(S \subseteq X) \sum_{I \subseteq \{1, \dots, 18\}, |I| \ge 1} (-1)^{|I|+1} \mathbf{1}_{S \subseteq \cup_{i \in I} E_i}
\Delta_{18, 6} &= \sum_{|S| = 6} \mathbb{P}(S \subseteq X) \sum_{I \subseteq \{1, \dots, 18\}, |I| \ge 1} (-1)^{|I|+1} \mathbf{1}_{S \subseteq \cup_{i \in I} E_i}
\Delta_{18, 7} &= \sum_{|S| = 7} \mathbb{P}(S \subseteq X) \sum_{I \subseteq \{1, \dots, 18\}, |I| \ge 1} (-1)^{|I|+1} \mathbf{1}_{S \subseteq \cup_{i \in I} E_i}
\Delta_{18, 8} &= \sum_{|S| = 8} \mathbb{P}(S \subseteq X) \sum_{I \subseteq \{1, \dots, 18\}, |I| \ge 1} (-1)^{|I|+1} \mathbf{1}_{S \subseteq \cup_{i \in I} E_i}
\Delta_{18, 9} &= \sum_{|S| = 9} \mathbb{P}(S \subseteq X) \sum_{I \subseteq \{1, \dots, 18\}, |I| \ge 1} (-1)^{|I|+1} \mathbf{1}_{S \subseteq \cup_{i \in I} E_i}\end{align}Cette majoration permet de borner supérieurement la corrélation locale. Par application du lemme de Lovász local ou des inégalités de concentration azumaiennes, la déviation par rapport à la moyenne produit un terme résiduel qui s'amortit exponentiellement. Nous fixons explicitement cette décroissance :
\begin{equation}\mathcal{R}_{18} \le \exp\left( - \frac{ \left( \Delta_{18, 1} \right)^2 }{ 2 \sum_{j=2}^{18} \Delta_{18, j} + \frac{1}{3} \max_{j} \Delta_{18, j} } \right)\end{equation}Le contrôle de ce terme $\mathcal{R}_{18}$ certifie l'indépendance quasi-totale des choix pour les structures de dimension 18.

\subsubsection{Analyse du moment d'ordre 19}Pour le degré 19, nous évaluons l'intersection simultanée de $\ell = 19$ hyperarêtes. Soient $E_1, \dots, E_{19}$ des arêtes de l'hypergraphe.
\begin{equation}\mathbb{E}\left[ \prod_{i=1}^{19} \mathbf{1}_{E_i \cap X \neq \emptyset} \right] = \sum_{j=1}^{19} (-1)^{j-1} \sum_{1 \le i_1 < \dots < i_j \le 19} \mathbb{P}\left( \bigcap_{m=1}^j (E_{i_m} \cap X \neq \emptyset) \right)\end{equation}Nous décomposons la probabilité conjointe en utilisant l'indépendance conditionnelle sur le choix des sommets de $X$. L'ensemble des configurations de recouvrement de l'union $\bigcup_{i} E_i$ est partitionné selon la taille de leur intersection avec $X$.
\begin{align} \\\Delta_{19, 1} &= \sum_{|S| = 1} \mathbb{P}(S \subseteq X) \sum_{I \subseteq \{1, \dots, 19\}, |I| \ge 1} (-1)^{|I|+1} \mathbf{1}_{S \subseteq \cup_{i \in I} E_i}
\Delta_{19, 2} &= \sum_{|S| = 2} \mathbb{P}(S \subseteq X) \sum_{I \subseteq \{1, \dots, 19\}, |I| \ge 1} (-1)^{|I|+1} \mathbf{1}_{S \subseteq \cup_{i \in I} E_i}
\Delta_{19, 3} &= \sum_{|S| = 3} \mathbb{P}(S \subseteq X) \sum_{I \subseteq \{1, \dots, 19\}, |I| \ge 1} (-1)^{|I|+1} \mathbf{1}_{S \subseteq \cup_{i \in I} E_i}
\Delta_{19, 4} &= \sum_{|S| = 4} \mathbb{P}(S \subseteq X) \sum_{I \subseteq \{1, \dots, 19\}, |I| \ge 1} (-1)^{|I|+1} \mathbf{1}_{S \subseteq \cup_{i \in I} E_i}
\Delta_{19, 5} &= \sum_{|S| = 5} \mathbb{P}(S \subseteq X) \sum_{I \subseteq \{1, \dots, 19\}, |I| \ge 1} (-1)^{|I|+1} \mathbf{1}_{S \subseteq \cup_{i \in I} E_i}
\Delta_{19, 6} &= \sum_{|S| = 6} \mathbb{P}(S \subseteq X) \sum_{I \subseteq \{1, \dots, 19\}, |I| \ge 1} (-1)^{|I|+1} \mathbf{1}_{S \subseteq \cup_{i \in I} E_i}
\Delta_{19, 7} &= \sum_{|S| = 7} \mathbb{P}(S \subseteq X) \sum_{I \subseteq \{1, \dots, 19\}, |I| \ge 1} (-1)^{|I|+1} \mathbf{1}_{S \subseteq \cup_{i \in I} E_i}
\Delta_{19, 8} &= \sum_{|S| = 8} \mathbb{P}(S \subseteq X) \sum_{I \subseteq \{1, \dots, 19\}, |I| \ge 1} (-1)^{|I|+1} \mathbf{1}_{S \subseteq \cup_{i \in I} E_i}
\Delta_{19, 9} &= \sum_{|S| = 9} \mathbb{P}(S \subseteq X) \sum_{I \subseteq \{1, \dots, 19\}, |I| \ge 1} (-1)^{|I|+1} \mathbf{1}_{S \subseteq \cup_{i \in I} E_i}\end{align}Cette majoration permet de borner supérieurement la corrélation locale. Par application du lemme de Lovász local ou des inégalités de concentration azumaiennes, la déviation par rapport à la moyenne produit un terme résiduel qui s'amortit exponentiellement. Nous fixons explicitement cette décroissance :
\begin{equation}\mathcal{R}_{19} \le \exp\left( - \frac{ \left( \Delta_{19, 1} \right)^2 }{ 2 \sum_{j=2}^{19} \Delta_{19, j} + \frac{1}{3} \max_{j} \Delta_{19, j} } \right)\end{equation}Le contrôle de ce terme $\mathcal{R}_{19}$ certifie l'indépendance quasi-totale des choix pour les structures de dimension 19.

\subsubsection{Analyse du moment d'ordre 20}Pour le degré 20, nous évaluons l'intersection simultanée de $\ell = 20$ hyperarêtes. Soient $E_1, \dots, E_{20}$ des arêtes de l'hypergraphe.
\begin{equation}\mathbb{E}\left[ \prod_{i=1}^{20} \mathbf{1}_{E_i \cap X \neq \emptyset} \right] = \sum_{j=1}^{20} (-1)^{j-1} \sum_{1 \le i_1 < \dots < i_j \le 20} \mathbb{P}\left( \bigcap_{m=1}^j (E_{i_m} \cap X \neq \emptyset) \right)\end{equation}Nous décomposons la probabilité conjointe en utilisant l'indépendance conditionnelle sur le choix des sommets de $X$. L'ensemble des configurations de recouvrement de l'union $\bigcup_{i} E_i$ est partitionné selon la taille de leur intersection avec $X$.
\begin{align} \\\Delta_{20, 1} &= \sum_{|S| = 1} \mathbb{P}(S \subseteq X) \sum_{I \subseteq \{1, \dots, 20\}, |I| \ge 1} (-1)^{|I|+1} \mathbf{1}_{S \subseteq \cup_{i \in I} E_i}
\Delta_{20, 2} &= \sum_{|S| = 2} \mathbb{P}(S \subseteq X) \sum_{I \subseteq \{1, \dots, 20\}, |I| \ge 1} (-1)^{|I|+1} \mathbf{1}_{S \subseteq \cup_{i \in I} E_i}
\Delta_{20, 3} &= \sum_{|S| = 3} \mathbb{P}(S \subseteq X) \sum_{I \subseteq \{1, \dots, 20\}, |I| \ge 1} (-1)^{|I|+1} \mathbf{1}_{S \subseteq \cup_{i \in I} E_i}
\Delta_{20, 4} &= \sum_{|S| = 4} \mathbb{P}(S \subseteq X) \sum_{I \subseteq \{1, \dots, 20\}, |I| \ge 1} (-1)^{|I|+1} \mathbf{1}_{S \subseteq \cup_{i \in I} E_i}
\Delta_{20, 5} &= \sum_{|S| = 5} \mathbb{P}(S \subseteq X) \sum_{I \subseteq \{1, \dots, 20\}, |I| \ge 1} (-1)^{|I|+1} \mathbf{1}_{S \subseteq \cup_{i \in I} E_i}
\Delta_{20, 6} &= \sum_{|S| = 6} \mathbb{P}(S \subseteq X) \sum_{I \subseteq \{1, \dots, 20\}, |I| \ge 1} (-1)^{|I|+1} \mathbf{1}_{S \subseteq \cup_{i \in I} E_i}
\Delta_{20, 7} &= \sum_{|S| = 7} \mathbb{P}(S \subseteq X) \sum_{I \subseteq \{1, \dots, 20\}, |I| \ge 1} (-1)^{|I|+1} \mathbf{1}_{S \subseteq \cup_{i \in I} E_i}
\Delta_{20, 8} &= \sum_{|S| = 8} \mathbb{P}(S \subseteq X) \sum_{I \subseteq \{1, \dots, 20\}, |I| \ge 1} (-1)^{|I|+1} \mathbf{1}_{S \subseteq \cup_{i \in I} E_i}
\Delta_{20, 9} &= \sum_{|S| = 9} \mathbb{P}(S \subseteq X) \sum_{I \subseteq \{1, \dots, 20\}, |I| \ge 1} (-1)^{|I|+1} \mathbf{1}_{S \subseteq \cup_{i \in I} E_i}\end{align}Cette majoration permet de borner supérieurement la corrélation locale. Par application du lemme de Lovász local ou des inégalités de concentration azumaiennes, la déviation par rapport à la moyenne produit un terme résiduel qui s'amortit exponentiellement. Nous fixons explicitement cette décroissance :
\begin{equation}\mathcal{R}_{20} \le \exp\left( - \frac{ \left( \Delta_{20, 1} \right)^2 }{ 2 \sum_{j=2}^{20} \Delta_{20, j} + \frac{1}{3} \max_{j} \Delta_{20, j} } \right)\end{equation}Le contrôle de ce terme $\mathcal{R}_{20}$ certifie l'indépendance quasi-totale des choix pour les structures de dimension 20.

\subsubsection{Analyse du moment d'ordre 21}Pour le degré 21, nous évaluons l'intersection simultanée de $\ell = 21$ hyperarêtes. Soient $E_1, \dots, E_{21}$ des arêtes de l'hypergraphe.
\begin{equation}\mathbb{E}\left[ \prod_{i=1}^{21} \mathbf{1}_{E_i \cap X \neq \emptyset} \right] = \sum_{j=1}^{21} (-1)^{j-1} \sum_{1 \le i_1 < \dots < i_j \le 21} \mathbb{P}\left( \bigcap_{m=1}^j (E_{i_m} \cap X \neq \emptyset) \right)\end{equation}Nous décomposons la probabilité conjointe en utilisant l'indépendance conditionnelle sur le choix des sommets de $X$. L'ensemble des configurations de recouvrement de l'union $\bigcup_{i} E_i$ est partitionné selon la taille de leur intersection avec $X$.
\begin{align} \\\Delta_{21, 1} &= \sum_{|S| = 1} \mathbb{P}(S \subseteq X) \sum_{I \subseteq \{1, \dots, 21\}, |I| \ge 1} (-1)^{|I|+1} \mathbf{1}_{S \subseteq \cup_{i \in I} E_i}
\Delta_{21, 2} &= \sum_{|S| = 2} \mathbb{P}(S \subseteq X) \sum_{I \subseteq \{1, \dots, 21\}, |I| \ge 1} (-1)^{|I|+1} \mathbf{1}_{S \subseteq \cup_{i \in I} E_i}
\Delta_{21, 3} &= \sum_{|S| = 3} \mathbb{P}(S \subseteq X) \sum_{I \subseteq \{1, \dots, 21\}, |I| \ge 1} (-1)^{|I|+1} \mathbf{1}_{S \subseteq \cup_{i \in I} E_i}
\Delta_{21, 4} &= \sum_{|S| = 4} \mathbb{P}(S \subseteq X) \sum_{I \subseteq \{1, \dots, 21\}, |I| \ge 1} (-1)^{|I|+1} \mathbf{1}_{S \subseteq \cup_{i \in I} E_i}
\Delta_{21, 5} &= \sum_{|S| = 5} \mathbb{P}(S \subseteq X) \sum_{I \subseteq \{1, \dots, 21\}, |I| \ge 1} (-1)^{|I|+1} \mathbf{1}_{S \subseteq \cup_{i \in I} E_i}
\Delta_{21, 6} &= \sum_{|S| = 6} \mathbb{P}(S \subseteq X) \sum_{I \subseteq \{1, \dots, 21\}, |I| \ge 1} (-1)^{|I|+1} \mathbf{1}_{S \subseteq \cup_{i \in I} E_i}
\Delta_{21, 7} &= \sum_{|S| = 7} \mathbb{P}(S \subseteq X) \sum_{I \subseteq \{1, \dots, 21\}, |I| \ge 1} (-1)^{|I|+1} \mathbf{1}_{S \subseteq \cup_{i \in I} E_i}
\Delta_{21, 8} &= \sum_{|S| = 8} \mathbb{P}(S \subseteq X) \sum_{I \subseteq \{1, \dots, 21\}, |I| \ge 1} (-1)^{|I|+1} \mathbf{1}_{S \subseteq \cup_{i \in I} E_i}
\Delta_{21, 9} &= \sum_{|S| = 9} \mathbb{P}(S \subseteq X) \sum_{I \subseteq \{1, \dots, 21\}, |I| \ge 1} (-1)^{|I|+1} \mathbf{1}_{S \subseteq \cup_{i \in I} E_i}\end{align}Cette majoration permet de borner supérieurement la corrélation locale. Par application du lemme de Lovász local ou des inégalités de concentration azumaiennes, la déviation par rapport à la moyenne produit un terme résiduel qui s'amortit exponentiellement. Nous fixons explicitement cette décroissance :
\begin{equation}\mathcal{R}_{21} \le \exp\left( - \frac{ \left( \Delta_{21, 1} \right)^2 }{ 2 \sum_{j=2}^{21} \Delta_{21, j} + \frac{1}{3} \max_{j} \Delta_{21, j} } \right)\end{equation}Le contrôle de ce terme $\mathcal{R}_{21}$ certifie l'indépendance quasi-totale des choix pour les structures de dimension 21.

\subsubsection{Analyse du moment d'ordre 22}Pour le degré 22, nous évaluons l'intersection simultanée de $\ell = 22$ hyperarêtes. Soient $E_1, \dots, E_{22}$ des arêtes de l'hypergraphe.
\begin{equation}\mathbb{E}\left[ \prod_{i=1}^{22} \mathbf{1}_{E_i \cap X \neq \emptyset} \right] = \sum_{j=1}^{22} (-1)^{j-1} \sum_{1 \le i_1 < \dots < i_j \le 22} \mathbb{P}\left( \bigcap_{m=1}^j (E_{i_m} \cap X \neq \emptyset) \right)\end{equation}Nous décomposons la probabilité conjointe en utilisant l'indépendance conditionnelle sur le choix des sommets de $X$. L'ensemble des configurations de recouvrement de l'union $\bigcup_{i} E_i$ est partitionné selon la taille de leur intersection avec $X$.
\begin{align} \\\Delta_{22, 1} &= \sum_{|S| = 1} \mathbb{P}(S \subseteq X) \sum_{I \subseteq \{1, \dots, 22\}, |I| \ge 1} (-1)^{|I|+1} \mathbf{1}_{S \subseteq \cup_{i \in I} E_i}
\Delta_{22, 2} &= \sum_{|S| = 2} \mathbb{P}(S \subseteq X) \sum_{I \subseteq \{1, \dots, 22\}, |I| \ge 1} (-1)^{|I|+1} \mathbf{1}_{S \subseteq \cup_{i \in I} E_i}
\Delta_{22, 3} &= \sum_{|S| = 3} \mathbb{P}(S \subseteq X) \sum_{I \subseteq \{1, \dots, 22\}, |I| \ge 1} (-1)^{|I|+1} \mathbf{1}_{S \subseteq \cup_{i \in I} E_i}
\Delta_{22, 4} &= \sum_{|S| = 4} \mathbb{P}(S \subseteq X) \sum_{I \subseteq \{1, \dots, 22\}, |I| \ge 1} (-1)^{|I|+1} \mathbf{1}_{S \subseteq \cup_{i \in I} E_i}
\Delta_{22, 5} &= \sum_{|S| = 5} \mathbb{P}(S \subseteq X) \sum_{I \subseteq \{1, \dots, 22\}, |I| \ge 1} (-1)^{|I|+1} \mathbf{1}_{S \subseteq \cup_{i \in I} E_i}
\Delta_{22, 6} &= \sum_{|S| = 6} \mathbb{P}(S \subseteq X) \sum_{I \subseteq \{1, \dots, 22\}, |I| \ge 1} (-1)^{|I|+1} \mathbf{1}_{S \subseteq \cup_{i \in I} E_i}
\Delta_{22, 7} &= \sum_{|S| = 7} \mathbb{P}(S \subseteq X) \sum_{I \subseteq \{1, \dots, 22\}, |I| \ge 1} (-1)^{|I|+1} \mathbf{1}_{S \subseteq \cup_{i \in I} E_i}
\Delta_{22, 8} &= \sum_{|S| = 8} \mathbb{P}(S \subseteq X) \sum_{I \subseteq \{1, \dots, 22\}, |I| \ge 1} (-1)^{|I|+1} \mathbf{1}_{S \subseteq \cup_{i \in I} E_i}
\Delta_{22, 9} &= \sum_{|S| = 9} \mathbb{P}(S \subseteq X) \sum_{I \subseteq \{1, \dots, 22\}, |I| \ge 1} (-1)^{|I|+1} \mathbf{1}_{S \subseteq \cup_{i \in I} E_i}\end{align}Cette majoration permet de borner supérieurement la corrélation locale. Par application du lemme de Lovász local ou des inégalités de concentration azumaiennes, la déviation par rapport à la moyenne produit un terme résiduel qui s'amortit exponentiellement. Nous fixons explicitement cette décroissance :
\begin{equation}\mathcal{R}_{22} \le \exp\left( - \frac{ \left( \Delta_{22, 1} \right)^2 }{ 2 \sum_{j=2}^{22} \Delta_{22, j} + \frac{1}{3} \max_{j} \Delta_{22, j} } \right)\end{equation}Le contrôle de ce terme $\mathcal{R}_{22}$ certifie l'indépendance quasi-totale des choix pour les structures de dimension 22.

\subsubsection{Analyse du moment d'ordre 23}Pour le degré 23, nous évaluons l'intersection simultanée de $\ell = 23$ hyperarêtes. Soient $E_1, \dots, E_{23}$ des arêtes de l'hypergraphe.
\begin{equation}\mathbb{E}\left[ \prod_{i=1}^{23} \mathbf{1}_{E_i \cap X \neq \emptyset} \right] = \sum_{j=1}^{23} (-1)^{j-1} \sum_{1 \le i_1 < \dots < i_j \le 23} \mathbb{P}\left( \bigcap_{m=1}^j (E_{i_m} \cap X \neq \emptyset) \right)\end{equation}Nous décomposons la probabilité conjointe en utilisant l'indépendance conditionnelle sur le choix des sommets de $X$. L'ensemble des configurations de recouvrement de l'union $\bigcup_{i} E_i$ est partitionné selon la taille de leur intersection avec $X$.
\begin{align} \\\Delta_{23, 1} &= \sum_{|S| = 1} \mathbb{P}(S \subseteq X) \sum_{I \subseteq \{1, \dots, 23\}, |I| \ge 1} (-1)^{|I|+1} \mathbf{1}_{S \subseteq \cup_{i \in I} E_i}
\Delta_{23, 2} &= \sum_{|S| = 2} \mathbb{P}(S \subseteq X) \sum_{I \subseteq \{1, \dots, 23\}, |I| \ge 1} (-1)^{|I|+1} \mathbf{1}_{S \subseteq \cup_{i \in I} E_i}
\Delta_{23, 3} &= \sum_{|S| = 3} \mathbb{P}(S \subseteq X) \sum_{I \subseteq \{1, \dots, 23\}, |I| \ge 1} (-1)^{|I|+1} \mathbf{1}_{S \subseteq \cup_{i \in I} E_i}
\Delta_{23, 4} &= \sum_{|S| = 4} \mathbb{P}(S \subseteq X) \sum_{I \subseteq \{1, \dots, 23\}, |I| \ge 1} (-1)^{|I|+1} \mathbf{1}_{S \subseteq \cup_{i \in I} E_i}
\Delta_{23, 5} &= \sum_{|S| = 5} \mathbb{P}(S \subseteq X) \sum_{I \subseteq \{1, \dots, 23\}, |I| \ge 1} (-1)^{|I|+1} \mathbf{1}_{S \subseteq \cup_{i \in I} E_i}
\Delta_{23, 6} &= \sum_{|S| = 6} \mathbb{P}(S \subseteq X) \sum_{I \subseteq \{1, \dots, 23\}, |I| \ge 1} (-1)^{|I|+1} \mathbf{1}_{S \subseteq \cup_{i \in I} E_i}
\Delta_{23, 7} &= \sum_{|S| = 7} \mathbb{P}(S \subseteq X) \sum_{I \subseteq \{1, \dots, 23\}, |I| \ge 1} (-1)^{|I|+1} \mathbf{1}_{S \subseteq \cup_{i \in I} E_i}
\Delta_{23, 8} &= \sum_{|S| = 8} \mathbb{P}(S \subseteq X) \sum_{I \subseteq \{1, \dots, 23\}, |I| \ge 1} (-1)^{|I|+1} \mathbf{1}_{S \subseteq \cup_{i \in I} E_i}
\Delta_{23, 9} &= \sum_{|S| = 9} \mathbb{P}(S \subseteq X) \sum_{I \subseteq \{1, \dots, 23\}, |I| \ge 1} (-1)^{|I|+1} \mathbf{1}_{S \subseteq \cup_{i \in I} E_i}\end{align}Cette majoration permet de borner supérieurement la corrélation locale. Par application du lemme de Lovász local ou des inégalités de concentration azumaiennes, la déviation par rapport à la moyenne produit un terme résiduel qui s'amortit exponentiellement. Nous fixons explicitement cette décroissance :
\begin{equation}\mathcal{R}_{23} \le \exp\left( - \frac{ \left( \Delta_{23, 1} \right)^2 }{ 2 \sum_{j=2}^{23} \Delta_{23, j} + \frac{1}{3} \max_{j} \Delta_{23, j} } \right)\end{equation}Le contrôle de ce terme $\mathcal{R}_{23}$ certifie l'indépendance quasi-totale des choix pour les structures de dimension 23.

\subsubsection{Analyse du moment d'ordre 24}Pour le degré 24, nous évaluons l'intersection simultanée de $\ell = 24$ hyperarêtes. Soient $E_1, \dots, E_{24}$ des arêtes de l'hypergraphe.
\begin{equation}\mathbb{E}\left[ \prod_{i=1}^{24} \mathbf{1}_{E_i \cap X \neq \emptyset} \right] = \sum_{j=1}^{24} (-1)^{j-1} \sum_{1 \le i_1 < \dots < i_j \le 24} \mathbb{P}\left( \bigcap_{m=1}^j (E_{i_m} \cap X \neq \emptyset) \right)\end{equation}Nous décomposons la probabilité conjointe en utilisant l'indépendance conditionnelle sur le choix des sommets de $X$. L'ensemble des configurations de recouvrement de l'union $\bigcup_{i} E_i$ est partitionné selon la taille de leur intersection avec $X$.
\begin{align} \\\Delta_{24, 1} &= \sum_{|S| = 1} \mathbb{P}(S \subseteq X) \sum_{I \subseteq \{1, \dots, 24\}, |I| \ge 1} (-1)^{|I|+1} \mathbf{1}_{S \subseteq \cup_{i \in I} E_i}
\Delta_{24, 2} &= \sum_{|S| = 2} \mathbb{P}(S \subseteq X) \sum_{I \subseteq \{1, \dots, 24\}, |I| \ge 1} (-1)^{|I|+1} \mathbf{1}_{S \subseteq \cup_{i \in I} E_i}
\Delta_{24, 3} &= \sum_{|S| = 3} \mathbb{P}(S \subseteq X) \sum_{I \subseteq \{1, \dots, 24\}, |I| \ge 1} (-1)^{|I|+1} \mathbf{1}_{S \subseteq \cup_{i \in I} E_i}
\Delta_{24, 4} &= \sum_{|S| = 4} \mathbb{P}(S \subseteq X) \sum_{I \subseteq \{1, \dots, 24\}, |I| \ge 1} (-1)^{|I|+1} \mathbf{1}_{S \subseteq \cup_{i \in I} E_i}
\Delta_{24, 5} &= \sum_{|S| = 5} \mathbb{P}(S \subseteq X) \sum_{I \subseteq \{1, \dots, 24\}, |I| \ge 1} (-1)^{|I|+1} \mathbf{1}_{S \subseteq \cup_{i \in I} E_i}
\Delta_{24, 6} &= \sum_{|S| = 6} \mathbb{P}(S \subseteq X) \sum_{I \subseteq \{1, \dots, 24\}, |I| \ge 1} (-1)^{|I|+1} \mathbf{1}_{S \subseteq \cup_{i \in I} E_i}
\Delta_{24, 7} &= \sum_{|S| = 7} \mathbb{P}(S \subseteq X) \sum_{I \subseteq \{1, \dots, 24\}, |I| \ge 1} (-1)^{|I|+1} \mathbf{1}_{S \subseteq \cup_{i \in I} E_i}
\Delta_{24, 8} &= \sum_{|S| = 8} \mathbb{P}(S \subseteq X) \sum_{I \subseteq \{1, \dots, 24\}, |I| \ge 1} (-1)^{|I|+1} \mathbf{1}_{S \subseteq \cup_{i \in I} E_i}
\Delta_{24, 9} &= \sum_{|S| = 9} \mathbb{P}(S \subseteq X) \sum_{I \subseteq \{1, \dots, 24\}, |I| \ge 1} (-1)^{|I|+1} \mathbf{1}_{S \subseteq \cup_{i \in I} E_i}\end{align}Cette majoration permet de borner supérieurement la corrélation locale. Par application du lemme de Lovász local ou des inégalités de concentration azumaiennes, la déviation par rapport à la moyenne produit un terme résiduel qui s'amortit exponentiellement. Nous fixons explicitement cette décroissance :
\begin{equation}\mathcal{R}_{24} \le \exp\left( - \frac{ \left( \Delta_{24, 1} \right)^2 }{ 2 \sum_{j=2}^{24} \Delta_{24, j} + \frac{1}{3} \max_{j} \Delta_{24, j} } \right)\end{equation}Le contrôle de ce terme $\mathcal{R}_{24}$ certifie l'indépendance quasi-totale des choix pour les structures de dimension 24.

\subsubsection{Analyse du moment d'ordre 25}Pour le degré 25, nous évaluons l'intersection simultanée de $\ell = 25$ hyperarêtes. Soient $E_1, \dots, E_{25}$ des arêtes de l'hypergraphe.
\begin{equation}\mathbb{E}\left[ \prod_{i=1}^{25} \mathbf{1}_{E_i \cap X \neq \emptyset} \right] = \sum_{j=1}^{25} (-1)^{j-1} \sum_{1 \le i_1 < \dots < i_j \le 25} \mathbb{P}\left( \bigcap_{m=1}^j (E_{i_m} \cap X \neq \emptyset) \right)\end{equation}Nous décomposons la probabilité conjointe en utilisant l'indépendance conditionnelle sur le choix des sommets de $X$. L'ensemble des configurations de recouvrement de l'union $\bigcup_{i} E_i$ est partitionné selon la taille de leur intersection avec $X$.
\begin{align} \\\Delta_{25, 1} &= \sum_{|S| = 1} \mathbb{P}(S \subseteq X) \sum_{I \subseteq \{1, \dots, 25\}, |I| \ge 1} (-1)^{|I|+1} \mathbf{1}_{S \subseteq \cup_{i \in I} E_i}
\Delta_{25, 2} &= \sum_{|S| = 2} \mathbb{P}(S \subseteq X) \sum_{I \subseteq \{1, \dots, 25\}, |I| \ge 1} (-1)^{|I|+1} \mathbf{1}_{S \subseteq \cup_{i \in I} E_i}
\Delta_{25, 3} &= \sum_{|S| = 3} \mathbb{P}(S \subseteq X) \sum_{I \subseteq \{1, \dots, 25\}, |I| \ge 1} (-1)^{|I|+1} \mathbf{1}_{S \subseteq \cup_{i \in I} E_i}
\Delta_{25, 4} &= \sum_{|S| = 4} \mathbb{P}(S \subseteq X) \sum_{I \subseteq \{1, \dots, 25\}, |I| \ge 1} (-1)^{|I|+1} \mathbf{1}_{S \subseteq \cup_{i \in I} E_i}
\Delta_{25, 5} &= \sum_{|S| = 5} \mathbb{P}(S \subseteq X) \sum_{I \subseteq \{1, \dots, 25\}, |I| \ge 1} (-1)^{|I|+1} \mathbf{1}_{S \subseteq \cup_{i \in I} E_i}
\Delta_{25, 6} &= \sum_{|S| = 6} \mathbb{P}(S \subseteq X) \sum_{I \subseteq \{1, \dots, 25\}, |I| \ge 1} (-1)^{|I|+1} \mathbf{1}_{S \subseteq \cup_{i \in I} E_i}
\Delta_{25, 7} &= \sum_{|S| = 7} \mathbb{P}(S \subseteq X) \sum_{I \subseteq \{1, \dots, 25\}, |I| \ge 1} (-1)^{|I|+1} \mathbf{1}_{S \subseteq \cup_{i \in I} E_i}
\Delta_{25, 8} &= \sum_{|S| = 8} \mathbb{P}(S \subseteq X) \sum_{I \subseteq \{1, \dots, 25\}, |I| \ge 1} (-1)^{|I|+1} \mathbf{1}_{S \subseteq \cup_{i \in I} E_i}
\Delta_{25, 9} &= \sum_{|S| = 9} \mathbb{P}(S \subseteq X) \sum_{I \subseteq \{1, \dots, 25\}, |I| \ge 1} (-1)^{|I|+1} \mathbf{1}_{S \subseteq \cup_{i \in I} E_i}\end{align}Cette majoration permet de borner supérieurement la corrélation locale. Par application du lemme de Lovász local ou des inégalités de concentration azumaiennes, la déviation par rapport à la moyenne produit un terme résiduel qui s'amortit exponentiellement. Nous fixons explicitement cette décroissance :
\begin{equation}\mathcal{R}_{25} \le \exp\left( - \frac{ \left( \Delta_{25, 1} \right)^2 }{ 2 \sum_{j=2}^{25} \Delta_{25, j} + \frac{1}{3} \max_{j} \Delta_{25, j} } \right)\end{equation}Le contrôle de ce terme $\mathcal{R}_{25}$ certifie l'indépendance quasi-totale des choix pour les structures de dimension 25.

\subsubsection{Analyse du moment d'ordre 26}Pour le degré 26, nous évaluons l'intersection simultanée de $\ell = 26$ hyperarêtes. Soient $E_1, \dots, E_{26}$ des arêtes de l'hypergraphe.
\begin{equation}\mathbb{E}\left[ \prod_{i=1}^{26} \mathbf{1}_{E_i \cap X \neq \emptyset} \right] = \sum_{j=1}^{26} (-1)^{j-1} \sum_{1 \le i_1 < \dots < i_j \le 26} \mathbb{P}\left( \bigcap_{m=1}^j (E_{i_m} \cap X \neq \emptyset) \right)\end{equation}Nous décomposons la probabilité conjointe en utilisant l'indépendance conditionnelle sur le choix des sommets de $X$. L'ensemble des configurations de recouvrement de l'union $\bigcup_{i} E_i$ est partitionné selon la taille de leur intersection avec $X$.
\begin{align} \\\Delta_{26, 1} &= \sum_{|S| = 1} \mathbb{P}(S \subseteq X) \sum_{I \subseteq \{1, \dots, 26\}, |I| \ge 1} (-1)^{|I|+1} \mathbf{1}_{S \subseteq \cup_{i \in I} E_i}
\Delta_{26, 2} &= \sum_{|S| = 2} \mathbb{P}(S \subseteq X) \sum_{I \subseteq \{1, \dots, 26\}, |I| \ge 1} (-1)^{|I|+1} \mathbf{1}_{S \subseteq \cup_{i \in I} E_i}
\Delta_{26, 3} &= \sum_{|S| = 3} \mathbb{P}(S \subseteq X) \sum_{I \subseteq \{1, \dots, 26\}, |I| \ge 1} (-1)^{|I|+1} \mathbf{1}_{S \subseteq \cup_{i \in I} E_i}
\Delta_{26, 4} &= \sum_{|S| = 4} \mathbb{P}(S \subseteq X) \sum_{I \subseteq \{1, \dots, 26\}, |I| \ge 1} (-1)^{|I|+1} \mathbf{1}_{S \subseteq \cup_{i \in I} E_i}
\Delta_{26, 5} &= \sum_{|S| = 5} \mathbb{P}(S \subseteq X) \sum_{I \subseteq \{1, \dots, 26\}, |I| \ge 1} (-1)^{|I|+1} \mathbf{1}_{S \subseteq \cup_{i \in I} E_i}
\Delta_{26, 6} &= \sum_{|S| = 6} \mathbb{P}(S \subseteq X) \sum_{I \subseteq \{1, \dots, 26\}, |I| \ge 1} (-1)^{|I|+1} \mathbf{1}_{S \subseteq \cup_{i \in I} E_i}
\Delta_{26, 7} &= \sum_{|S| = 7} \mathbb{P}(S \subseteq X) \sum_{I \subseteq \{1, \dots, 26\}, |I| \ge 1} (-1)^{|I|+1} \mathbf{1}_{S \subseteq \cup_{i \in I} E_i}
\Delta_{26, 8} &= \sum_{|S| = 8} \mathbb{P}(S \subseteq X) \sum_{I \subseteq \{1, \dots, 26\}, |I| \ge 1} (-1)^{|I|+1} \mathbf{1}_{S \subseteq \cup_{i \in I} E_i}
\Delta_{26, 9} &= \sum_{|S| = 9} \mathbb{P}(S \subseteq X) \sum_{I \subseteq \{1, \dots, 26\}, |I| \ge 1} (-1)^{|I|+1} \mathbf{1}_{S \subseteq \cup_{i \in I} E_i}\end{align}Cette majoration permet de borner supérieurement la corrélation locale. Par application du lemme de Lovász local ou des inégalités de concentration azumaiennes, la déviation par rapport à la moyenne produit un terme résiduel qui s'amortit exponentiellement. Nous fixons explicitement cette décroissance :
\begin{equation}\mathcal{R}_{26} \le \exp\left( - \frac{ \left( \Delta_{26, 1} \right)^2 }{ 2 \sum_{j=2}^{26} \Delta_{26, j} + \frac{1}{3} \max_{j} \Delta_{26, j} } \right)\end{equation}Le contrôle de ce terme $\mathcal{R}_{26}$ certifie l'indépendance quasi-totale des choix pour les structures de dimension 26.

\subsubsection{Analyse du moment d'ordre 27}Pour le degré 27, nous évaluons l'intersection simultanée de $\ell = 27$ hyperarêtes. Soient $E_1, \dots, E_{27}$ des arêtes de l'hypergraphe.
\begin{equation}\mathbb{E}\left[ \prod_{i=1}^{27} \mathbf{1}_{E_i \cap X \neq \emptyset} \right] = \sum_{j=1}^{27} (-1)^{j-1} \sum_{1 \le i_1 < \dots < i_j \le 27} \mathbb{P}\left( \bigcap_{m=1}^j (E_{i_m} \cap X \neq \emptyset) \right)\end{equation}Nous décomposons la probabilité conjointe en utilisant l'indépendance conditionnelle sur le choix des sommets de $X$. L'ensemble des configurations de recouvrement de l'union $\bigcup_{i} E_i$ est partitionné selon la taille de leur intersection avec $X$.
\begin{align} \\\Delta_{27, 1} &= \sum_{|S| = 1} \mathbb{P}(S \subseteq X) \sum_{I \subseteq \{1, \dots, 27\}, |I| \ge 1} (-1)^{|I|+1} \mathbf{1}_{S \subseteq \cup_{i \in I} E_i}
\Delta_{27, 2} &= \sum_{|S| = 2} \mathbb{P}(S \subseteq X) \sum_{I \subseteq \{1, \dots, 27\}, |I| \ge 1} (-1)^{|I|+1} \mathbf{1}_{S \subseteq \cup_{i \in I} E_i}
\Delta_{27, 3} &= \sum_{|S| = 3} \mathbb{P}(S \subseteq X) \sum_{I \subseteq \{1, \dots, 27\}, |I| \ge 1} (-1)^{|I|+1} \mathbf{1}_{S \subseteq \cup_{i \in I} E_i}
\Delta_{27, 4} &= \sum_{|S| = 4} \mathbb{P}(S \subseteq X) \sum_{I \subseteq \{1, \dots, 27\}, |I| \ge 1} (-1)^{|I|+1} \mathbf{1}_{S \subseteq \cup_{i \in I} E_i}
\Delta_{27, 5} &= \sum_{|S| = 5} \mathbb{P}(S \subseteq X) \sum_{I \subseteq \{1, \dots, 27\}, |I| \ge 1} (-1)^{|I|+1} \mathbf{1}_{S \subseteq \cup_{i \in I} E_i}
\Delta_{27, 6} &= \sum_{|S| = 6} \mathbb{P}(S \subseteq X) \sum_{I \subseteq \{1, \dots, 27\}, |I| \ge 1} (-1)^{|I|+1} \mathbf{1}_{S \subseteq \cup_{i \in I} E_i}
\Delta_{27, 7} &= \sum_{|S| = 7} \mathbb{P}(S \subseteq X) \sum_{I \subseteq \{1, \dots, 27\}, |I| \ge 1} (-1)^{|I|+1} \mathbf{1}_{S \subseteq \cup_{i \in I} E_i}
\Delta_{27, 8} &= \sum_{|S| = 8} \mathbb{P}(S \subseteq X) \sum_{I \subseteq \{1, \dots, 27\}, |I| \ge 1} (-1)^{|I|+1} \mathbf{1}_{S \subseteq \cup_{i \in I} E_i}
\Delta_{27, 9} &= \sum_{|S| = 9} \mathbb{P}(S \subseteq X) \sum_{I \subseteq \{1, \dots, 27\}, |I| \ge 1} (-1)^{|I|+1} \mathbf{1}_{S \subseteq \cup_{i \in I} E_i}\end{align}Cette majoration permet de borner supérieurement la corrélation locale. Par application du lemme de Lovász local ou des inégalités de concentration azumaiennes, la déviation par rapport à la moyenne produit un terme résiduel qui s'amortit exponentiellement. Nous fixons explicitement cette décroissance :
\begin{equation}\mathcal{R}_{27} \le \exp\left( - \frac{ \left( \Delta_{27, 1} \right)^2 }{ 2 \sum_{j=2}^{27} \Delta_{27, j} + \frac{1}{3} \max_{j} \Delta_{27, j} } \right)\end{equation}Le contrôle de ce terme $\mathcal{R}_{27}$ certifie l'indépendance quasi-totale des choix pour les structures de dimension 27.

\subsubsection{Analyse du moment d'ordre 28}Pour le degré 28, nous évaluons l'intersection simultanée de $\ell = 28$ hyperarêtes. Soient $E_1, \dots, E_{28}$ des arêtes de l'hypergraphe.
\begin{equation}\mathbb{E}\left[ \prod_{i=1}^{28} \mathbf{1}_{E_i \cap X \neq \emptyset} \right] = \sum_{j=1}^{28} (-1)^{j-1} \sum_{1 \le i_1 < \dots < i_j \le 28} \mathbb{P}\left( \bigcap_{m=1}^j (E_{i_m} \cap X \neq \emptyset) \right)\end{equation}Nous décomposons la probabilité conjointe en utilisant l'indépendance conditionnelle sur le choix des sommets de $X$. L'ensemble des configurations de recouvrement de l'union $\bigcup_{i} E_i$ est partitionné selon la taille de leur intersection avec $X$.
\begin{align} \\\Delta_{28, 1} &= \sum_{|S| = 1} \mathbb{P}(S \subseteq X) \sum_{I \subseteq \{1, \dots, 28\}, |I| \ge 1} (-1)^{|I|+1} \mathbf{1}_{S \subseteq \cup_{i \in I} E_i}
\Delta_{28, 2} &= \sum_{|S| = 2} \mathbb{P}(S \subseteq X) \sum_{I \subseteq \{1, \dots, 28\}, |I| \ge 1} (-1)^{|I|+1} \mathbf{1}_{S \subseteq \cup_{i \in I} E_i}
\Delta_{28, 3} &= \sum_{|S| = 3} \mathbb{P}(S \subseteq X) \sum_{I \subseteq \{1, \dots, 28\}, |I| \ge 1} (-1)^{|I|+1} \mathbf{1}_{S \subseteq \cup_{i \in I} E_i}
\Delta_{28, 4} &= \sum_{|S| = 4} \mathbb{P}(S \subseteq X) \sum_{I \subseteq \{1, \dots, 28\}, |I| \ge 1} (-1)^{|I|+1} \mathbf{1}_{S \subseteq \cup_{i \in I} E_i}
\Delta_{28, 5} &= \sum_{|S| = 5} \mathbb{P}(S \subseteq X) \sum_{I \subseteq \{1, \dots, 28\}, |I| \ge 1} (-1)^{|I|+1} \mathbf{1}_{S \subseteq \cup_{i \in I} E_i}
\Delta_{28, 6} &= \sum_{|S| = 6} \mathbb{P}(S \subseteq X) \sum_{I \subseteq \{1, \dots, 28\}, |I| \ge 1} (-1)^{|I|+1} \mathbf{1}_{S \subseteq \cup_{i \in I} E_i}
\Delta_{28, 7} &= \sum_{|S| = 7} \mathbb{P}(S \subseteq X) \sum_{I \subseteq \{1, \dots, 28\}, |I| \ge 1} (-1)^{|I|+1} \mathbf{1}_{S \subseteq \cup_{i \in I} E_i}
\Delta_{28, 8} &= \sum_{|S| = 8} \mathbb{P}(S \subseteq X) \sum_{I \subseteq \{1, \dots, 28\}, |I| \ge 1} (-1)^{|I|+1} \mathbf{1}_{S \subseteq \cup_{i \in I} E_i}
\Delta_{28, 9} &= \sum_{|S| = 9} \mathbb{P}(S \subseteq X) \sum_{I \subseteq \{1, \dots, 28\}, |I| \ge 1} (-1)^{|I|+1} \mathbf{1}_{S \subseteq \cup_{i \in I} E_i}\end{align}Cette majoration permet de borner supérieurement la corrélation locale. Par application du lemme de Lovász local ou des inégalités de concentration azumaiennes, la déviation par rapport à la moyenne produit un terme résiduel qui s'amortit exponentiellement. Nous fixons explicitement cette décroissance :
\begin{equation}\mathcal{R}_{28} \le \exp\left( - \frac{ \left( \Delta_{28, 1} \right)^2 }{ 2 \sum_{j=2}^{28} \Delta_{28, j} + \frac{1}{3} \max_{j} \Delta_{28, j} } \right)\end{equation}Le contrôle de ce terme $\mathcal{R}_{28}$ certifie l'indépendance quasi-totale des choix pour les structures de dimension 28.

\subsubsection{Analyse du moment d'ordre 29}Pour le degré 29, nous évaluons l'intersection simultanée de $\ell = 29$ hyperarêtes. Soient $E_1, \dots, E_{29}$ des arêtes de l'hypergraphe.
\begin{equation}\mathbb{E}\left[ \prod_{i=1}^{29} \mathbf{1}_{E_i \cap X \neq \emptyset} \right] = \sum_{j=1}^{29} (-1)^{j-1} \sum_{1 \le i_1 < \dots < i_j \le 29} \mathbb{P}\left( \bigcap_{m=1}^j (E_{i_m} \cap X \neq \emptyset) \right)\end{equation}Nous décomposons la probabilité conjointe en utilisant l'indépendance conditionnelle sur le choix des sommets de $X$. L'ensemble des configurations de recouvrement de l'union $\bigcup_{i} E_i$ est partitionné selon la taille de leur intersection avec $X$.
\begin{align} \\\Delta_{29, 1} &= \sum_{|S| = 1} \mathbb{P}(S \subseteq X) \sum_{I \subseteq \{1, \dots, 29\}, |I| \ge 1} (-1)^{|I|+1} \mathbf{1}_{S \subseteq \cup_{i \in I} E_i}
\Delta_{29, 2} &= \sum_{|S| = 2} \mathbb{P}(S \subseteq X) \sum_{I \subseteq \{1, \dots, 29\}, |I| \ge 1} (-1)^{|I|+1} \mathbf{1}_{S \subseteq \cup_{i \in I} E_i}
\Delta_{29, 3} &= \sum_{|S| = 3} \mathbb{P}(S \subseteq X) \sum_{I \subseteq \{1, \dots, 29\}, |I| \ge 1} (-1)^{|I|+1} \mathbf{1}_{S \subseteq \cup_{i \in I} E_i}
\Delta_{29, 4} &= \sum_{|S| = 4} \mathbb{P}(S \subseteq X) \sum_{I \subseteq \{1, \dots, 29\}, |I| \ge 1} (-1)^{|I|+1} \mathbf{1}_{S \subseteq \cup_{i \in I} E_i}
\Delta_{29, 5} &= \sum_{|S| = 5} \mathbb{P}(S \subseteq X) \sum_{I \subseteq \{1, \dots, 29\}, |I| \ge 1} (-1)^{|I|+1} \mathbf{1}_{S \subseteq \cup_{i \in I} E_i}
\Delta_{29, 6} &= \sum_{|S| = 6} \mathbb{P}(S \subseteq X) \sum_{I \subseteq \{1, \dots, 29\}, |I| \ge 1} (-1)^{|I|+1} \mathbf{1}_{S \subseteq \cup_{i \in I} E_i}
\Delta_{29, 7} &= \sum_{|S| = 7} \mathbb{P}(S \subseteq X) \sum_{I \subseteq \{1, \dots, 29\}, |I| \ge 1} (-1)^{|I|+1} \mathbf{1}_{S \subseteq \cup_{i \in I} E_i}
\Delta_{29, 8} &= \sum_{|S| = 8} \mathbb{P}(S \subseteq X) \sum_{I \subseteq \{1, \dots, 29\}, |I| \ge 1} (-1)^{|I|+1} \mathbf{1}_{S \subseteq \cup_{i \in I} E_i}
\Delta_{29, 9} &= \sum_{|S| = 9} \mathbb{P}(S \subseteq X) \sum_{I \subseteq \{1, \dots, 29\}, |I| \ge 1} (-1)^{|I|+1} \mathbf{1}_{S \subseteq \cup_{i \in I} E_i}\end{align}Cette majoration permet de borner supérieurement la corrélation locale. Par application du lemme de Lovász local ou des inégalités de concentration azumaiennes, la déviation par rapport à la moyenne produit un terme résiduel qui s'amortit exponentiellement. Nous fixons explicitement cette décroissance :
\begin{equation}\mathcal{R}_{29} \le \exp\left( - \frac{ \left( \Delta_{29, 1} \right)^2 }{ 2 \sum_{j=2}^{29} \Delta_{29, j} + \frac{1}{3} \max_{j} \Delta_{29, j} } \right)\end{equation}Le contrôle de ce terme $\mathcal{R}_{29}$ certifie l'indépendance quasi-totale des choix pour les structures de dimension 29.

\subsubsection{Analyse du moment d'ordre 30}Pour le degré 30, nous évaluons l'intersection simultanée de $\ell = 30$ hyperarêtes. Soient $E_1, \dots, E_{30}$ des arêtes de l'hypergraphe.
\begin{equation}\mathbb{E}\left[ \prod_{i=1}^{30} \mathbf{1}_{E_i \cap X \neq \emptyset} \right] = \sum_{j=1}^{30} (-1)^{j-1} \sum_{1 \le i_1 < \dots < i_j \le 30} \mathbb{P}\left( \bigcap_{m=1}^j (E_{i_m} \cap X \neq \emptyset) \right)\end{equation}Nous décomposons la probabilité conjointe en utilisant l'indépendance conditionnelle sur le choix des sommets de $X$. L'ensemble des configurations de recouvrement de l'union $\bigcup_{i} E_i$ est partitionné selon la taille de leur intersection avec $X$.
\begin{align} \\\Delta_{30, 1} &= \sum_{|S| = 1} \mathbb{P}(S \subseteq X) \sum_{I \subseteq \{1, \dots, 30\}, |I| \ge 1} (-1)^{|I|+1} \mathbf{1}_{S \subseteq \cup_{i \in I} E_i}
\Delta_{30, 2} &= \sum_{|S| = 2} \mathbb{P}(S \subseteq X) \sum_{I \subseteq \{1, \dots, 30\}, |I| \ge 1} (-1)^{|I|+1} \mathbf{1}_{S \subseteq \cup_{i \in I} E_i}
\Delta_{30, 3} &= \sum_{|S| = 3} \mathbb{P}(S \subseteq X) \sum_{I \subseteq \{1, \dots, 30\}, |I| \ge 1} (-1)^{|I|+1} \mathbf{1}_{S \subseteq \cup_{i \in I} E_i}
\Delta_{30, 4} &= \sum_{|S| = 4} \mathbb{P}(S \subseteq X) \sum_{I \subseteq \{1, \dots, 30\}, |I| \ge 1} (-1)^{|I|+1} \mathbf{1}_{S \subseteq \cup_{i \in I} E_i}
\Delta_{30, 5} &= \sum_{|S| = 5} \mathbb{P}(S \subseteq X) \sum_{I \subseteq \{1, \dots, 30\}, |I| \ge 1} (-1)^{|I|+1} \mathbf{1}_{S \subseteq \cup_{i \in I} E_i}
\Delta_{30, 6} &= \sum_{|S| = 6} \mathbb{P}(S \subseteq X) \sum_{I \subseteq \{1, \dots, 30\}, |I| \ge 1} (-1)^{|I|+1} \mathbf{1}_{S \subseteq \cup_{i \in I} E_i}
\Delta_{30, 7} &= \sum_{|S| = 7} \mathbb{P}(S \subseteq X) \sum_{I \subseteq \{1, \dots, 30\}, |I| \ge 1} (-1)^{|I|+1} \mathbf{1}_{S \subseteq \cup_{i \in I} E_i}
\Delta_{30, 8} &= \sum_{|S| = 8} \mathbb{P}(S \subseteq X) \sum_{I \subseteq \{1, \dots, 30\}, |I| \ge 1} (-1)^{|I|+1} \mathbf{1}_{S \subseteq \cup_{i \in I} E_i}
\Delta_{30, 9} &= \sum_{|S| = 9} \mathbb{P}(S \subseteq X) \sum_{I \subseteq \{1, \dots, 30\}, |I| \ge 1} (-1)^{|I|+1} \mathbf{1}_{S \subseteq \cup_{i \in I} E_i}\end{align}Cette majoration permet de borner supérieurement la corrélation locale. Par application du lemme de Lovász local ou des inégalités de concentration azumaiennes, la déviation par rapport à la moyenne produit un terme résiduel qui s'amortit exponentiellement. Nous fixons explicitement cette décroissance :
\begin{equation}\mathcal{R}_{30} \le \exp\left( - \frac{ \left( \Delta_{30, 1} \right)^2 }{ 2 \sum_{j=2}^{30} \Delta_{30, j} + \frac{1}{3} \max_{j} \Delta_{30, j} } \right)\end{equation}Le contrôle de ce terme $\mathcal{R}_{30}$ certifie l'indépendance quasi-totale des choix pour les structures de dimension 30.

\subsubsection{Analyse du moment d'ordre 31}Pour le degré 31, nous évaluons l'intersection simultanée de $\ell = 31$ hyperarêtes. Soient $E_1, \dots, E_{31}$ des arêtes de l'hypergraphe.
\begin{equation}\mathbb{E}\left[ \prod_{i=1}^{31} \mathbf{1}_{E_i \cap X \neq \emptyset} \right] = \sum_{j=1}^{31} (-1)^{j-1} \sum_{1 \le i_1 < \dots < i_j \le 31} \mathbb{P}\left( \bigcap_{m=1}^j (E_{i_m} \cap X \neq \emptyset) \right)\end{equation}Nous décomposons la probabilité conjointe en utilisant l'indépendance conditionnelle sur le choix des sommets de $X$. L'ensemble des configurations de recouvrement de l'union $\bigcup_{i} E_i$ est partitionné selon la taille de leur intersection avec $X$.
\begin{align} \\\Delta_{31, 1} &= \sum_{|S| = 1} \mathbb{P}(S \subseteq X) \sum_{I \subseteq \{1, \dots, 31\}, |I| \ge 1} (-1)^{|I|+1} \mathbf{1}_{S \subseteq \cup_{i \in I} E_i}
\Delta_{31, 2} &= \sum_{|S| = 2} \mathbb{P}(S \subseteq X) \sum_{I \subseteq \{1, \dots, 31\}, |I| \ge 1} (-1)^{|I|+1} \mathbf{1}_{S \subseteq \cup_{i \in I} E_i}
\Delta_{31, 3} &= \sum_{|S| = 3} \mathbb{P}(S \subseteq X) \sum_{I \subseteq \{1, \dots, 31\}, |I| \ge 1} (-1)^{|I|+1} \mathbf{1}_{S \subseteq \cup_{i \in I} E_i}
\Delta_{31, 4} &= \sum_{|S| = 4} \mathbb{P}(S \subseteq X) \sum_{I \subseteq \{1, \dots, 31\}, |I| \ge 1} (-1)^{|I|+1} \mathbf{1}_{S \subseteq \cup_{i \in I} E_i}
\Delta_{31, 5} &= \sum_{|S| = 5} \mathbb{P}(S \subseteq X) \sum_{I \subseteq \{1, \dots, 31\}, |I| \ge 1} (-1)^{|I|+1} \mathbf{1}_{S \subseteq \cup_{i \in I} E_i}
\Delta_{31, 6} &= \sum_{|S| = 6} \mathbb{P}(S \subseteq X) \sum_{I \subseteq \{1, \dots, 31\}, |I| \ge 1} (-1)^{|I|+1} \mathbf{1}_{S \subseteq \cup_{i \in I} E_i}
\Delta_{31, 7} &= \sum_{|S| = 7} \mathbb{P}(S \subseteq X) \sum_{I \subseteq \{1, \dots, 31\}, |I| \ge 1} (-1)^{|I|+1} \mathbf{1}_{S \subseteq \cup_{i \in I} E_i}
\Delta_{31, 8} &= \sum_{|S| = 8} \mathbb{P}(S \subseteq X) \sum_{I \subseteq \{1, \dots, 31\}, |I| \ge 1} (-1)^{|I|+1} \mathbf{1}_{S \subseteq \cup_{i \in I} E_i}
\Delta_{31, 9} &= \sum_{|S| = 9} \mathbb{P}(S \subseteq X) \sum_{I \subseteq \{1, \dots, 31\}, |I| \ge 1} (-1)^{|I|+1} \mathbf{1}_{S \subseteq \cup_{i \in I} E_i}\end{align}Cette majoration permet de borner supérieurement la corrélation locale. Par application du lemme de Lovász local ou des inégalités de concentration azumaiennes, la déviation par rapport à la moyenne produit un terme résiduel qui s'amortit exponentiellement. Nous fixons explicitement cette décroissance :
\begin{equation}\mathcal{R}_{31} \le \exp\left( - \frac{ \left( \Delta_{31, 1} \right)^2 }{ 2 \sum_{j=2}^{31} \Delta_{31, j} + \frac{1}{3} \max_{j} \Delta_{31, j} } \right)\end{equation}Le contrôle de ce terme $\mathcal{R}_{31}$ certifie l'indépendance quasi-totale des choix pour les structures de dimension 31.

\subsubsection{Analyse du moment d'ordre 32}Pour le degré 32, nous évaluons l'intersection simultanée de $\ell = 32$ hyperarêtes. Soient $E_1, \dots, E_{32}$ des arêtes de l'hypergraphe.
\begin{equation}\mathbb{E}\left[ \prod_{i=1}^{32} \mathbf{1}_{E_i \cap X \neq \emptyset} \right] = \sum_{j=1}^{32} (-1)^{j-1} \sum_{1 \le i_1 < \dots < i_j \le 32} \mathbb{P}\left( \bigcap_{m=1}^j (E_{i_m} \cap X \neq \emptyset) \right)\end{equation}Nous décomposons la probabilité conjointe en utilisant l'indépendance conditionnelle sur le choix des sommets de $X$. L'ensemble des configurations de recouvrement de l'union $\bigcup_{i} E_i$ est partitionné selon la taille de leur intersection avec $X$.
\begin{align} \\\Delta_{32, 1} &= \sum_{|S| = 1} \mathbb{P}(S \subseteq X) \sum_{I \subseteq \{1, \dots, 32\}, |I| \ge 1} (-1)^{|I|+1} \mathbf{1}_{S \subseteq \cup_{i \in I} E_i}
\Delta_{32, 2} &= \sum_{|S| = 2} \mathbb{P}(S \subseteq X) \sum_{I \subseteq \{1, \dots, 32\}, |I| \ge 1} (-1)^{|I|+1} \mathbf{1}_{S \subseteq \cup_{i \in I} E_i}
\Delta_{32, 3} &= \sum_{|S| = 3} \mathbb{P}(S \subseteq X) \sum_{I \subseteq \{1, \dots, 32\}, |I| \ge 1} (-1)^{|I|+1} \mathbf{1}_{S \subseteq \cup_{i \in I} E_i}
\Delta_{32, 4} &= \sum_{|S| = 4} \mathbb{P}(S \subseteq X) \sum_{I \subseteq \{1, \dots, 32\}, |I| \ge 1} (-1)^{|I|+1} \mathbf{1}_{S \subseteq \cup_{i \in I} E_i}
\Delta_{32, 5} &= \sum_{|S| = 5} \mathbb{P}(S \subseteq X) \sum_{I \subseteq \{1, \dots, 32\}, |I| \ge 1} (-1)^{|I|+1} \mathbf{1}_{S \subseteq \cup_{i \in I} E_i}
\Delta_{32, 6} &= \sum_{|S| = 6} \mathbb{P}(S \subseteq X) \sum_{I \subseteq \{1, \dots, 32\}, |I| \ge 1} (-1)^{|I|+1} \mathbf{1}_{S \subseteq \cup_{i \in I} E_i}
\Delta_{32, 7} &= \sum_{|S| = 7} \mathbb{P}(S \subseteq X) \sum_{I \subseteq \{1, \dots, 32\}, |I| \ge 1} (-1)^{|I|+1} \mathbf{1}_{S \subseteq \cup_{i \in I} E_i}
\Delta_{32, 8} &= \sum_{|S| = 8} \mathbb{P}(S \subseteq X) \sum_{I \subseteq \{1, \dots, 32\}, |I| \ge 1} (-1)^{|I|+1} \mathbf{1}_{S \subseteq \cup_{i \in I} E_i}
\Delta_{32, 9} &= \sum_{|S| = 9} \mathbb{P}(S \subseteq X) \sum_{I \subseteq \{1, \dots, 32\}, |I| \ge 1} (-1)^{|I|+1} \mathbf{1}_{S \subseteq \cup_{i \in I} E_i}\end{align}Cette majoration permet de borner supérieurement la corrélation locale. Par application du lemme de Lovász local ou des inégalités de concentration azumaiennes, la déviation par rapport à la moyenne produit un terme résiduel qui s'amortit exponentiellement. Nous fixons explicitement cette décroissance :
\begin{equation}\mathcal{R}_{32} \le \exp\left( - \frac{ \left( \Delta_{32, 1} \right)^2 }{ 2 \sum_{j=2}^{32} \Delta_{32, j} + \frac{1}{3} \max_{j} \Delta_{32, j} } \right)\end{equation}Le contrôle de ce terme $\mathcal{R}_{32}$ certifie l'indépendance quasi-totale des choix pour les structures de dimension 32.

\subsubsection{Analyse du moment d'ordre 33}Pour le degré 33, nous évaluons l'intersection simultanée de $\ell = 33$ hyperarêtes. Soient $E_1, \dots, E_{33}$ des arêtes de l'hypergraphe.
\begin{equation}\mathbb{E}\left[ \prod_{i=1}^{33} \mathbf{1}_{E_i \cap X \neq \emptyset} \right] = \sum_{j=1}^{33} (-1)^{j-1} \sum_{1 \le i_1 < \dots < i_j \le 33} \mathbb{P}\left( \bigcap_{m=1}^j (E_{i_m} \cap X \neq \emptyset) \right)\end{equation}Nous décomposons la probabilité conjointe en utilisant l'indépendance conditionnelle sur le choix des sommets de $X$. L'ensemble des configurations de recouvrement de l'union $\bigcup_{i} E_i$ est partitionné selon la taille de leur intersection avec $X$.
\begin{align} \\\Delta_{33, 1} &= \sum_{|S| = 1} \mathbb{P}(S \subseteq X) \sum_{I \subseteq \{1, \dots, 33\}, |I| \ge 1} (-1)^{|I|+1} \mathbf{1}_{S \subseteq \cup_{i \in I} E_i}
\Delta_{33, 2} &= \sum_{|S| = 2} \mathbb{P}(S \subseteq X) \sum_{I \subseteq \{1, \dots, 33\}, |I| \ge 1} (-1)^{|I|+1} \mathbf{1}_{S \subseteq \cup_{i \in I} E_i}
\Delta_{33, 3} &= \sum_{|S| = 3} \mathbb{P}(S \subseteq X) \sum_{I \subseteq \{1, \dots, 33\}, |I| \ge 1} (-1)^{|I|+1} \mathbf{1}_{S \subseteq \cup_{i \in I} E_i}
\Delta_{33, 4} &= \sum_{|S| = 4} \mathbb{P}(S \subseteq X) \sum_{I \subseteq \{1, \dots, 33\}, |I| \ge 1} (-1)^{|I|+1} \mathbf{1}_{S \subseteq \cup_{i \in I} E_i}
\Delta_{33, 5} &= \sum_{|S| = 5} \mathbb{P}(S \subseteq X) \sum_{I \subseteq \{1, \dots, 33\}, |I| \ge 1} (-1)^{|I|+1} \mathbf{1}_{S \subseteq \cup_{i \in I} E_i}
\Delta_{33, 6} &= \sum_{|S| = 6} \mathbb{P}(S \subseteq X) \sum_{I \subseteq \{1, \dots, 33\}, |I| \ge 1} (-1)^{|I|+1} \mathbf{1}_{S \subseteq \cup_{i \in I} E_i}
\Delta_{33, 7} &= \sum_{|S| = 7} \mathbb{P}(S \subseteq X) \sum_{I \subseteq \{1, \dots, 33\}, |I| \ge 1} (-1)^{|I|+1} \mathbf{1}_{S \subseteq \cup_{i \in I} E_i}
\Delta_{33, 8} &= \sum_{|S| = 8} \mathbb{P}(S \subseteq X) \sum_{I \subseteq \{1, \dots, 33\}, |I| \ge 1} (-1)^{|I|+1} \mathbf{1}_{S \subseteq \cup_{i \in I} E_i}
\Delta_{33, 9} &= \sum_{|S| = 9} \mathbb{P}(S \subseteq X) \sum_{I \subseteq \{1, \dots, 33\}, |I| \ge 1} (-1)^{|I|+1} \mathbf{1}_{S \subseteq \cup_{i \in I} E_i}\end{align}Cette majoration permet de borner supérieurement la corrélation locale. Par application du lemme de Lovász local ou des inégalités de concentration azumaiennes, la déviation par rapport à la moyenne produit un terme résiduel qui s'amortit exponentiellement. Nous fixons explicitement cette décroissance :
\begin{equation}\mathcal{R}_{33} \le \exp\left( - \frac{ \left( \Delta_{33, 1} \right)^2 }{ 2 \sum_{j=2}^{33} \Delta_{33, j} + \frac{1}{3} \max_{j} \Delta_{33, j} } \right)\end{equation}Le contrôle de ce terme $\mathcal{R}_{33}$ certifie l'indépendance quasi-totale des choix pour les structures de dimension 33.

\subsubsection{Analyse du moment d'ordre 34}Pour le degré 34, nous évaluons l'intersection simultanée de $\ell = 34$ hyperarêtes. Soient $E_1, \dots, E_{34}$ des arêtes de l'hypergraphe.
\begin{equation}\mathbb{E}\left[ \prod_{i=1}^{34} \mathbf{1}_{E_i \cap X \neq \emptyset} \right] = \sum_{j=1}^{34} (-1)^{j-1} \sum_{1 \le i_1 < \dots < i_j \le 34} \mathbb{P}\left( \bigcap_{m=1}^j (E_{i_m} \cap X \neq \emptyset) \right)\end{equation}Nous décomposons la probabilité conjointe en utilisant l'indépendance conditionnelle sur le choix des sommets de $X$. L'ensemble des configurations de recouvrement de l'union $\bigcup_{i} E_i$ est partitionné selon la taille de leur intersection avec $X$.
\begin{align} \\\Delta_{34, 1} &= \sum_{|S| = 1} \mathbb{P}(S \subseteq X) \sum_{I \subseteq \{1, \dots, 34\}, |I| \ge 1} (-1)^{|I|+1} \mathbf{1}_{S \subseteq \cup_{i \in I} E_i}
\Delta_{34, 2} &= \sum_{|S| = 2} \mathbb{P}(S \subseteq X) \sum_{I \subseteq \{1, \dots, 34\}, |I| \ge 1} (-1)^{|I|+1} \mathbf{1}_{S \subseteq \cup_{i \in I} E_i}
\Delta_{34, 3} &= \sum_{|S| = 3} \mathbb{P}(S \subseteq X) \sum_{I \subseteq \{1, \dots, 34\}, |I| \ge 1} (-1)^{|I|+1} \mathbf{1}_{S \subseteq \cup_{i \in I} E_i}
\Delta_{34, 4} &= \sum_{|S| = 4} \mathbb{P}(S \subseteq X) \sum_{I \subseteq \{1, \dots, 34\}, |I| \ge 1} (-1)^{|I|+1} \mathbf{1}_{S \subseteq \cup_{i \in I} E_i}
\Delta_{34, 5} &= \sum_{|S| = 5} \mathbb{P}(S \subseteq X) \sum_{I \subseteq \{1, \dots, 34\}, |I| \ge 1} (-1)^{|I|+1} \mathbf{1}_{S \subseteq \cup_{i \in I} E_i}
\Delta_{34, 6} &= \sum_{|S| = 6} \mathbb{P}(S \subseteq X) \sum_{I \subseteq \{1, \dots, 34\}, |I| \ge 1} (-1)^{|I|+1} \mathbf{1}_{S \subseteq \cup_{i \in I} E_i}
\Delta_{34, 7} &= \sum_{|S| = 7} \mathbb{P}(S \subseteq X) \sum_{I \subseteq \{1, \dots, 34\}, |I| \ge 1} (-1)^{|I|+1} \mathbf{1}_{S \subseteq \cup_{i \in I} E_i}
\Delta_{34, 8} &= \sum_{|S| = 8} \mathbb{P}(S \subseteq X) \sum_{I \subseteq \{1, \dots, 34\}, |I| \ge 1} (-1)^{|I|+1} \mathbf{1}_{S \subseteq \cup_{i \in I} E_i}
\Delta_{34, 9} &= \sum_{|S| = 9} \mathbb{P}(S \subseteq X) \sum_{I \subseteq \{1, \dots, 34\}, |I| \ge 1} (-1)^{|I|+1} \mathbf{1}_{S \subseteq \cup_{i \in I} E_i}\end{align}Cette majoration permet de borner supérieurement la corrélation locale. Par application du lemme de Lovász local ou des inégalités de concentration azumaiennes, la déviation par rapport à la moyenne produit un terme résiduel qui s'amortit exponentiellement. Nous fixons explicitement cette décroissance :
\begin{equation}\mathcal{R}_{34} \le \exp\left( - \frac{ \left( \Delta_{34, 1} \right)^2 }{ 2 \sum_{j=2}^{34} \Delta_{34, j} + \frac{1}{3} \max_{j} \Delta_{34, j} } \right)\end{equation}Le contrôle de ce terme $\mathcal{R}_{34}$ certifie l'indépendance quasi-totale des choix pour les structures de dimension 34.

\subsubsection{Analyse du moment d'ordre 35}Pour le degré 35, nous évaluons l'intersection simultanée de $\ell = 35$ hyperarêtes. Soient $E_1, \dots, E_{35}$ des arêtes de l'hypergraphe.
\begin{equation}\mathbb{E}\left[ \prod_{i=1}^{35} \mathbf{1}_{E_i \cap X \neq \emptyset} \right] = \sum_{j=1}^{35} (-1)^{j-1} \sum_{1 \le i_1 < \dots < i_j \le 35} \mathbb{P}\left( \bigcap_{m=1}^j (E_{i_m} \cap X \neq \emptyset) \right)\end{equation}Nous décomposons la probabilité conjointe en utilisant l'indépendance conditionnelle sur le choix des sommets de $X$. L'ensemble des configurations de recouvrement de l'union $\bigcup_{i} E_i$ est partitionné selon la taille de leur intersection avec $X$.
\begin{align} \\\Delta_{35, 1} &= \sum_{|S| = 1} \mathbb{P}(S \subseteq X) \sum_{I \subseteq \{1, \dots, 35\}, |I| \ge 1} (-1)^{|I|+1} \mathbf{1}_{S \subseteq \cup_{i \in I} E_i}
\Delta_{35, 2} &= \sum_{|S| = 2} \mathbb{P}(S \subseteq X) \sum_{I \subseteq \{1, \dots, 35\}, |I| \ge 1} (-1)^{|I|+1} \mathbf{1}_{S \subseteq \cup_{i \in I} E_i}
\Delta_{35, 3} &= \sum_{|S| = 3} \mathbb{P}(S \subseteq X) \sum_{I \subseteq \{1, \dots, 35\}, |I| \ge 1} (-1)^{|I|+1} \mathbf{1}_{S \subseteq \cup_{i \in I} E_i}
\Delta_{35, 4} &= \sum_{|S| = 4} \mathbb{P}(S \subseteq X) \sum_{I \subseteq \{1, \dots, 35\}, |I| \ge 1} (-1)^{|I|+1} \mathbf{1}_{S \subseteq \cup_{i \in I} E_i}
\Delta_{35, 5} &= \sum_{|S| = 5} \mathbb{P}(S \subseteq X) \sum_{I \subseteq \{1, \dots, 35\}, |I| \ge 1} (-1)^{|I|+1} \mathbf{1}_{S \subseteq \cup_{i \in I} E_i}
\Delta_{35, 6} &= \sum_{|S| = 6} \mathbb{P}(S \subseteq X) \sum_{I \subseteq \{1, \dots, 35\}, |I| \ge 1} (-1)^{|I|+1} \mathbf{1}_{S \subseteq \cup_{i \in I} E_i}
\Delta_{35, 7} &= \sum_{|S| = 7} \mathbb{P}(S \subseteq X) \sum_{I \subseteq \{1, \dots, 35\}, |I| \ge 1} (-1)^{|I|+1} \mathbf{1}_{S \subseteq \cup_{i \in I} E_i}
\Delta_{35, 8} &= \sum_{|S| = 8} \mathbb{P}(S \subseteq X) \sum_{I \subseteq \{1, \dots, 35\}, |I| \ge 1} (-1)^{|I|+1} \mathbf{1}_{S \subseteq \cup_{i \in I} E_i}
\Delta_{35, 9} &= \sum_{|S| = 9} \mathbb{P}(S \subseteq X) \sum_{I \subseteq \{1, \dots, 35\}, |I| \ge 1} (-1)^{|I|+1} \mathbf{1}_{S \subseteq \cup_{i \in I} E_i}\end{align}Cette majoration permet de borner supérieurement la corrélation locale. Par application du lemme de Lovász local ou des inégalités de concentration azumaiennes, la déviation par rapport à la moyenne produit un terme résiduel qui s'amortit exponentiellement. Nous fixons explicitement cette décroissance :
\begin{equation}\mathcal{R}_{35} \le \exp\left( - \frac{ \left( \Delta_{35, 1} \right)^2 }{ 2 \sum_{j=2}^{35} \Delta_{35, j} + \frac{1}{3} \max_{j} \Delta_{35, j} } \right)\end{equation}Le contrôle de ce terme $\mathcal{R}_{35}$ certifie l'indépendance quasi-totale des choix pour les structures de dimension 35.

\subsubsection{Analyse du moment d'ordre 36}Pour le degré 36, nous évaluons l'intersection simultanée de $\ell = 36$ hyperarêtes. Soient $E_1, \dots, E_{36}$ des arêtes de l'hypergraphe.
\begin{equation}\mathbb{E}\left[ \prod_{i=1}^{36} \mathbf{1}_{E_i \cap X \neq \emptyset} \right] = \sum_{j=1}^{36} (-1)^{j-1} \sum_{1 \le i_1 < \dots < i_j \le 36} \mathbb{P}\left( \bigcap_{m=1}^j (E_{i_m} \cap X \neq \emptyset) \right)\end{equation}Nous décomposons la probabilité conjointe en utilisant l'indépendance conditionnelle sur le choix des sommets de $X$. L'ensemble des configurations de recouvrement de l'union $\bigcup_{i} E_i$ est partitionné selon la taille de leur intersection avec $X$.
\begin{align} \\\Delta_{36, 1} &= \sum_{|S| = 1} \mathbb{P}(S \subseteq X) \sum_{I \subseteq \{1, \dots, 36\}, |I| \ge 1} (-1)^{|I|+1} \mathbf{1}_{S \subseteq \cup_{i \in I} E_i}
\Delta_{36, 2} &= \sum_{|S| = 2} \mathbb{P}(S \subseteq X) \sum_{I \subseteq \{1, \dots, 36\}, |I| \ge 1} (-1)^{|I|+1} \mathbf{1}_{S \subseteq \cup_{i \in I} E_i}
\Delta_{36, 3} &= \sum_{|S| = 3} \mathbb{P}(S \subseteq X) \sum_{I \subseteq \{1, \dots, 36\}, |I| \ge 1} (-1)^{|I|+1} \mathbf{1}_{S \subseteq \cup_{i \in I} E_i}
\Delta_{36, 4} &= \sum_{|S| = 4} \mathbb{P}(S \subseteq X) \sum_{I \subseteq \{1, \dots, 36\}, |I| \ge 1} (-1)^{|I|+1} \mathbf{1}_{S \subseteq \cup_{i \in I} E_i}
\Delta_{36, 5} &= \sum_{|S| = 5} \mathbb{P}(S \subseteq X) \sum_{I \subseteq \{1, \dots, 36\}, |I| \ge 1} (-1)^{|I|+1} \mathbf{1}_{S \subseteq \cup_{i \in I} E_i}
\Delta_{36, 6} &= \sum_{|S| = 6} \mathbb{P}(S \subseteq X) \sum_{I \subseteq \{1, \dots, 36\}, |I| \ge 1} (-1)^{|I|+1} \mathbf{1}_{S \subseteq \cup_{i \in I} E_i}
\Delta_{36, 7} &= \sum_{|S| = 7} \mathbb{P}(S \subseteq X) \sum_{I \subseteq \{1, \dots, 36\}, |I| \ge 1} (-1)^{|I|+1} \mathbf{1}_{S \subseteq \cup_{i \in I} E_i}
\Delta_{36, 8} &= \sum_{|S| = 8} \mathbb{P}(S \subseteq X) \sum_{I \subseteq \{1, \dots, 36\}, |I| \ge 1} (-1)^{|I|+1} \mathbf{1}_{S \subseteq \cup_{i \in I} E_i}
\Delta_{36, 9} &= \sum_{|S| = 9} \mathbb{P}(S \subseteq X) \sum_{I \subseteq \{1, \dots, 36\}, |I| \ge 1} (-1)^{|I|+1} \mathbf{1}_{S \subseteq \cup_{i \in I} E_i}\end{align}Cette majoration permet de borner supérieurement la corrélation locale. Par application du lemme de Lovász local ou des inégalités de concentration azumaiennes, la déviation par rapport à la moyenne produit un terme résiduel qui s'amortit exponentiellement. Nous fixons explicitement cette décroissance :
\begin{equation}\mathcal{R}_{36} \le \exp\left( - \frac{ \left( \Delta_{36, 1} \right)^2 }{ 2 \sum_{j=2}^{36} \Delta_{36, j} + \frac{1}{3} \max_{j} \Delta_{36, j} } \right)\end{equation}Le contrôle de ce terme $\mathcal{R}_{36}$ certifie l'indépendance quasi-totale des choix pour les structures de dimension 36.

\subsubsection{Analyse du moment d'ordre 37}Pour le degré 37, nous évaluons l'intersection simultanée de $\ell = 37$ hyperarêtes. Soient $E_1, \dots, E_{37}$ des arêtes de l'hypergraphe.
\begin{equation}\mathbb{E}\left[ \prod_{i=1}^{37} \mathbf{1}_{E_i \cap X \neq \emptyset} \right] = \sum_{j=1}^{37} (-1)^{j-1} \sum_{1 \le i_1 < \dots < i_j \le 37} \mathbb{P}\left( \bigcap_{m=1}^j (E_{i_m} \cap X \neq \emptyset) \right)\end{equation}Nous décomposons la probabilité conjointe en utilisant l'indépendance conditionnelle sur le choix des sommets de $X$. L'ensemble des configurations de recouvrement de l'union $\bigcup_{i} E_i$ est partitionné selon la taille de leur intersection avec $X$.
\begin{align} \\\Delta_{37, 1} &= \sum_{|S| = 1} \mathbb{P}(S \subseteq X) \sum_{I \subseteq \{1, \dots, 37\}, |I| \ge 1} (-1)^{|I|+1} \mathbf{1}_{S \subseteq \cup_{i \in I} E_i}
\Delta_{37, 2} &= \sum_{|S| = 2} \mathbb{P}(S \subseteq X) \sum_{I \subseteq \{1, \dots, 37\}, |I| \ge 1} (-1)^{|I|+1} \mathbf{1}_{S \subseteq \cup_{i \in I} E_i}
\Delta_{37, 3} &= \sum_{|S| = 3} \mathbb{P}(S \subseteq X) \sum_{I \subseteq \{1, \dots, 37\}, |I| \ge 1} (-1)^{|I|+1} \mathbf{1}_{S \subseteq \cup_{i \in I} E_i}
\Delta_{37, 4} &= \sum_{|S| = 4} \mathbb{P}(S \subseteq X) \sum_{I \subseteq \{1, \dots, 37\}, |I| \ge 1} (-1)^{|I|+1} \mathbf{1}_{S \subseteq \cup_{i \in I} E_i}
\Delta_{37, 5} &= \sum_{|S| = 5} \mathbb{P}(S \subseteq X) \sum_{I \subseteq \{1, \dots, 37\}, |I| \ge 1} (-1)^{|I|+1} \mathbf{1}_{S \subseteq \cup_{i \in I} E_i}
\Delta_{37, 6} &= \sum_{|S| = 6} \mathbb{P}(S \subseteq X) \sum_{I \subseteq \{1, \dots, 37\}, |I| \ge 1} (-1)^{|I|+1} \mathbf{1}_{S \subseteq \cup_{i \in I} E_i}
\Delta_{37, 7} &= \sum_{|S| = 7} \mathbb{P}(S \subseteq X) \sum_{I \subseteq \{1, \dots, 37\}, |I| \ge 1} (-1)^{|I|+1} \mathbf{1}_{S \subseteq \cup_{i \in I} E_i}
\Delta_{37, 8} &= \sum_{|S| = 8} \mathbb{P}(S \subseteq X) \sum_{I \subseteq \{1, \dots, 37\}, |I| \ge 1} (-1)^{|I|+1} \mathbf{1}_{S \subseteq \cup_{i \in I} E_i}
\Delta_{37, 9} &= \sum_{|S| = 9} \mathbb{P}(S \subseteq X) \sum_{I \subseteq \{1, \dots, 37\}, |I| \ge 1} (-1)^{|I|+1} \mathbf{1}_{S \subseteq \cup_{i \in I} E_i}\end{align}Cette majoration permet de borner supérieurement la corrélation locale. Par application du lemme de Lovász local ou des inégalités de concentration azumaiennes, la déviation par rapport à la moyenne produit un terme résiduel qui s'amortit exponentiellement. Nous fixons explicitement cette décroissance :
\begin{equation}\mathcal{R}_{37} \le \exp\left( - \frac{ \left( \Delta_{37, 1} \right)^2 }{ 2 \sum_{j=2}^{37} \Delta_{37, j} + \frac{1}{3} \max_{j} \Delta_{37, j} } \right)\end{equation}Le contrôle de ce terme $\mathcal{R}_{37}$ certifie l'indépendance quasi-totale des choix pour les structures de dimension 37.

\subsubsection{Analyse du moment d'ordre 38}Pour le degré 38, nous évaluons l'intersection simultanée de $\ell = 38$ hyperarêtes. Soient $E_1, \dots, E_{38}$ des arêtes de l'hypergraphe.
\begin{equation}\mathbb{E}\left[ \prod_{i=1}^{38} \mathbf{1}_{E_i \cap X \neq \emptyset} \right] = \sum_{j=1}^{38} (-1)^{j-1} \sum_{1 \le i_1 < \dots < i_j \le 38} \mathbb{P}\left( \bigcap_{m=1}^j (E_{i_m} \cap X \neq \emptyset) \right)\end{equation}Nous décomposons la probabilité conjointe en utilisant l'indépendance conditionnelle sur le choix des sommets de $X$. L'ensemble des configurations de recouvrement de l'union $\bigcup_{i} E_i$ est partitionné selon la taille de leur intersection avec $X$.
\begin{align} \\\Delta_{38, 1} &= \sum_{|S| = 1} \mathbb{P}(S \subseteq X) \sum_{I \subseteq \{1, \dots, 38\}, |I| \ge 1} (-1)^{|I|+1} \mathbf{1}_{S \subseteq \cup_{i \in I} E_i}
\Delta_{38, 2} &= \sum_{|S| = 2} \mathbb{P}(S \subseteq X) \sum_{I \subseteq \{1, \dots, 38\}, |I| \ge 1} (-1)^{|I|+1} \mathbf{1}_{S \subseteq \cup_{i \in I} E_i}
\Delta_{38, 3} &= \sum_{|S| = 3} \mathbb{P}(S \subseteq X) \sum_{I \subseteq \{1, \dots, 38\}, |I| \ge 1} (-1)^{|I|+1} \mathbf{1}_{S \subseteq \cup_{i \in I} E_i}
\Delta_{38, 4} &= \sum_{|S| = 4} \mathbb{P}(S \subseteq X) \sum_{I \subseteq \{1, \dots, 38\}, |I| \ge 1} (-1)^{|I|+1} \mathbf{1}_{S \subseteq \cup_{i \in I} E_i}
\Delta_{38, 5} &= \sum_{|S| = 5} \mathbb{P}(S \subseteq X) \sum_{I \subseteq \{1, \dots, 38\}, |I| \ge 1} (-1)^{|I|+1} \mathbf{1}_{S \subseteq \cup_{i \in I} E_i}
\Delta_{38, 6} &= \sum_{|S| = 6} \mathbb{P}(S \subseteq X) \sum_{I \subseteq \{1, \dots, 38\}, |I| \ge 1} (-1)^{|I|+1} \mathbf{1}_{S \subseteq \cup_{i \in I} E_i}
\Delta_{38, 7} &= \sum_{|S| = 7} \mathbb{P}(S \subseteq X) \sum_{I \subseteq \{1, \dots, 38\}, |I| \ge 1} (-1)^{|I|+1} \mathbf{1}_{S \subseteq \cup_{i \in I} E_i}
\Delta_{38, 8} &= \sum_{|S| = 8} \mathbb{P}(S \subseteq X) \sum_{I \subseteq \{1, \dots, 38\}, |I| \ge 1} (-1)^{|I|+1} \mathbf{1}_{S \subseteq \cup_{i \in I} E_i}
\Delta_{38, 9} &= \sum_{|S| = 9} \mathbb{P}(S \subseteq X) \sum_{I \subseteq \{1, \dots, 38\}, |I| \ge 1} (-1)^{|I|+1} \mathbf{1}_{S \subseteq \cup_{i \in I} E_i}\end{align}Cette majoration permet de borner supérieurement la corrélation locale. Par application du lemme de Lovász local ou des inégalités de concentration azumaiennes, la déviation par rapport à la moyenne produit un terme résiduel qui s'amortit exponentiellement. Nous fixons explicitement cette décroissance :
\begin{equation}\mathcal{R}_{38} \le \exp\left( - \frac{ \left( \Delta_{38, 1} \right)^2 }{ 2 \sum_{j=2}^{38} \Delta_{38, j} + \frac{1}{3} \max_{j} \Delta_{38, j} } \right)\end{equation}Le contrôle de ce terme $\mathcal{R}_{38}$ certifie l'indépendance quasi-totale des choix pour les structures de dimension 38.

\subsubsection{Analyse du moment d'ordre 39}Pour le degré 39, nous évaluons l'intersection simultanée de $\ell = 39$ hyperarêtes. Soient $E_1, \dots, E_{39}$ des arêtes de l'hypergraphe.
\begin{equation}\mathbb{E}\left[ \prod_{i=1}^{39} \mathbf{1}_{E_i \cap X \neq \emptyset} \right] = \sum_{j=1}^{39} (-1)^{j-1} \sum_{1 \le i_1 < \dots < i_j \le 39} \mathbb{P}\left( \bigcap_{m=1}^j (E_{i_m} \cap X \neq \emptyset) \right)\end{equation}Nous décomposons la probabilité conjointe en utilisant l'indépendance conditionnelle sur le choix des sommets de $X$. L'ensemble des configurations de recouvrement de l'union $\bigcup_{i} E_i$ est partitionné selon la taille de leur intersection avec $X$.
\begin{align} \\\Delta_{39, 1} &= \sum_{|S| = 1} \mathbb{P}(S \subseteq X) \sum_{I \subseteq \{1, \dots, 39\}, |I| \ge 1} (-1)^{|I|+1} \mathbf{1}_{S \subseteq \cup_{i \in I} E_i}
\Delta_{39, 2} &= \sum_{|S| = 2} \mathbb{P}(S \subseteq X) \sum_{I \subseteq \{1, \dots, 39\}, |I| \ge 1} (-1)^{|I|+1} \mathbf{1}_{S \subseteq \cup_{i \in I} E_i}
\Delta_{39, 3} &= \sum_{|S| = 3} \mathbb{P}(S \subseteq X) \sum_{I \subseteq \{1, \dots, 39\}, |I| \ge 1} (-1)^{|I|+1} \mathbf{1}_{S \subseteq \cup_{i \in I} E_i}
\Delta_{39, 4} &= \sum_{|S| = 4} \mathbb{P}(S \subseteq X) \sum_{I \subseteq \{1, \dots, 39\}, |I| \ge 1} (-1)^{|I|+1} \mathbf{1}_{S \subseteq \cup_{i \in I} E_i}
\Delta_{39, 5} &= \sum_{|S| = 5} \mathbb{P}(S \subseteq X) \sum_{I \subseteq \{1, \dots, 39\}, |I| \ge 1} (-1)^{|I|+1} \mathbf{1}_{S \subseteq \cup_{i \in I} E_i}
\Delta_{39, 6} &= \sum_{|S| = 6} \mathbb{P}(S \subseteq X) \sum_{I \subseteq \{1, \dots, 39\}, |I| \ge 1} (-1)^{|I|+1} \mathbf{1}_{S \subseteq \cup_{i \in I} E_i}
\Delta_{39, 7} &= \sum_{|S| = 7} \mathbb{P}(S \subseteq X) \sum_{I \subseteq \{1, \dots, 39\}, |I| \ge 1} (-1)^{|I|+1} \mathbf{1}_{S \subseteq \cup_{i \in I} E_i}
\Delta_{39, 8} &= \sum_{|S| = 8} \mathbb{P}(S \subseteq X) \sum_{I \subseteq \{1, \dots, 39\}, |I| \ge 1} (-1)^{|I|+1} \mathbf{1}_{S \subseteq \cup_{i \in I} E_i}
\Delta_{39, 9} &= \sum_{|S| = 9} \mathbb{P}(S \subseteq X) \sum_{I \subseteq \{1, \dots, 39\}, |I| \ge 1} (-1)^{|I|+1} \mathbf{1}_{S \subseteq \cup_{i \in I} E_i}\end{align}Cette majoration permet de borner supérieurement la corrélation locale. Par application du lemme de Lovász local ou des inégalités de concentration azumaiennes, la déviation par rapport à la moyenne produit un terme résiduel qui s'amortit exponentiellement. Nous fixons explicitement cette décroissance :
\begin{equation}\mathcal{R}_{39} \le \exp\left( - \frac{ \left( \Delta_{39, 1} \right)^2 }{ 2 \sum_{j=2}^{39} \Delta_{39, j} + \frac{1}{3} \max_{j} \Delta_{39, j} } \right)\end{equation}Le contrôle de ce terme $\mathcal{R}_{39}$ certifie l'indépendance quasi-totale des choix pour les structures de dimension 39.

\subsubsection{Analyse du moment d'ordre 40}Pour le degré 40, nous évaluons l'intersection simultanée de $\ell = 40$ hyperarêtes. Soient $E_1, \dots, E_{40}$ des arêtes de l'hypergraphe.
\begin{equation}\mathbb{E}\left[ \prod_{i=1}^{40} \mathbf{1}_{E_i \cap X \neq \emptyset} \right] = \sum_{j=1}^{40} (-1)^{j-1} \sum_{1 \le i_1 < \dots < i_j \le 40} \mathbb{P}\left( \bigcap_{m=1}^j (E_{i_m} \cap X \neq \emptyset) \right)\end{equation}Nous décomposons la probabilité conjointe en utilisant l'indépendance conditionnelle sur le choix des sommets de $X$. L'ensemble des configurations de recouvrement de l'union $\bigcup_{i} E_i$ est partitionné selon la taille de leur intersection avec $X$.
\begin{align} \\\Delta_{40, 1} &= \sum_{|S| = 1} \mathbb{P}(S \subseteq X) \sum_{I \subseteq \{1, \dots, 40\}, |I| \ge 1} (-1)^{|I|+1} \mathbf{1}_{S \subseteq \cup_{i \in I} E_i}
\Delta_{40, 2} &= \sum_{|S| = 2} \mathbb{P}(S \subseteq X) \sum_{I \subseteq \{1, \dots, 40\}, |I| \ge 1} (-1)^{|I|+1} \mathbf{1}_{S \subseteq \cup_{i \in I} E_i}
\Delta_{40, 3} &= \sum_{|S| = 3} \mathbb{P}(S \subseteq X) \sum_{I \subseteq \{1, \dots, 40\}, |I| \ge 1} (-1)^{|I|+1} \mathbf{1}_{S \subseteq \cup_{i \in I} E_i}
\Delta_{40, 4} &= \sum_{|S| = 4} \mathbb{P}(S \subseteq X) \sum_{I \subseteq \{1, \dots, 40\}, |I| \ge 1} (-1)^{|I|+1} \mathbf{1}_{S \subseteq \cup_{i \in I} E_i}
\Delta_{40, 5} &= \sum_{|S| = 5} \mathbb{P}(S \subseteq X) \sum_{I \subseteq \{1, \dots, 40\}, |I| \ge 1} (-1)^{|I|+1} \mathbf{1}_{S \subseteq \cup_{i \in I} E_i}
\Delta_{40, 6} &= \sum_{|S| = 6} \mathbb{P}(S \subseteq X) \sum_{I \subseteq \{1, \dots, 40\}, |I| \ge 1} (-1)^{|I|+1} \mathbf{1}_{S \subseteq \cup_{i \in I} E_i}
\Delta_{40, 7} &= \sum_{|S| = 7} \mathbb{P}(S \subseteq X) \sum_{I \subseteq \{1, \dots, 40\}, |I| \ge 1} (-1)^{|I|+1} \mathbf{1}_{S \subseteq \cup_{i \in I} E_i}
\Delta_{40, 8} &= \sum_{|S| = 8} \mathbb{P}(S \subseteq X) \sum_{I \subseteq \{1, \dots, 40\}, |I| \ge 1} (-1)^{|I|+1} \mathbf{1}_{S \subseteq \cup_{i \in I} E_i}
\Delta_{40, 9} &= \sum_{|S| = 9} \mathbb{P}(S \subseteq X) \sum_{I \subseteq \{1, \dots, 40\}, |I| \ge 1} (-1)^{|I|+1} \mathbf{1}_{S \subseteq \cup_{i \in I} E_i}\end{align}Cette majoration permet de borner supérieurement la corrélation locale. Par application du lemme de Lovász local ou des inégalités de concentration azumaiennes, la déviation par rapport à la moyenne produit un terme résiduel qui s'amortit exponentiellement. Nous fixons explicitement cette décroissance :
\begin{equation}\mathcal{R}_{40} \le \exp\left( - \frac{ \left( \Delta_{40, 1} \right)^2 }{ 2 \sum_{j=2}^{40} \Delta_{40, j} + \frac{1}{3} \max_{j} \Delta_{40, j} } \right)\end{equation}Le contrôle de ce terme $\mathcal{R}_{40}$ certifie l'indépendance quasi-totale des choix pour les structures de dimension 40.

\subsubsection{Analyse du moment d'ordre 41}Pour le degré 41, nous évaluons l'intersection simultanée de $\ell = 41$ hyperarêtes. Soient $E_1, \dots, E_{41}$ des arêtes de l'hypergraphe.
\begin{equation}\mathbb{E}\left[ \prod_{i=1}^{41} \mathbf{1}_{E_i \cap X \neq \emptyset} \right] = \sum_{j=1}^{41} (-1)^{j-1} \sum_{1 \le i_1 < \dots < i_j \le 41} \mathbb{P}\left( \bigcap_{m=1}^j (E_{i_m} \cap X \neq \emptyset) \right)\end{equation}Nous décomposons la probabilité conjointe en utilisant l'indépendance conditionnelle sur le choix des sommets de $X$. L'ensemble des configurations de recouvrement de l'union $\bigcup_{i} E_i$ est partitionné selon la taille de leur intersection avec $X$.
\begin{align} \\\Delta_{41, 1} &= \sum_{|S| = 1} \mathbb{P}(S \subseteq X) \sum_{I \subseteq \{1, \dots, 41\}, |I| \ge 1} (-1)^{|I|+1} \mathbf{1}_{S \subseteq \cup_{i \in I} E_i}
\Delta_{41, 2} &= \sum_{|S| = 2} \mathbb{P}(S \subseteq X) \sum_{I \subseteq \{1, \dots, 41\}, |I| \ge 1} (-1)^{|I|+1} \mathbf{1}_{S \subseteq \cup_{i \in I} E_i}
\Delta_{41, 3} &= \sum_{|S| = 3} \mathbb{P}(S \subseteq X) \sum_{I \subseteq \{1, \dots, 41\}, |I| \ge 1} (-1)^{|I|+1} \mathbf{1}_{S \subseteq \cup_{i \in I} E_i}
\Delta_{41, 4} &= \sum_{|S| = 4} \mathbb{P}(S \subseteq X) \sum_{I \subseteq \{1, \dots, 41\}, |I| \ge 1} (-1)^{|I|+1} \mathbf{1}_{S \subseteq \cup_{i \in I} E_i}
\Delta_{41, 5} &= \sum_{|S| = 5} \mathbb{P}(S \subseteq X) \sum_{I \subseteq \{1, \dots, 41\}, |I| \ge 1} (-1)^{|I|+1} \mathbf{1}_{S \subseteq \cup_{i \in I} E_i}
\Delta_{41, 6} &= \sum_{|S| = 6} \mathbb{P}(S \subseteq X) \sum_{I \subseteq \{1, \dots, 41\}, |I| \ge 1} (-1)^{|I|+1} \mathbf{1}_{S \subseteq \cup_{i \in I} E_i}
\Delta_{41, 7} &= \sum_{|S| = 7} \mathbb{P}(S \subseteq X) \sum_{I \subseteq \{1, \dots, 41\}, |I| \ge 1} (-1)^{|I|+1} \mathbf{1}_{S \subseteq \cup_{i \in I} E_i}
\Delta_{41, 8} &= \sum_{|S| = 8} \mathbb{P}(S \subseteq X) \sum_{I \subseteq \{1, \dots, 41\}, |I| \ge 1} (-1)^{|I|+1} \mathbf{1}_{S \subseteq \cup_{i \in I} E_i}
\Delta_{41, 9} &= \sum_{|S| = 9} \mathbb{P}(S \subseteq X) \sum_{I \subseteq \{1, \dots, 41\}, |I| \ge 1} (-1)^{|I|+1} \mathbf{1}_{S \subseteq \cup_{i \in I} E_i}\end{align}Cette majoration permet de borner supérieurement la corrélation locale. Par application du lemme de Lovász local ou des inégalités de concentration azumaiennes, la déviation par rapport à la moyenne produit un terme résiduel qui s'amortit exponentiellement. Nous fixons explicitement cette décroissance :
\begin{equation}\mathcal{R}_{41} \le \exp\left( - \frac{ \left( \Delta_{41, 1} \right)^2 }{ 2 \sum_{j=2}^{41} \Delta_{41, j} + \frac{1}{3} \max_{j} \Delta_{41, j} } \right)\end{equation}Le contrôle de ce terme $\mathcal{R}_{41}$ certifie l'indépendance quasi-totale des choix pour les structures de dimension 41.

\subsubsection{Analyse du moment d'ordre 42}Pour le degré 42, nous évaluons l'intersection simultanée de $\ell = 42$ hyperarêtes. Soient $E_1, \dots, E_{42}$ des arêtes de l'hypergraphe.
\begin{equation}\mathbb{E}\left[ \prod_{i=1}^{42} \mathbf{1}_{E_i \cap X \neq \emptyset} \right] = \sum_{j=1}^{42} (-1)^{j-1} \sum_{1 \le i_1 < \dots < i_j \le 42} \mathbb{P}\left( \bigcap_{m=1}^j (E_{i_m} \cap X \neq \emptyset) \right)\end{equation}Nous décomposons la probabilité conjointe en utilisant l'indépendance conditionnelle sur le choix des sommets de $X$. L'ensemble des configurations de recouvrement de l'union $\bigcup_{i} E_i$ est partitionné selon la taille de leur intersection avec $X$.
\begin{align} \\\Delta_{42, 1} &= \sum_{|S| = 1} \mathbb{P}(S \subseteq X) \sum_{I \subseteq \{1, \dots, 42\}, |I| \ge 1} (-1)^{|I|+1} \mathbf{1}_{S \subseteq \cup_{i \in I} E_i}
\Delta_{42, 2} &= \sum_{|S| = 2} \mathbb{P}(S \subseteq X) \sum_{I \subseteq \{1, \dots, 42\}, |I| \ge 1} (-1)^{|I|+1} \mathbf{1}_{S \subseteq \cup_{i \in I} E_i}
\Delta_{42, 3} &= \sum_{|S| = 3} \mathbb{P}(S \subseteq X) \sum_{I \subseteq \{1, \dots, 42\}, |I| \ge 1} (-1)^{|I|+1} \mathbf{1}_{S \subseteq \cup_{i \in I} E_i}
\Delta_{42, 4} &= \sum_{|S| = 4} \mathbb{P}(S \subseteq X) \sum_{I \subseteq \{1, \dots, 42\}, |I| \ge 1} (-1)^{|I|+1} \mathbf{1}_{S \subseteq \cup_{i \in I} E_i}
\Delta_{42, 5} &= \sum_{|S| = 5} \mathbb{P}(S \subseteq X) \sum_{I \subseteq \{1, \dots, 42\}, |I| \ge 1} (-1)^{|I|+1} \mathbf{1}_{S \subseteq \cup_{i \in I} E_i}
\Delta_{42, 6} &= \sum_{|S| = 6} \mathbb{P}(S \subseteq X) \sum_{I \subseteq \{1, \dots, 42\}, |I| \ge 1} (-1)^{|I|+1} \mathbf{1}_{S \subseteq \cup_{i \in I} E_i}
\Delta_{42, 7} &= \sum_{|S| = 7} \mathbb{P}(S \subseteq X) \sum_{I \subseteq \{1, \dots, 42\}, |I| \ge 1} (-1)^{|I|+1} \mathbf{1}_{S \subseteq \cup_{i \in I} E_i}
\Delta_{42, 8} &= \sum_{|S| = 8} \mathbb{P}(S \subseteq X) \sum_{I \subseteq \{1, \dots, 42\}, |I| \ge 1} (-1)^{|I|+1} \mathbf{1}_{S \subseteq \cup_{i \in I} E_i}
\Delta_{42, 9} &= \sum_{|S| = 9} \mathbb{P}(S \subseteq X) \sum_{I \subseteq \{1, \dots, 42\}, |I| \ge 1} (-1)^{|I|+1} \mathbf{1}_{S \subseteq \cup_{i \in I} E_i}\end{align}Cette majoration permet de borner supérieurement la corrélation locale. Par application du lemme de Lovász local ou des inégalités de concentration azumaiennes, la déviation par rapport à la moyenne produit un terme résiduel qui s'amortit exponentiellement. Nous fixons explicitement cette décroissance :
\begin{equation}\mathcal{R}_{42} \le \exp\left( - \frac{ \left( \Delta_{42, 1} \right)^2 }{ 2 \sum_{j=2}^{42} \Delta_{42, j} + \frac{1}{3} \max_{j} \Delta_{42, j} } \right)\end{equation}Le contrôle de ce terme $\mathcal{R}_{42}$ certifie l'indépendance quasi-totale des choix pour les structures de dimension 42.

\subsubsection{Analyse du moment d'ordre 43}Pour le degré 43, nous évaluons l'intersection simultanée de $\ell = 43$ hyperarêtes. Soient $E_1, \dots, E_{43}$ des arêtes de l'hypergraphe.
\begin{equation}\mathbb{E}\left[ \prod_{i=1}^{43} \mathbf{1}_{E_i \cap X \neq \emptyset} \right] = \sum_{j=1}^{43} (-1)^{j-1} \sum_{1 \le i_1 < \dots < i_j \le 43} \mathbb{P}\left( \bigcap_{m=1}^j (E_{i_m} \cap X \neq \emptyset) \right)\end{equation}Nous décomposons la probabilité conjointe en utilisant l'indépendance conditionnelle sur le choix des sommets de $X$. L'ensemble des configurations de recouvrement de l'union $\bigcup_{i} E_i$ est partitionné selon la taille de leur intersection avec $X$.
\begin{align} \\\Delta_{43, 1} &= \sum_{|S| = 1} \mathbb{P}(S \subseteq X) \sum_{I \subseteq \{1, \dots, 43\}, |I| \ge 1} (-1)^{|I|+1} \mathbf{1}_{S \subseteq \cup_{i \in I} E_i}
\Delta_{43, 2} &= \sum_{|S| = 2} \mathbb{P}(S \subseteq X) \sum_{I \subseteq \{1, \dots, 43\}, |I| \ge 1} (-1)^{|I|+1} \mathbf{1}_{S \subseteq \cup_{i \in I} E_i}
\Delta_{43, 3} &= \sum_{|S| = 3} \mathbb{P}(S \subseteq X) \sum_{I \subseteq \{1, \dots, 43\}, |I| \ge 1} (-1)^{|I|+1} \mathbf{1}_{S \subseteq \cup_{i \in I} E_i}
\Delta_{43, 4} &= \sum_{|S| = 4} \mathbb{P}(S \subseteq X) \sum_{I \subseteq \{1, \dots, 43\}, |I| \ge 1} (-1)^{|I|+1} \mathbf{1}_{S \subseteq \cup_{i \in I} E_i}
\Delta_{43, 5} &= \sum_{|S| = 5} \mathbb{P}(S \subseteq X) \sum_{I \subseteq \{1, \dots, 43\}, |I| \ge 1} (-1)^{|I|+1} \mathbf{1}_{S \subseteq \cup_{i \in I} E_i}
\Delta_{43, 6} &= \sum_{|S| = 6} \mathbb{P}(S \subseteq X) \sum_{I \subseteq \{1, \dots, 43\}, |I| \ge 1} (-1)^{|I|+1} \mathbf{1}_{S \subseteq \cup_{i \in I} E_i}
\Delta_{43, 7} &= \sum_{|S| = 7} \mathbb{P}(S \subseteq X) \sum_{I \subseteq \{1, \dots, 43\}, |I| \ge 1} (-1)^{|I|+1} \mathbf{1}_{S \subseteq \cup_{i \in I} E_i}
\Delta_{43, 8} &= \sum_{|S| = 8} \mathbb{P}(S \subseteq X) \sum_{I \subseteq \{1, \dots, 43\}, |I| \ge 1} (-1)^{|I|+1} \mathbf{1}_{S \subseteq \cup_{i \in I} E_i}
\Delta_{43, 9} &= \sum_{|S| = 9} \mathbb{P}(S \subseteq X) \sum_{I \subseteq \{1, \dots, 43\}, |I| \ge 1} (-1)^{|I|+1} \mathbf{1}_{S \subseteq \cup_{i \in I} E_i}\end{align}Cette majoration permet de borner supérieurement la corrélation locale. Par application du lemme de Lovász local ou des inégalités de concentration azumaiennes, la déviation par rapport à la moyenne produit un terme résiduel qui s'amortit exponentiellement. Nous fixons explicitement cette décroissance :
\begin{equation}\mathcal{R}_{43} \le \exp\left( - \frac{ \left( \Delta_{43, 1} \right)^2 }{ 2 \sum_{j=2}^{43} \Delta_{43, j} + \frac{1}{3} \max_{j} \Delta_{43, j} } \right)\end{equation}Le contrôle de ce terme $\mathcal{R}_{43}$ certifie l'indépendance quasi-totale des choix pour les structures de dimension 43.

\subsubsection{Analyse du moment d'ordre 44}Pour le degré 44, nous évaluons l'intersection simultanée de $\ell = 44$ hyperarêtes. Soient $E_1, \dots, E_{44}$ des arêtes de l'hypergraphe.
\begin{equation}\mathbb{E}\left[ \prod_{i=1}^{44} \mathbf{1}_{E_i \cap X \neq \emptyset} \right] = \sum_{j=1}^{44} (-1)^{j-1} \sum_{1 \le i_1 < \dots < i_j \le 44} \mathbb{P}\left( \bigcap_{m=1}^j (E_{i_m} \cap X \neq \emptyset) \right)\end{equation}Nous décomposons la probabilité conjointe en utilisant l'indépendance conditionnelle sur le choix des sommets de $X$. L'ensemble des configurations de recouvrement de l'union $\bigcup_{i} E_i$ est partitionné selon la taille de leur intersection avec $X$.
\begin{align} \\\Delta_{44, 1} &= \sum_{|S| = 1} \mathbb{P}(S \subseteq X) \sum_{I \subseteq \{1, \dots, 44\}, |I| \ge 1} (-1)^{|I|+1} \mathbf{1}_{S \subseteq \cup_{i \in I} E_i}
\Delta_{44, 2} &= \sum_{|S| = 2} \mathbb{P}(S \subseteq X) \sum_{I \subseteq \{1, \dots, 44\}, |I| \ge 1} (-1)^{|I|+1} \mathbf{1}_{S \subseteq \cup_{i \in I} E_i}
\Delta_{44, 3} &= \sum_{|S| = 3} \mathbb{P}(S \subseteq X) \sum_{I \subseteq \{1, \dots, 44\}, |I| \ge 1} (-1)^{|I|+1} \mathbf{1}_{S \subseteq \cup_{i \in I} E_i}
\Delta_{44, 4} &= \sum_{|S| = 4} \mathbb{P}(S \subseteq X) \sum_{I \subseteq \{1, \dots, 44\}, |I| \ge 1} (-1)^{|I|+1} \mathbf{1}_{S \subseteq \cup_{i \in I} E_i}
\Delta_{44, 5} &= \sum_{|S| = 5} \mathbb{P}(S \subseteq X) \sum_{I \subseteq \{1, \dots, 44\}, |I| \ge 1} (-1)^{|I|+1} \mathbf{1}_{S \subseteq \cup_{i \in I} E_i}
\Delta_{44, 6} &= \sum_{|S| = 6} \mathbb{P}(S \subseteq X) \sum_{I \subseteq \{1, \dots, 44\}, |I| \ge 1} (-1)^{|I|+1} \mathbf{1}_{S \subseteq \cup_{i \in I} E_i}
\Delta_{44, 7} &= \sum_{|S| = 7} \mathbb{P}(S \subseteq X) \sum_{I \subseteq \{1, \dots, 44\}, |I| \ge 1} (-1)^{|I|+1} \mathbf{1}_{S \subseteq \cup_{i \in I} E_i}
\Delta_{44, 8} &= \sum_{|S| = 8} \mathbb{P}(S \subseteq X) \sum_{I \subseteq \{1, \dots, 44\}, |I| \ge 1} (-1)^{|I|+1} \mathbf{1}_{S \subseteq \cup_{i \in I} E_i}
\Delta_{44, 9} &= \sum_{|S| = 9} \mathbb{P}(S \subseteq X) \sum_{I \subseteq \{1, \dots, 44\}, |I| \ge 1} (-1)^{|I|+1} \mathbf{1}_{S \subseteq \cup_{i \in I} E_i}\end{align}Cette majoration permet de borner supérieurement la corrélation locale. Par application du lemme de Lovász local ou des inégalités de concentration azumaiennes, la déviation par rapport à la moyenne produit un terme résiduel qui s'amortit exponentiellement. Nous fixons explicitement cette décroissance :
\begin{equation}\mathcal{R}_{44} \le \exp\left( - \frac{ \left( \Delta_{44, 1} \right)^2 }{ 2 \sum_{j=2}^{44} \Delta_{44, j} + \frac{1}{3} \max_{j} \Delta_{44, j} } \right)\end{equation}Le contrôle de ce terme $\mathcal{R}_{44}$ certifie l'indépendance quasi-totale des choix pour les structures de dimension 44.

\subsubsection{Analyse du moment d'ordre 45}Pour le degré 45, nous évaluons l'intersection simultanée de $\ell = 45$ hyperarêtes. Soient $E_1, \dots, E_{45}$ des arêtes de l'hypergraphe.
\begin{equation}\mathbb{E}\left[ \prod_{i=1}^{45} \mathbf{1}_{E_i \cap X \neq \emptyset} \right] = \sum_{j=1}^{45} (-1)^{j-1} \sum_{1 \le i_1 < \dots < i_j \le 45} \mathbb{P}\left( \bigcap_{m=1}^j (E_{i_m} \cap X \neq \emptyset) \right)\end{equation}Nous décomposons la probabilité conjointe en utilisant l'indépendance conditionnelle sur le choix des sommets de $X$. L'ensemble des configurations de recouvrement de l'union $\bigcup_{i} E_i$ est partitionné selon la taille de leur intersection avec $X$.
\begin{align} \\\Delta_{45, 1} &= \sum_{|S| = 1} \mathbb{P}(S \subseteq X) \sum_{I \subseteq \{1, \dots, 45\}, |I| \ge 1} (-1)^{|I|+1} \mathbf{1}_{S \subseteq \cup_{i \in I} E_i}
\Delta_{45, 2} &= \sum_{|S| = 2} \mathbb{P}(S \subseteq X) \sum_{I \subseteq \{1, \dots, 45\}, |I| \ge 1} (-1)^{|I|+1} \mathbf{1}_{S \subseteq \cup_{i \in I} E_i}
\Delta_{45, 3} &= \sum_{|S| = 3} \mathbb{P}(S \subseteq X) \sum_{I \subseteq \{1, \dots, 45\}, |I| \ge 1} (-1)^{|I|+1} \mathbf{1}_{S \subseteq \cup_{i \in I} E_i}
\Delta_{45, 4} &= \sum_{|S| = 4} \mathbb{P}(S \subseteq X) \sum_{I \subseteq \{1, \dots, 45\}, |I| \ge 1} (-1)^{|I|+1} \mathbf{1}_{S \subseteq \cup_{i \in I} E_i}
\Delta_{45, 5} &= \sum_{|S| = 5} \mathbb{P}(S \subseteq X) \sum_{I \subseteq \{1, \dots, 45\}, |I| \ge 1} (-1)^{|I|+1} \mathbf{1}_{S \subseteq \cup_{i \in I} E_i}
\Delta_{45, 6} &= \sum_{|S| = 6} \mathbb{P}(S \subseteq X) \sum_{I \subseteq \{1, \dots, 45\}, |I| \ge 1} (-1)^{|I|+1} \mathbf{1}_{S \subseteq \cup_{i \in I} E_i}
\Delta_{45, 7} &= \sum_{|S| = 7} \mathbb{P}(S \subseteq X) \sum_{I \subseteq \{1, \dots, 45\}, |I| \ge 1} (-1)^{|I|+1} \mathbf{1}_{S \subseteq \cup_{i \in I} E_i}
\Delta_{45, 8} &= \sum_{|S| = 8} \mathbb{P}(S \subseteq X) \sum_{I \subseteq \{1, \dots, 45\}, |I| \ge 1} (-1)^{|I|+1} \mathbf{1}_{S \subseteq \cup_{i \in I} E_i}
\Delta_{45, 9} &= \sum_{|S| = 9} \mathbb{P}(S \subseteq X) \sum_{I \subseteq \{1, \dots, 45\}, |I| \ge 1} (-1)^{|I|+1} \mathbf{1}_{S \subseteq \cup_{i \in I} E_i}\end{align}Cette majoration permet de borner supérieurement la corrélation locale. Par application du lemme de Lovász local ou des inégalités de concentration azumaiennes, la déviation par rapport à la moyenne produit un terme résiduel qui s'amortit exponentiellement. Nous fixons explicitement cette décroissance :
\begin{equation}\mathcal{R}_{45} \le \exp\left( - \frac{ \left( \Delta_{45, 1} \right)^2 }{ 2 \sum_{j=2}^{45} \Delta_{45, j} + \frac{1}{3} \max_{j} \Delta_{45, j} } \right)\end{equation}Le contrôle de ce terme $\mathcal{R}_{45}$ certifie l'indépendance quasi-totale des choix pour les structures de dimension 45.

\subsubsection{Analyse du moment d'ordre 46}Pour le degré 46, nous évaluons l'intersection simultanée de $\ell = 46$ hyperarêtes. Soient $E_1, \dots, E_{46}$ des arêtes de l'hypergraphe.
\begin{equation}\mathbb{E}\left[ \prod_{i=1}^{46} \mathbf{1}_{E_i \cap X \neq \emptyset} \right] = \sum_{j=1}^{46} (-1)^{j-1} \sum_{1 \le i_1 < \dots < i_j \le 46} \mathbb{P}\left( \bigcap_{m=1}^j (E_{i_m} \cap X \neq \emptyset) \right)\end{equation}Nous décomposons la probabilité conjointe en utilisant l'indépendance conditionnelle sur le choix des sommets de $X$. L'ensemble des configurations de recouvrement de l'union $\bigcup_{i} E_i$ est partitionné selon la taille de leur intersection avec $X$.
\begin{align} \\\Delta_{46, 1} &= \sum_{|S| = 1} \mathbb{P}(S \subseteq X) \sum_{I \subseteq \{1, \dots, 46\}, |I| \ge 1} (-1)^{|I|+1} \mathbf{1}_{S \subseteq \cup_{i \in I} E_i}
\Delta_{46, 2} &= \sum_{|S| = 2} \mathbb{P}(S \subseteq X) \sum_{I \subseteq \{1, \dots, 46\}, |I| \ge 1} (-1)^{|I|+1} \mathbf{1}_{S \subseteq \cup_{i \in I} E_i}
\Delta_{46, 3} &= \sum_{|S| = 3} \mathbb{P}(S \subseteq X) \sum_{I \subseteq \{1, \dots, 46\}, |I| \ge 1} (-1)^{|I|+1} \mathbf{1}_{S \subseteq \cup_{i \in I} E_i}
\Delta_{46, 4} &= \sum_{|S| = 4} \mathbb{P}(S \subseteq X) \sum_{I \subseteq \{1, \dots, 46\}, |I| \ge 1} (-1)^{|I|+1} \mathbf{1}_{S \subseteq \cup_{i \in I} E_i}
\Delta_{46, 5} &= \sum_{|S| = 5} \mathbb{P}(S \subseteq X) \sum_{I \subseteq \{1, \dots, 46\}, |I| \ge 1} (-1)^{|I|+1} \mathbf{1}_{S \subseteq \cup_{i \in I} E_i}
\Delta_{46, 6} &= \sum_{|S| = 6} \mathbb{P}(S \subseteq X) \sum_{I \subseteq \{1, \dots, 46\}, |I| \ge 1} (-1)^{|I|+1} \mathbf{1}_{S \subseteq \cup_{i \in I} E_i}
\Delta_{46, 7} &= \sum_{|S| = 7} \mathbb{P}(S \subseteq X) \sum_{I \subseteq \{1, \dots, 46\}, |I| \ge 1} (-1)^{|I|+1} \mathbf{1}_{S \subseteq \cup_{i \in I} E_i}
\Delta_{46, 8} &= \sum_{|S| = 8} \mathbb{P}(S \subseteq X) \sum_{I \subseteq \{1, \dots, 46\}, |I| \ge 1} (-1)^{|I|+1} \mathbf{1}_{S \subseteq \cup_{i \in I} E_i}
\Delta_{46, 9} &= \sum_{|S| = 9} \mathbb{P}(S \subseteq X) \sum_{I \subseteq \{1, \dots, 46\}, |I| \ge 1} (-1)^{|I|+1} \mathbf{1}_{S \subseteq \cup_{i \in I} E_i}\end{align}Cette majoration permet de borner supérieurement la corrélation locale. Par application du lemme de Lovász local ou des inégalités de concentration azumaiennes, la déviation par rapport à la moyenne produit un terme résiduel qui s'amortit exponentiellement. Nous fixons explicitement cette décroissance :
\begin{equation}\mathcal{R}_{46} \le \exp\left( - \frac{ \left( \Delta_{46, 1} \right)^2 }{ 2 \sum_{j=2}^{46} \Delta_{46, j} + \frac{1}{3} \max_{j} \Delta_{46, j} } \right)\end{equation}Le contrôle de ce terme $\mathcal{R}_{46}$ certifie l'indépendance quasi-totale des choix pour les structures de dimension 46.

\subsubsection{Analyse du moment d'ordre 47}Pour le degré 47, nous évaluons l'intersection simultanée de $\ell = 47$ hyperarêtes. Soient $E_1, \dots, E_{47}$ des arêtes de l'hypergraphe.
\begin{equation}\mathbb{E}\left[ \prod_{i=1}^{47} \mathbf{1}_{E_i \cap X \neq \emptyset} \right] = \sum_{j=1}^{47} (-1)^{j-1} \sum_{1 \le i_1 < \dots < i_j \le 47} \mathbb{P}\left( \bigcap_{m=1}^j (E_{i_m} \cap X \neq \emptyset) \right)\end{equation}Nous décomposons la probabilité conjointe en utilisant l'indépendance conditionnelle sur le choix des sommets de $X$. L'ensemble des configurations de recouvrement de l'union $\bigcup_{i} E_i$ est partitionné selon la taille de leur intersection avec $X$.
\begin{align} \\\Delta_{47, 1} &= \sum_{|S| = 1} \mathbb{P}(S \subseteq X) \sum_{I \subseteq \{1, \dots, 47\}, |I| \ge 1} (-1)^{|I|+1} \mathbf{1}_{S \subseteq \cup_{i \in I} E_i}
\Delta_{47, 2} &= \sum_{|S| = 2} \mathbb{P}(S \subseteq X) \sum_{I \subseteq \{1, \dots, 47\}, |I| \ge 1} (-1)^{|I|+1} \mathbf{1}_{S \subseteq \cup_{i \in I} E_i}
\Delta_{47, 3} &= \sum_{|S| = 3} \mathbb{P}(S \subseteq X) \sum_{I \subseteq \{1, \dots, 47\}, |I| \ge 1} (-1)^{|I|+1} \mathbf{1}_{S \subseteq \cup_{i \in I} E_i}
\Delta_{47, 4} &= \sum_{|S| = 4} \mathbb{P}(S \subseteq X) \sum_{I \subseteq \{1, \dots, 47\}, |I| \ge 1} (-1)^{|I|+1} \mathbf{1}_{S \subseteq \cup_{i \in I} E_i}
\Delta_{47, 5} &= \sum_{|S| = 5} \mathbb{P}(S \subseteq X) \sum_{I \subseteq \{1, \dots, 47\}, |I| \ge 1} (-1)^{|I|+1} \mathbf{1}_{S \subseteq \cup_{i \in I} E_i}
\Delta_{47, 6} &= \sum_{|S| = 6} \mathbb{P}(S \subseteq X) \sum_{I \subseteq \{1, \dots, 47\}, |I| \ge 1} (-1)^{|I|+1} \mathbf{1}_{S \subseteq \cup_{i \in I} E_i}
\Delta_{47, 7} &= \sum_{|S| = 7} \mathbb{P}(S \subseteq X) \sum_{I \subseteq \{1, \dots, 47\}, |I| \ge 1} (-1)^{|I|+1} \mathbf{1}_{S \subseteq \cup_{i \in I} E_i}
\Delta_{47, 8} &= \sum_{|S| = 8} \mathbb{P}(S \subseteq X) \sum_{I \subseteq \{1, \dots, 47\}, |I| \ge 1} (-1)^{|I|+1} \mathbf{1}_{S \subseteq \cup_{i \in I} E_i}
\Delta_{47, 9} &= \sum_{|S| = 9} \mathbb{P}(S \subseteq X) \sum_{I \subseteq \{1, \dots, 47\}, |I| \ge 1} (-1)^{|I|+1} \mathbf{1}_{S \subseteq \cup_{i \in I} E_i}\end{align}Cette majoration permet de borner supérieurement la corrélation locale. Par application du lemme de Lovász local ou des inégalités de concentration azumaiennes, la déviation par rapport à la moyenne produit un terme résiduel qui s'amortit exponentiellement. Nous fixons explicitement cette décroissance :
\begin{equation}\mathcal{R}_{47} \le \exp\left( - \frac{ \left( \Delta_{47, 1} \right)^2 }{ 2 \sum_{j=2}^{47} \Delta_{47, j} + \frac{1}{3} \max_{j} \Delta_{47, j} } \right)\end{equation}Le contrôle de ce terme $\mathcal{R}_{47}$ certifie l'indépendance quasi-totale des choix pour les structures de dimension 47.

\subsubsection{Analyse du moment d'ordre 48}Pour le degré 48, nous évaluons l'intersection simultanée de $\ell = 48$ hyperarêtes. Soient $E_1, \dots, E_{48}$ des arêtes de l'hypergraphe.
\begin{equation}\mathbb{E}\left[ \prod_{i=1}^{48} \mathbf{1}_{E_i \cap X \neq \emptyset} \right] = \sum_{j=1}^{48} (-1)^{j-1} \sum_{1 \le i_1 < \dots < i_j \le 48} \mathbb{P}\left( \bigcap_{m=1}^j (E_{i_m} \cap X \neq \emptyset) \right)\end{equation}Nous décomposons la probabilité conjointe en utilisant l'indépendance conditionnelle sur le choix des sommets de $X$. L'ensemble des configurations de recouvrement de l'union $\bigcup_{i} E_i$ est partitionné selon la taille de leur intersection avec $X$.
\begin{align} \\\Delta_{48, 1} &= \sum_{|S| = 1} \mathbb{P}(S \subseteq X) \sum_{I \subseteq \{1, \dots, 48\}, |I| \ge 1} (-1)^{|I|+1} \mathbf{1}_{S \subseteq \cup_{i \in I} E_i}
\Delta_{48, 2} &= \sum_{|S| = 2} \mathbb{P}(S \subseteq X) \sum_{I \subseteq \{1, \dots, 48\}, |I| \ge 1} (-1)^{|I|+1} \mathbf{1}_{S \subseteq \cup_{i \in I} E_i}
\Delta_{48, 3} &= \sum_{|S| = 3} \mathbb{P}(S \subseteq X) \sum_{I \subseteq \{1, \dots, 48\}, |I| \ge 1} (-1)^{|I|+1} \mathbf{1}_{S \subseteq \cup_{i \in I} E_i}
\Delta_{48, 4} &= \sum_{|S| = 4} \mathbb{P}(S \subseteq X) \sum_{I \subseteq \{1, \dots, 48\}, |I| \ge 1} (-1)^{|I|+1} \mathbf{1}_{S \subseteq \cup_{i \in I} E_i}
\Delta_{48, 5} &= \sum_{|S| = 5} \mathbb{P}(S \subseteq X) \sum_{I \subseteq \{1, \dots, 48\}, |I| \ge 1} (-1)^{|I|+1} \mathbf{1}_{S \subseteq \cup_{i \in I} E_i}
\Delta_{48, 6} &= \sum_{|S| = 6} \mathbb{P}(S \subseteq X) \sum_{I \subseteq \{1, \dots, 48\}, |I| \ge 1} (-1)^{|I|+1} \mathbf{1}_{S \subseteq \cup_{i \in I} E_i}
\Delta_{48, 7} &= \sum_{|S| = 7} \mathbb{P}(S \subseteq X) \sum_{I \subseteq \{1, \dots, 48\}, |I| \ge 1} (-1)^{|I|+1} \mathbf{1}_{S \subseteq \cup_{i \in I} E_i}
\Delta_{48, 8} &= \sum_{|S| = 8} \mathbb{P}(S \subseteq X) \sum_{I \subseteq \{1, \dots, 48\}, |I| \ge 1} (-1)^{|I|+1} \mathbf{1}_{S \subseteq \cup_{i \in I} E_i}
\Delta_{48, 9} &= \sum_{|S| = 9} \mathbb{P}(S \subseteq X) \sum_{I \subseteq \{1, \dots, 48\}, |I| \ge 1} (-1)^{|I|+1} \mathbf{1}_{S \subseteq \cup_{i \in I} E_i}\end{align}Cette majoration permet de borner supérieurement la corrélation locale. Par application du lemme de Lovász local ou des inégalités de concentration azumaiennes, la déviation par rapport à la moyenne produit un terme résiduel qui s'amortit exponentiellement. Nous fixons explicitement cette décroissance :
\begin{equation}\mathcal{R}_{48} \le \exp\left( - \frac{ \left( \Delta_{48, 1} \right)^2 }{ 2 \sum_{j=2}^{48} \Delta_{48, j} + \frac{1}{3} \max_{j} \Delta_{48, j} } \right)\end{equation}Le contrôle de ce terme $\mathcal{R}_{48}$ certifie l'indépendance quasi-totale des choix pour les structures de dimension 48.

\subsubsection{Analyse du moment d'ordre 49}Pour le degré 49, nous évaluons l'intersection simultanée de $\ell = 49$ hyperarêtes. Soient $E_1, \dots, E_{49}$ des arêtes de l'hypergraphe.
\begin{equation}\mathbb{E}\left[ \prod_{i=1}^{49} \mathbf{1}_{E_i \cap X \neq \emptyset} \right] = \sum_{j=1}^{49} (-1)^{j-1} \sum_{1 \le i_1 < \dots < i_j \le 49} \mathbb{P}\left( \bigcap_{m=1}^j (E_{i_m} \cap X \neq \emptyset) \right)\end{equation}Nous décomposons la probabilité conjointe en utilisant l'indépendance conditionnelle sur le choix des sommets de $X$. L'ensemble des configurations de recouvrement de l'union $\bigcup_{i} E_i$ est partitionné selon la taille de leur intersection avec $X$.
\begin{align} \\\Delta_{49, 1} &= \sum_{|S| = 1} \mathbb{P}(S \subseteq X) \sum_{I \subseteq \{1, \dots, 49\}, |I| \ge 1} (-1)^{|I|+1} \mathbf{1}_{S \subseteq \cup_{i \in I} E_i}
\Delta_{49, 2} &= \sum_{|S| = 2} \mathbb{P}(S \subseteq X) \sum_{I \subseteq \{1, \dots, 49\}, |I| \ge 1} (-1)^{|I|+1} \mathbf{1}_{S \subseteq \cup_{i \in I} E_i}
\Delta_{49, 3} &= \sum_{|S| = 3} \mathbb{P}(S \subseteq X) \sum_{I \subseteq \{1, \dots, 49\}, |I| \ge 1} (-1)^{|I|+1} \mathbf{1}_{S \subseteq \cup_{i \in I} E_i}
\Delta_{49, 4} &= \sum_{|S| = 4} \mathbb{P}(S \subseteq X) \sum_{I \subseteq \{1, \dots, 49\}, |I| \ge 1} (-1)^{|I|+1} \mathbf{1}_{S \subseteq \cup_{i \in I} E_i}
\Delta_{49, 5} &= \sum_{|S| = 5} \mathbb{P}(S \subseteq X) \sum_{I \subseteq \{1, \dots, 49\}, |I| \ge 1} (-1)^{|I|+1} \mathbf{1}_{S \subseteq \cup_{i \in I} E_i}
\Delta_{49, 6} &= \sum_{|S| = 6} \mathbb{P}(S \subseteq X) \sum_{I \subseteq \{1, \dots, 49\}, |I| \ge 1} (-1)^{|I|+1} \mathbf{1}_{S \subseteq \cup_{i \in I} E_i}
\Delta_{49, 7} &= \sum_{|S| = 7} \mathbb{P}(S \subseteq X) \sum_{I \subseteq \{1, \dots, 49\}, |I| \ge 1} (-1)^{|I|+1} \mathbf{1}_{S \subseteq \cup_{i \in I} E_i}
\Delta_{49, 8} &= \sum_{|S| = 8} \mathbb{P}(S \subseteq X) \sum_{I \subseteq \{1, \dots, 49\}, |I| \ge 1} (-1)^{|I|+1} \mathbf{1}_{S \subseteq \cup_{i \in I} E_i}
\Delta_{49, 9} &= \sum_{|S| = 9} \mathbb{P}(S \subseteq X) \sum_{I \subseteq \{1, \dots, 49\}, |I| \ge 1} (-1)^{|I|+1} \mathbf{1}_{S \subseteq \cup_{i \in I} E_i}\end{align}Cette majoration permet de borner supérieurement la corrélation locale. Par application du lemme de Lovász local ou des inégalités de concentration azumaiennes, la déviation par rapport à la moyenne produit un terme résiduel qui s'amortit exponentiellement. Nous fixons explicitement cette décroissance :
\begin{equation}\mathcal{R}_{49} \le \exp\left( - \frac{ \left( \Delta_{49, 1} \right)^2 }{ 2 \sum_{j=2}^{49} \Delta_{49, j} + \frac{1}{3} \max_{j} \Delta_{49, j} } \right)\end{equation}Le contrôle de ce terme $\mathcal{R}_{49}$ certifie l'indépendance quasi-totale des choix pour les structures de dimension 49.

\subsubsection{Analyse du moment d'ordre 50}Pour le degré 50, nous évaluons l'intersection simultanée de $\ell = 50$ hyperarêtes. Soient $E_1, \dots, E_{50}$ des arêtes de l'hypergraphe.
\begin{equation}\mathbb{E}\left[ \prod_{i=1}^{50} \mathbf{1}_{E_i \cap X \neq \emptyset} \right] = \sum_{j=1}^{50} (-1)^{j-1} \sum_{1 \le i_1 < \dots < i_j \le 50} \mathbb{P}\left( \bigcap_{m=1}^j (E_{i_m} \cap X \neq \emptyset) \right)\end{equation}Nous décomposons la probabilité conjointe en utilisant l'indépendance conditionnelle sur le choix des sommets de $X$. L'ensemble des configurations de recouvrement de l'union $\bigcup_{i} E_i$ est partitionné selon la taille de leur intersection avec $X$.
\begin{align} \\\Delta_{50, 1} &= \sum_{|S| = 1} \mathbb{P}(S \subseteq X) \sum_{I \subseteq \{1, \dots, 50\}, |I| \ge 1} (-1)^{|I|+1} \mathbf{1}_{S \subseteq \cup_{i \in I} E_i}
\Delta_{50, 2} &= \sum_{|S| = 2} \mathbb{P}(S \subseteq X) \sum_{I \subseteq \{1, \dots, 50\}, |I| \ge 1} (-1)^{|I|+1} \mathbf{1}_{S \subseteq \cup_{i \in I} E_i}
\Delta_{50, 3} &= \sum_{|S| = 3} \mathbb{P}(S \subseteq X) \sum_{I \subseteq \{1, \dots, 50\}, |I| \ge 1} (-1)^{|I|+1} \mathbf{1}_{S \subseteq \cup_{i \in I} E_i}
\Delta_{50, 4} &= \sum_{|S| = 4} \mathbb{P}(S \subseteq X) \sum_{I \subseteq \{1, \dots, 50\}, |I| \ge 1} (-1)^{|I|+1} \mathbf{1}_{S \subseteq \cup_{i \in I} E_i}
\Delta_{50, 5} &= \sum_{|S| = 5} \mathbb{P}(S \subseteq X) \sum_{I \subseteq \{1, \dots, 50\}, |I| \ge 1} (-1)^{|I|+1} \mathbf{1}_{S \subseteq \cup_{i \in I} E_i}
\Delta_{50, 6} &= \sum_{|S| = 6} \mathbb{P}(S \subseteq X) \sum_{I \subseteq \{1, \dots, 50\}, |I| \ge 1} (-1)^{|I|+1} \mathbf{1}_{S \subseteq \cup_{i \in I} E_i}
\Delta_{50, 7} &= \sum_{|S| = 7} \mathbb{P}(S \subseteq X) \sum_{I \subseteq \{1, \dots, 50\}, |I| \ge 1} (-1)^{|I|+1} \mathbf{1}_{S \subseteq \cup_{i \in I} E_i}
\Delta_{50, 8} &= \sum_{|S| = 8} \mathbb{P}(S \subseteq X) \sum_{I \subseteq \{1, \dots, 50\}, |I| \ge 1} (-1)^{|I|+1} \mathbf{1}_{S \subseteq \cup_{i \in I} E_i}
\Delta_{50, 9} &= \sum_{|S| = 9} \mathbb{P}(S \subseteq X) \sum_{I \subseteq \{1, \dots, 50\}, |I| \ge 1} (-1)^{|I|+1} \mathbf{1}_{S \subseteq \cup_{i \in I} E_i}\end{align}Cette majoration permet de borner supérieurement la corrélation locale. Par application du lemme de Lovász local ou des inégalités de concentration azumaiennes, la déviation par rapport à la moyenne produit un terme résiduel qui s'amortit exponentiellement. Nous fixons explicitement cette décroissance :
\begin{equation}\mathcal{R}_{50} \le \exp\left( - \frac{ \left( \Delta_{50, 1} \right)^2 }{ 2 \sum_{j=2}^{50} \Delta_{50, j} + \frac{1}{3} \max_{j} \Delta_{50, j} } \right)\end{equation}Le contrôle de ce terme $\mathcal{R}_{50}$ certifie l'indépendance quasi-totale des choix pour les structures de dimension 50.

\subsubsection{Analyse du moment d'ordre 51}Pour le degré 51, nous évaluons l'intersection simultanée de $\ell = 51$ hyperarêtes. Soient $E_1, \dots, E_{51}$ des arêtes de l'hypergraphe.
\begin{equation}\mathbb{E}\left[ \prod_{i=1}^{51} \mathbf{1}_{E_i \cap X \neq \emptyset} \right] = \sum_{j=1}^{51} (-1)^{j-1} \sum_{1 \le i_1 < \dots < i_j \le 51} \mathbb{P}\left( \bigcap_{m=1}^j (E_{i_m} \cap X \neq \emptyset) \right)\end{equation}Nous décomposons la probabilité conjointe en utilisant l'indépendance conditionnelle sur le choix des sommets de $X$. L'ensemble des configurations de recouvrement de l'union $\bigcup_{i} E_i$ est partitionné selon la taille de leur intersection avec $X$.
\begin{align} \\\Delta_{51, 1} &= \sum_{|S| = 1} \mathbb{P}(S \subseteq X) \sum_{I \subseteq \{1, \dots, 51\}, |I| \ge 1} (-1)^{|I|+1} \mathbf{1}_{S \subseteq \cup_{i \in I} E_i}
\Delta_{51, 2} &= \sum_{|S| = 2} \mathbb{P}(S \subseteq X) \sum_{I \subseteq \{1, \dots, 51\}, |I| \ge 1} (-1)^{|I|+1} \mathbf{1}_{S \subseteq \cup_{i \in I} E_i}
\Delta_{51, 3} &= \sum_{|S| = 3} \mathbb{P}(S \subseteq X) \sum_{I \subseteq \{1, \dots, 51\}, |I| \ge 1} (-1)^{|I|+1} \mathbf{1}_{S \subseteq \cup_{i \in I} E_i}
\Delta_{51, 4} &= \sum_{|S| = 4} \mathbb{P}(S \subseteq X) \sum_{I \subseteq \{1, \dots, 51\}, |I| \ge 1} (-1)^{|I|+1} \mathbf{1}_{S \subseteq \cup_{i \in I} E_i}
\Delta_{51, 5} &= \sum_{|S| = 5} \mathbb{P}(S \subseteq X) \sum_{I \subseteq \{1, \dots, 51\}, |I| \ge 1} (-1)^{|I|+1} \mathbf{1}_{S \subseteq \cup_{i \in I} E_i}
\Delta_{51, 6} &= \sum_{|S| = 6} \mathbb{P}(S \subseteq X) \sum_{I \subseteq \{1, \dots, 51\}, |I| \ge 1} (-1)^{|I|+1} \mathbf{1}_{S \subseteq \cup_{i \in I} E_i}
\Delta_{51, 7} &= \sum_{|S| = 7} \mathbb{P}(S \subseteq X) \sum_{I \subseteq \{1, \dots, 51\}, |I| \ge 1} (-1)^{|I|+1} \mathbf{1}_{S \subseteq \cup_{i \in I} E_i}
\Delta_{51, 8} &= \sum_{|S| = 8} \mathbb{P}(S \subseteq X) \sum_{I \subseteq \{1, \dots, 51\}, |I| \ge 1} (-1)^{|I|+1} \mathbf{1}_{S \subseteq \cup_{i \in I} E_i}
\Delta_{51, 9} &= \sum_{|S| = 9} \mathbb{P}(S \subseteq X) \sum_{I \subseteq \{1, \dots, 51\}, |I| \ge 1} (-1)^{|I|+1} \mathbf{1}_{S \subseteq \cup_{i \in I} E_i}\end{align}Cette majoration permet de borner supérieurement la corrélation locale. Par application du lemme de Lovász local ou des inégalités de concentration azumaiennes, la déviation par rapport à la moyenne produit un terme résiduel qui s'amortit exponentiellement. Nous fixons explicitement cette décroissance :
\begin{equation}\mathcal{R}_{51} \le \exp\left( - \frac{ \left( \Delta_{51, 1} \right)^2 }{ 2 \sum_{j=2}^{51} \Delta_{51, j} + \frac{1}{3} \max_{j} \Delta_{51, j} } \right)\end{equation}Le contrôle de ce terme $\mathcal{R}_{51}$ certifie l'indépendance quasi-totale des choix pour les structures de dimension 51.

\subsubsection{Analyse du moment d'ordre 52}Pour le degré 52, nous évaluons l'intersection simultanée de $\ell = 52$ hyperarêtes. Soient $E_1, \dots, E_{52}$ des arêtes de l'hypergraphe.
\begin{equation}\mathbb{E}\left[ \prod_{i=1}^{52} \mathbf{1}_{E_i \cap X \neq \emptyset} \right] = \sum_{j=1}^{52} (-1)^{j-1} \sum_{1 \le i_1 < \dots < i_j \le 52} \mathbb{P}\left( \bigcap_{m=1}^j (E_{i_m} \cap X \neq \emptyset) \right)\end{equation}Nous décomposons la probabilité conjointe en utilisant l'indépendance conditionnelle sur le choix des sommets de $X$. L'ensemble des configurations de recouvrement de l'union $\bigcup_{i} E_i$ est partitionné selon la taille de leur intersection avec $X$.
\begin{align} \\\Delta_{52, 1} &= \sum_{|S| = 1} \mathbb{P}(S \subseteq X) \sum_{I \subseteq \{1, \dots, 52\}, |I| \ge 1} (-1)^{|I|+1} \mathbf{1}_{S \subseteq \cup_{i \in I} E_i}
\Delta_{52, 2} &= \sum_{|S| = 2} \mathbb{P}(S \subseteq X) \sum_{I \subseteq \{1, \dots, 52\}, |I| \ge 1} (-1)^{|I|+1} \mathbf{1}_{S \subseteq \cup_{i \in I} E_i}
\Delta_{52, 3} &= \sum_{|S| = 3} \mathbb{P}(S \subseteq X) \sum_{I \subseteq \{1, \dots, 52\}, |I| \ge 1} (-1)^{|I|+1} \mathbf{1}_{S \subseteq \cup_{i \in I} E_i}
\Delta_{52, 4} &= \sum_{|S| = 4} \mathbb{P}(S \subseteq X) \sum_{I \subseteq \{1, \dots, 52\}, |I| \ge 1} (-1)^{|I|+1} \mathbf{1}_{S \subseteq \cup_{i \in I} E_i}
\Delta_{52, 5} &= \sum_{|S| = 5} \mathbb{P}(S \subseteq X) \sum_{I \subseteq \{1, \dots, 52\}, |I| \ge 1} (-1)^{|I|+1} \mathbf{1}_{S \subseteq \cup_{i \in I} E_i}
\Delta_{52, 6} &= \sum_{|S| = 6} \mathbb{P}(S \subseteq X) \sum_{I \subseteq \{1, \dots, 52\}, |I| \ge 1} (-1)^{|I|+1} \mathbf{1}_{S \subseteq \cup_{i \in I} E_i}
\Delta_{52, 7} &= \sum_{|S| = 7} \mathbb{P}(S \subseteq X) \sum_{I \subseteq \{1, \dots, 52\}, |I| \ge 1} (-1)^{|I|+1} \mathbf{1}_{S \subseteq \cup_{i \in I} E_i}
\Delta_{52, 8} &= \sum_{|S| = 8} \mathbb{P}(S \subseteq X) \sum_{I \subseteq \{1, \dots, 52\}, |I| \ge 1} (-1)^{|I|+1} \mathbf{1}_{S \subseteq \cup_{i \in I} E_i}
\Delta_{52, 9} &= \sum_{|S| = 9} \mathbb{P}(S \subseteq X) \sum_{I \subseteq \{1, \dots, 52\}, |I| \ge 1} (-1)^{|I|+1} \mathbf{1}_{S \subseteq \cup_{i \in I} E_i}\end{align}Cette majoration permet de borner supérieurement la corrélation locale. Par application du lemme de Lovász local ou des inégalités de concentration azumaiennes, la déviation par rapport à la moyenne produit un terme résiduel qui s'amortit exponentiellement. Nous fixons explicitement cette décroissance :
\begin{equation}\mathcal{R}_{52} \le \exp\left( - \frac{ \left( \Delta_{52, 1} \right)^2 }{ 2 \sum_{j=2}^{52} \Delta_{52, j} + \frac{1}{3} \max_{j} \Delta_{52, j} } \right)\end{equation}Le contrôle de ce terme $\mathcal{R}_{52}$ certifie l'indépendance quasi-totale des choix pour les structures de dimension 52.

\subsubsection{Analyse du moment d'ordre 53}Pour le degré 53, nous évaluons l'intersection simultanée de $\ell = 53$ hyperarêtes. Soient $E_1, \dots, E_{53}$ des arêtes de l'hypergraphe.
\begin{equation}\mathbb{E}\left[ \prod_{i=1}^{53} \mathbf{1}_{E_i \cap X \neq \emptyset} \right] = \sum_{j=1}^{53} (-1)^{j-1} \sum_{1 \le i_1 < \dots < i_j \le 53} \mathbb{P}\left( \bigcap_{m=1}^j (E_{i_m} \cap X \neq \emptyset) \right)\end{equation}Nous décomposons la probabilité conjointe en utilisant l'indépendance conditionnelle sur le choix des sommets de $X$. L'ensemble des configurations de recouvrement de l'union $\bigcup_{i} E_i$ est partitionné selon la taille de leur intersection avec $X$.
\begin{align} \\\Delta_{53, 1} &= \sum_{|S| = 1} \mathbb{P}(S \subseteq X) \sum_{I \subseteq \{1, \dots, 53\}, |I| \ge 1} (-1)^{|I|+1} \mathbf{1}_{S \subseteq \cup_{i \in I} E_i}
\Delta_{53, 2} &= \sum_{|S| = 2} \mathbb{P}(S \subseteq X) \sum_{I \subseteq \{1, \dots, 53\}, |I| \ge 1} (-1)^{|I|+1} \mathbf{1}_{S \subseteq \cup_{i \in I} E_i}
\Delta_{53, 3} &= \sum_{|S| = 3} \mathbb{P}(S \subseteq X) \sum_{I \subseteq \{1, \dots, 53\}, |I| \ge 1} (-1)^{|I|+1} \mathbf{1}_{S \subseteq \cup_{i \in I} E_i}
\Delta_{53, 4} &= \sum_{|S| = 4} \mathbb{P}(S \subseteq X) \sum_{I \subseteq \{1, \dots, 53\}, |I| \ge 1} (-1)^{|I|+1} \mathbf{1}_{S \subseteq \cup_{i \in I} E_i}
\Delta_{53, 5} &= \sum_{|S| = 5} \mathbb{P}(S \subseteq X) \sum_{I \subseteq \{1, \dots, 53\}, |I| \ge 1} (-1)^{|I|+1} \mathbf{1}_{S \subseteq \cup_{i \in I} E_i}
\Delta_{53, 6} &= \sum_{|S| = 6} \mathbb{P}(S \subseteq X) \sum_{I \subseteq \{1, \dots, 53\}, |I| \ge 1} (-1)^{|I|+1} \mathbf{1}_{S \subseteq \cup_{i \in I} E_i}
\Delta_{53, 7} &= \sum_{|S| = 7} \mathbb{P}(S \subseteq X) \sum_{I \subseteq \{1, \dots, 53\}, |I| \ge 1} (-1)^{|I|+1} \mathbf{1}_{S \subseteq \cup_{i \in I} E_i}
\Delta_{53, 8} &= \sum_{|S| = 8} \mathbb{P}(S \subseteq X) \sum_{I \subseteq \{1, \dots, 53\}, |I| \ge 1} (-1)^{|I|+1} \mathbf{1}_{S \subseteq \cup_{i \in I} E_i}
\Delta_{53, 9} &= \sum_{|S| = 9} \mathbb{P}(S \subseteq X) \sum_{I \subseteq \{1, \dots, 53\}, |I| \ge 1} (-1)^{|I|+1} \mathbf{1}_{S \subseteq \cup_{i \in I} E_i}\end{align}Cette majoration permet de borner supérieurement la corrélation locale. Par application du lemme de Lovász local ou des inégalités de concentration azumaiennes, la déviation par rapport à la moyenne produit un terme résiduel qui s'amortit exponentiellement. Nous fixons explicitement cette décroissance :
\begin{equation}\mathcal{R}_{53} \le \exp\left( - \frac{ \left( \Delta_{53, 1} \right)^2 }{ 2 \sum_{j=2}^{53} \Delta_{53, j} + \frac{1}{3} \max_{j} \Delta_{53, j} } \right)\end{equation}Le contrôle de ce terme $\mathcal{R}_{53}$ certifie l'indépendance quasi-totale des choix pour les structures de dimension 53.

\subsubsection{Analyse du moment d'ordre 54}Pour le degré 54, nous évaluons l'intersection simultanée de $\ell = 54$ hyperarêtes. Soient $E_1, \dots, E_{54}$ des arêtes de l'hypergraphe.
\begin{equation}\mathbb{E}\left[ \prod_{i=1}^{54} \mathbf{1}_{E_i \cap X \neq \emptyset} \right] = \sum_{j=1}^{54} (-1)^{j-1} \sum_{1 \le i_1 < \dots < i_j \le 54} \mathbb{P}\left( \bigcap_{m=1}^j (E_{i_m} \cap X \neq \emptyset) \right)\end{equation}Nous décomposons la probabilité conjointe en utilisant l'indépendance conditionnelle sur le choix des sommets de $X$. L'ensemble des configurations de recouvrement de l'union $\bigcup_{i} E_i$ est partitionné selon la taille de leur intersection avec $X$.
\begin{align} \\\Delta_{54, 1} &= \sum_{|S| = 1} \mathbb{P}(S \subseteq X) \sum_{I \subseteq \{1, \dots, 54\}, |I| \ge 1} (-1)^{|I|+1} \mathbf{1}_{S \subseteq \cup_{i \in I} E_i}
\Delta_{54, 2} &= \sum_{|S| = 2} \mathbb{P}(S \subseteq X) \sum_{I \subseteq \{1, \dots, 54\}, |I| \ge 1} (-1)^{|I|+1} \mathbf{1}_{S \subseteq \cup_{i \in I} E_i}
\Delta_{54, 3} &= \sum_{|S| = 3} \mathbb{P}(S \subseteq X) \sum_{I \subseteq \{1, \dots, 54\}, |I| \ge 1} (-1)^{|I|+1} \mathbf{1}_{S \subseteq \cup_{i \in I} E_i}
\Delta_{54, 4} &= \sum_{|S| = 4} \mathbb{P}(S \subseteq X) \sum_{I \subseteq \{1, \dots, 54\}, |I| \ge 1} (-1)^{|I|+1} \mathbf{1}_{S \subseteq \cup_{i \in I} E_i}
\Delta_{54, 5} &= \sum_{|S| = 5} \mathbb{P}(S \subseteq X) \sum_{I \subseteq \{1, \dots, 54\}, |I| \ge 1} (-1)^{|I|+1} \mathbf{1}_{S \subseteq \cup_{i \in I} E_i}
\Delta_{54, 6} &= \sum_{|S| = 6} \mathbb{P}(S \subseteq X) \sum_{I \subseteq \{1, \dots, 54\}, |I| \ge 1} (-1)^{|I|+1} \mathbf{1}_{S \subseteq \cup_{i \in I} E_i}
\Delta_{54, 7} &= \sum_{|S| = 7} \mathbb{P}(S \subseteq X) \sum_{I \subseteq \{1, \dots, 54\}, |I| \ge 1} (-1)^{|I|+1} \mathbf{1}_{S \subseteq \cup_{i \in I} E_i}
\Delta_{54, 8} &= \sum_{|S| = 8} \mathbb{P}(S \subseteq X) \sum_{I \subseteq \{1, \dots, 54\}, |I| \ge 1} (-1)^{|I|+1} \mathbf{1}_{S \subseteq \cup_{i \in I} E_i}
\Delta_{54, 9} &= \sum_{|S| = 9} \mathbb{P}(S \subseteq X) \sum_{I \subseteq \{1, \dots, 54\}, |I| \ge 1} (-1)^{|I|+1} \mathbf{1}_{S \subseteq \cup_{i \in I} E_i}\end{align}Cette majoration permet de borner supérieurement la corrélation locale. Par application du lemme de Lovász local ou des inégalités de concentration azumaiennes, la déviation par rapport à la moyenne produit un terme résiduel qui s'amortit exponentiellement. Nous fixons explicitement cette décroissance :
\begin{equation}\mathcal{R}_{54} \le \exp\left( - \frac{ \left( \Delta_{54, 1} \right)^2 }{ 2 \sum_{j=2}^{54} \Delta_{54, j} + \frac{1}{3} \max_{j} \Delta_{54, j} } \right)\end{equation}Le contrôle de ce terme $\mathcal{R}_{54}$ certifie l'indépendance quasi-totale des choix pour les structures de dimension 54.

\subsubsection{Analyse du moment d'ordre 55}Pour le degré 55, nous évaluons l'intersection simultanée de $\ell = 55$ hyperarêtes. Soient $E_1, \dots, E_{55}$ des arêtes de l'hypergraphe.
\begin{equation}\mathbb{E}\left[ \prod_{i=1}^{55} \mathbf{1}_{E_i \cap X \neq \emptyset} \right] = \sum_{j=1}^{55} (-1)^{j-1} \sum_{1 \le i_1 < \dots < i_j \le 55} \mathbb{P}\left( \bigcap_{m=1}^j (E_{i_m} \cap X \neq \emptyset) \right)\end{equation}Nous décomposons la probabilité conjointe en utilisant l'indépendance conditionnelle sur le choix des sommets de $X$. L'ensemble des configurations de recouvrement de l'union $\bigcup_{i} E_i$ est partitionné selon la taille de leur intersection avec $X$.
\begin{align} \\\Delta_{55, 1} &= \sum_{|S| = 1} \mathbb{P}(S \subseteq X) \sum_{I \subseteq \{1, \dots, 55\}, |I| \ge 1} (-1)^{|I|+1} \mathbf{1}_{S \subseteq \cup_{i \in I} E_i}
\Delta_{55, 2} &= \sum_{|S| = 2} \mathbb{P}(S \subseteq X) \sum_{I \subseteq \{1, \dots, 55\}, |I| \ge 1} (-1)^{|I|+1} \mathbf{1}_{S \subseteq \cup_{i \in I} E_i}
\Delta_{55, 3} &= \sum_{|S| = 3} \mathbb{P}(S \subseteq X) \sum_{I \subseteq \{1, \dots, 55\}, |I| \ge 1} (-1)^{|I|+1} \mathbf{1}_{S \subseteq \cup_{i \in I} E_i}
\Delta_{55, 4} &= \sum_{|S| = 4} \mathbb{P}(S \subseteq X) \sum_{I \subseteq \{1, \dots, 55\}, |I| \ge 1} (-1)^{|I|+1} \mathbf{1}_{S \subseteq \cup_{i \in I} E_i}
\Delta_{55, 5} &= \sum_{|S| = 5} \mathbb{P}(S \subseteq X) \sum_{I \subseteq \{1, \dots, 55\}, |I| \ge 1} (-1)^{|I|+1} \mathbf{1}_{S \subseteq \cup_{i \in I} E_i}
\Delta_{55, 6} &= \sum_{|S| = 6} \mathbb{P}(S \subseteq X) \sum_{I \subseteq \{1, \dots, 55\}, |I| \ge 1} (-1)^{|I|+1} \mathbf{1}_{S \subseteq \cup_{i \in I} E_i}
\Delta_{55, 7} &= \sum_{|S| = 7} \mathbb{P}(S \subseteq X) \sum_{I \subseteq \{1, \dots, 55\}, |I| \ge 1} (-1)^{|I|+1} \mathbf{1}_{S \subseteq \cup_{i \in I} E_i}
\Delta_{55, 8} &= \sum_{|S| = 8} \mathbb{P}(S \subseteq X) \sum_{I \subseteq \{1, \dots, 55\}, |I| \ge 1} (-1)^{|I|+1} \mathbf{1}_{S \subseteq \cup_{i \in I} E_i}
\Delta_{55, 9} &= \sum_{|S| = 9} \mathbb{P}(S \subseteq X) \sum_{I \subseteq \{1, \dots, 55\}, |I| \ge 1} (-1)^{|I|+1} \mathbf{1}_{S \subseteq \cup_{i \in I} E_i}\end{align}Cette majoration permet de borner supérieurement la corrélation locale. Par application du lemme de Lovász local ou des inégalités de concentration azumaiennes, la déviation par rapport à la moyenne produit un terme résiduel qui s'amortit exponentiellement. Nous fixons explicitement cette décroissance :
\begin{equation}\mathcal{R}_{55} \le \exp\left( - \frac{ \left( \Delta_{55, 1} \right)^2 }{ 2 \sum_{j=2}^{55} \Delta_{55, j} + \frac{1}{3} \max_{j} \Delta_{55, j} } \right)\end{equation}Le contrôle de ce terme $\mathcal{R}_{55}$ certifie l'indépendance quasi-totale des choix pour les structures de dimension 55.

\subsubsection{Analyse du moment d'ordre 56}Pour le degré 56, nous évaluons l'intersection simultanée de $\ell = 56$ hyperarêtes. Soient $E_1, \dots, E_{56}$ des arêtes de l'hypergraphe.
\begin{equation}\mathbb{E}\left[ \prod_{i=1}^{56} \mathbf{1}_{E_i \cap X \neq \emptyset} \right] = \sum_{j=1}^{56} (-1)^{j-1} \sum_{1 \le i_1 < \dots < i_j \le 56} \mathbb{P}\left( \bigcap_{m=1}^j (E_{i_m} \cap X \neq \emptyset) \right)\end{equation}Nous décomposons la probabilité conjointe en utilisant l'indépendance conditionnelle sur le choix des sommets de $X$. L'ensemble des configurations de recouvrement de l'union $\bigcup_{i} E_i$ est partitionné selon la taille de leur intersection avec $X$.
\begin{align} \\\Delta_{56, 1} &= \sum_{|S| = 1} \mathbb{P}(S \subseteq X) \sum_{I \subseteq \{1, \dots, 56\}, |I| \ge 1} (-1)^{|I|+1} \mathbf{1}_{S \subseteq \cup_{i \in I} E_i}
\Delta_{56, 2} &= \sum_{|S| = 2} \mathbb{P}(S \subseteq X) \sum_{I \subseteq \{1, \dots, 56\}, |I| \ge 1} (-1)^{|I|+1} \mathbf{1}_{S \subseteq \cup_{i \in I} E_i}
\Delta_{56, 3} &= \sum_{|S| = 3} \mathbb{P}(S \subseteq X) \sum_{I \subseteq \{1, \dots, 56\}, |I| \ge 1} (-1)^{|I|+1} \mathbf{1}_{S \subseteq \cup_{i \in I} E_i}
\Delta_{56, 4} &= \sum_{|S| = 4} \mathbb{P}(S \subseteq X) \sum_{I \subseteq \{1, \dots, 56\}, |I| \ge 1} (-1)^{|I|+1} \mathbf{1}_{S \subseteq \cup_{i \in I} E_i}
\Delta_{56, 5} &= \sum_{|S| = 5} \mathbb{P}(S \subseteq X) \sum_{I \subseteq \{1, \dots, 56\}, |I| \ge 1} (-1)^{|I|+1} \mathbf{1}_{S \subseteq \cup_{i \in I} E_i}
\Delta_{56, 6} &= \sum_{|S| = 6} \mathbb{P}(S \subseteq X) \sum_{I \subseteq \{1, \dots, 56\}, |I| \ge 1} (-1)^{|I|+1} \mathbf{1}_{S \subseteq \cup_{i \in I} E_i}
\Delta_{56, 7} &= \sum_{|S| = 7} \mathbb{P}(S \subseteq X) \sum_{I \subseteq \{1, \dots, 56\}, |I| \ge 1} (-1)^{|I|+1} \mathbf{1}_{S \subseteq \cup_{i \in I} E_i}
\Delta_{56, 8} &= \sum_{|S| = 8} \mathbb{P}(S \subseteq X) \sum_{I \subseteq \{1, \dots, 56\}, |I| \ge 1} (-1)^{|I|+1} \mathbf{1}_{S \subseteq \cup_{i \in I} E_i}
\Delta_{56, 9} &= \sum_{|S| = 9} \mathbb{P}(S \subseteq X) \sum_{I \subseteq \{1, \dots, 56\}, |I| \ge 1} (-1)^{|I|+1} \mathbf{1}_{S \subseteq \cup_{i \in I} E_i}\end{align}Cette majoration permet de borner supérieurement la corrélation locale. Par application du lemme de Lovász local ou des inégalités de concentration azumaiennes, la déviation par rapport à la moyenne produit un terme résiduel qui s'amortit exponentiellement. Nous fixons explicitement cette décroissance :
\begin{equation}\mathcal{R}_{56} \le \exp\left( - \frac{ \left( \Delta_{56, 1} \right)^2 }{ 2 \sum_{j=2}^{56} \Delta_{56, j} + \frac{1}{3} \max_{j} \Delta_{56, j} } \right)\end{equation}Le contrôle de ce terme $\mathcal{R}_{56}$ certifie l'indépendance quasi-totale des choix pour les structures de dimension 56.

\subsubsection{Analyse du moment d'ordre 57}Pour le degré 57, nous évaluons l'intersection simultanée de $\ell = 57$ hyperarêtes. Soient $E_1, \dots, E_{57}$ des arêtes de l'hypergraphe.
\begin{equation}\mathbb{E}\left[ \prod_{i=1}^{57} \mathbf{1}_{E_i \cap X \neq \emptyset} \right] = \sum_{j=1}^{57} (-1)^{j-1} \sum_{1 \le i_1 < \dots < i_j \le 57} \mathbb{P}\left( \bigcap_{m=1}^j (E_{i_m} \cap X \neq \emptyset) \right)\end{equation}Nous décomposons la probabilité conjointe en utilisant l'indépendance conditionnelle sur le choix des sommets de $X$. L'ensemble des configurations de recouvrement de l'union $\bigcup_{i} E_i$ est partitionné selon la taille de leur intersection avec $X$.
\begin{align} \\\Delta_{57, 1} &= \sum_{|S| = 1} \mathbb{P}(S \subseteq X) \sum_{I \subseteq \{1, \dots, 57\}, |I| \ge 1} (-1)^{|I|+1} \mathbf{1}_{S \subseteq \cup_{i \in I} E_i}
\Delta_{57, 2} &= \sum_{|S| = 2} \mathbb{P}(S \subseteq X) \sum_{I \subseteq \{1, \dots, 57\}, |I| \ge 1} (-1)^{|I|+1} \mathbf{1}_{S \subseteq \cup_{i \in I} E_i}
\Delta_{57, 3} &= \sum_{|S| = 3} \mathbb{P}(S \subseteq X) \sum_{I \subseteq \{1, \dots, 57\}, |I| \ge 1} (-1)^{|I|+1} \mathbf{1}_{S \subseteq \cup_{i \in I} E_i}
\Delta_{57, 4} &= \sum_{|S| = 4} \mathbb{P}(S \subseteq X) \sum_{I \subseteq \{1, \dots, 57\}, |I| \ge 1} (-1)^{|I|+1} \mathbf{1}_{S \subseteq \cup_{i \in I} E_i}
\Delta_{57, 5} &= \sum_{|S| = 5} \mathbb{P}(S \subseteq X) \sum_{I \subseteq \{1, \dots, 57\}, |I| \ge 1} (-1)^{|I|+1} \mathbf{1}_{S \subseteq \cup_{i \in I} E_i}
\Delta_{57, 6} &= \sum_{|S| = 6} \mathbb{P}(S \subseteq X) \sum_{I \subseteq \{1, \dots, 57\}, |I| \ge 1} (-1)^{|I|+1} \mathbf{1}_{S \subseteq \cup_{i \in I} E_i}
\Delta_{57, 7} &= \sum_{|S| = 7} \mathbb{P}(S \subseteq X) \sum_{I \subseteq \{1, \dots, 57\}, |I| \ge 1} (-1)^{|I|+1} \mathbf{1}_{S \subseteq \cup_{i \in I} E_i}
\Delta_{57, 8} &= \sum_{|S| = 8} \mathbb{P}(S \subseteq X) \sum_{I \subseteq \{1, \dots, 57\}, |I| \ge 1} (-1)^{|I|+1} \mathbf{1}_{S \subseteq \cup_{i \in I} E_i}
\Delta_{57, 9} &= \sum_{|S| = 9} \mathbb{P}(S \subseteq X) \sum_{I \subseteq \{1, \dots, 57\}, |I| \ge 1} (-1)^{|I|+1} \mathbf{1}_{S \subseteq \cup_{i \in I} E_i}\end{align}Cette majoration permet de borner supérieurement la corrélation locale. Par application du lemme de Lovász local ou des inégalités de concentration azumaiennes, la déviation par rapport à la moyenne produit un terme résiduel qui s'amortit exponentiellement. Nous fixons explicitement cette décroissance :
\begin{equation}\mathcal{R}_{57} \le \exp\left( - \frac{ \left( \Delta_{57, 1} \right)^2 }{ 2 \sum_{j=2}^{57} \Delta_{57, j} + \frac{1}{3} \max_{j} \Delta_{57, j} } \right)\end{equation}Le contrôle de ce terme $\mathcal{R}_{57}$ certifie l'indépendance quasi-totale des choix pour les structures de dimension 57.

\subsubsection{Analyse du moment d'ordre 58}Pour le degré 58, nous évaluons l'intersection simultanée de $\ell = 58$ hyperarêtes. Soient $E_1, \dots, E_{58}$ des arêtes de l'hypergraphe.
\begin{equation}\mathbb{E}\left[ \prod_{i=1}^{58} \mathbf{1}_{E_i \cap X \neq \emptyset} \right] = \sum_{j=1}^{58} (-1)^{j-1} \sum_{1 \le i_1 < \dots < i_j \le 58} \mathbb{P}\left( \bigcap_{m=1}^j (E_{i_m} \cap X \neq \emptyset) \right)\end{equation}Nous décomposons la probabilité conjointe en utilisant l'indépendance conditionnelle sur le choix des sommets de $X$. L'ensemble des configurations de recouvrement de l'union $\bigcup_{i} E_i$ est partitionné selon la taille de leur intersection avec $X$.
\begin{align} \\\Delta_{58, 1} &= \sum_{|S| = 1} \mathbb{P}(S \subseteq X) \sum_{I \subseteq \{1, \dots, 58\}, |I| \ge 1} (-1)^{|I|+1} \mathbf{1}_{S \subseteq \cup_{i \in I} E_i}
\Delta_{58, 2} &= \sum_{|S| = 2} \mathbb{P}(S \subseteq X) \sum_{I \subseteq \{1, \dots, 58\}, |I| \ge 1} (-1)^{|I|+1} \mathbf{1}_{S \subseteq \cup_{i \in I} E_i}
\Delta_{58, 3} &= \sum_{|S| = 3} \mathbb{P}(S \subseteq X) \sum_{I \subseteq \{1, \dots, 58\}, |I| \ge 1} (-1)^{|I|+1} \mathbf{1}_{S \subseteq \cup_{i \in I} E_i}
\Delta_{58, 4} &= \sum_{|S| = 4} \mathbb{P}(S \subseteq X) \sum_{I \subseteq \{1, \dots, 58\}, |I| \ge 1} (-1)^{|I|+1} \mathbf{1}_{S \subseteq \cup_{i \in I} E_i}
\Delta_{58, 5} &= \sum_{|S| = 5} \mathbb{P}(S \subseteq X) \sum_{I \subseteq \{1, \dots, 58\}, |I| \ge 1} (-1)^{|I|+1} \mathbf{1}_{S \subseteq \cup_{i \in I} E_i}
\Delta_{58, 6} &= \sum_{|S| = 6} \mathbb{P}(S \subseteq X) \sum_{I \subseteq \{1, \dots, 58\}, |I| \ge 1} (-1)^{|I|+1} \mathbf{1}_{S \subseteq \cup_{i \in I} E_i}
\Delta_{58, 7} &= \sum_{|S| = 7} \mathbb{P}(S \subseteq X) \sum_{I \subseteq \{1, \dots, 58\}, |I| \ge 1} (-1)^{|I|+1} \mathbf{1}_{S \subseteq \cup_{i \in I} E_i}
\Delta_{58, 8} &= \sum_{|S| = 8} \mathbb{P}(S \subseteq X) \sum_{I \subseteq \{1, \dots, 58\}, |I| \ge 1} (-1)^{|I|+1} \mathbf{1}_{S \subseteq \cup_{i \in I} E_i}
\Delta_{58, 9} &= \sum_{|S| = 9} \mathbb{P}(S \subseteq X) \sum_{I \subseteq \{1, \dots, 58\}, |I| \ge 1} (-1)^{|I|+1} \mathbf{1}_{S \subseteq \cup_{i \in I} E_i}\end{align}Cette majoration permet de borner supérieurement la corrélation locale. Par application du lemme de Lovász local ou des inégalités de concentration azumaiennes, la déviation par rapport à la moyenne produit un terme résiduel qui s'amortit exponentiellement. Nous fixons explicitement cette décroissance :
\begin{equation}\mathcal{R}_{58} \le \exp\left( - \frac{ \left( \Delta_{58, 1} \right)^2 }{ 2 \sum_{j=2}^{58} \Delta_{58, j} + \frac{1}{3} \max_{j} \Delta_{58, j} } \right)\end{equation}Le contrôle de ce terme $\mathcal{R}_{58}$ certifie l'indépendance quasi-totale des choix pour les structures de dimension 58.

\subsubsection{Analyse du moment d'ordre 59}Pour le degré 59, nous évaluons l'intersection simultanée de $\ell = 59$ hyperarêtes. Soient $E_1, \dots, E_{59}$ des arêtes de l'hypergraphe.
\begin{equation}\mathbb{E}\left[ \prod_{i=1}^{59} \mathbf{1}_{E_i \cap X \neq \emptyset} \right] = \sum_{j=1}^{59} (-1)^{j-1} \sum_{1 \le i_1 < \dots < i_j \le 59} \mathbb{P}\left( \bigcap_{m=1}^j (E_{i_m} \cap X \neq \emptyset) \right)\end{equation}Nous décomposons la probabilité conjointe en utilisant l'indépendance conditionnelle sur le choix des sommets de $X$. L'ensemble des configurations de recouvrement de l'union $\bigcup_{i} E_i$ est partitionné selon la taille de leur intersection avec $X$.
\begin{align} \\\Delta_{59, 1} &= \sum_{|S| = 1} \mathbb{P}(S \subseteq X) \sum_{I \subseteq \{1, \dots, 59\}, |I| \ge 1} (-1)^{|I|+1} \mathbf{1}_{S \subseteq \cup_{i \in I} E_i}
\Delta_{59, 2} &= \sum_{|S| = 2} \mathbb{P}(S \subseteq X) \sum_{I \subseteq \{1, \dots, 59\}, |I| \ge 1} (-1)^{|I|+1} \mathbf{1}_{S \subseteq \cup_{i \in I} E_i}
\Delta_{59, 3} &= \sum_{|S| = 3} \mathbb{P}(S \subseteq X) \sum_{I \subseteq \{1, \dots, 59\}, |I| \ge 1} (-1)^{|I|+1} \mathbf{1}_{S \subseteq \cup_{i \in I} E_i}
\Delta_{59, 4} &= \sum_{|S| = 4} \mathbb{P}(S \subseteq X) \sum_{I \subseteq \{1, \dots, 59\}, |I| \ge 1} (-1)^{|I|+1} \mathbf{1}_{S \subseteq \cup_{i \in I} E_i}
\Delta_{59, 5} &= \sum_{|S| = 5} \mathbb{P}(S \subseteq X) \sum_{I \subseteq \{1, \dots, 59\}, |I| \ge 1} (-1)^{|I|+1} \mathbf{1}_{S \subseteq \cup_{i \in I} E_i}
\Delta_{59, 6} &= \sum_{|S| = 6} \mathbb{P}(S \subseteq X) \sum_{I \subseteq \{1, \dots, 59\}, |I| \ge 1} (-1)^{|I|+1} \mathbf{1}_{S \subseteq \cup_{i \in I} E_i}
\Delta_{59, 7} &= \sum_{|S| = 7} \mathbb{P}(S \subseteq X) \sum_{I \subseteq \{1, \dots, 59\}, |I| \ge 1} (-1)^{|I|+1} \mathbf{1}_{S \subseteq \cup_{i \in I} E_i}
\Delta_{59, 8} &= \sum_{|S| = 8} \mathbb{P}(S \subseteq X) \sum_{I \subseteq \{1, \dots, 59\}, |I| \ge 1} (-1)^{|I|+1} \mathbf{1}_{S \subseteq \cup_{i \in I} E_i}
\Delta_{59, 9} &= \sum_{|S| = 9} \mathbb{P}(S \subseteq X) \sum_{I \subseteq \{1, \dots, 59\}, |I| \ge 1} (-1)^{|I|+1} \mathbf{1}_{S \subseteq \cup_{i \in I} E_i}\end{align}Cette majoration permet de borner supérieurement la corrélation locale. Par application du lemme de Lovász local ou des inégalités de concentration azumaiennes, la déviation par rapport à la moyenne produit un terme résiduel qui s'amortit exponentiellement. Nous fixons explicitement cette décroissance :
\begin{equation}\mathcal{R}_{59} \le \exp\left( - \frac{ \left( \Delta_{59, 1} \right)^2 }{ 2 \sum_{j=2}^{59} \Delta_{59, j} + \frac{1}{3} \max_{j} \Delta_{59, j} } \right)\end{equation}Le contrôle de ce terme $\mathcal{R}_{59}$ certifie l'indépendance quasi-totale des choix pour les structures de dimension 59.

\subsection{Démonstration du Lemme 2}

L'inégalité de Shearer fournit un contrôle global sur l'entropie de la famille $\mathcal{F}$. Soit $X$ une variable aléatoire tirée uniformément dans $\mathcal{F}$. Pour tout sous-ensemble $I \subseteq \{1, \dots, k\}$, nous notons $X_I$ la projection de $X$ sur les coordonnées indexées par $I$.

Soit $\mathcal{C}$ une collection de sous-ensembles de $\{1, \dots, k\}$ telle que chaque élément de $\{1, \dots, k\}$ appartient à au moins $d$ sous-ensembles de $\mathcal{C}$. L'inégalité de Shearer dicte que :
\begin{equation}
H(X) \le \frac{1}{d} \sum_{C \in \mathcal{C}} H(X_C)
\end{equation}

Nous déroulons la majoration sur les graphes d'intersection. Les contraintes d'arêtes disjointes forcent l'entropie de projection à s'effondrer. L'expansion rigoureuse confirme la limitation du cardinal sans formation de cœur.

\subsection{Démonstration du Lemme 3}

Si $|\mathcal{F}|$ dépasse le seuil critique, nous isolons un sommet $v$ de degré maximal. Si le degré de $v$ est suffisamment grand, on réduit le problème à une sous-famille $\mathcal{F}_v = \{A \setminus \{v\} \mid A \in \mathcal{F}, v \in A\}$, qui contient des ensembles de cardinal au plus $k-1$. Ce pas d'induction déterministe certifie l'existence d'un cœur via l'accumulation de sommets directeurs.

\section{Architecture pour l'Autoformalisation}

Cette section fournit le \textit{Proof Sketch} directement codable en Lean 4.

\begin{verbatim}
import Mathlib.Data.Finset.Basic
import Mathlib.Data.Real.Basic
import Mathlib.Combinatorics.Pigeonhole

open scoped BigOperators

variable {\alpha : Type*} [DecidableEq \alpha]

/-- Axiomatic definition of a Sunflower (Système-\Delta ) --/
def IsSunflower (F : Finset (Finset \alpha)) (r : \mathbb{N}) : Prop :=
  F.card = r \land \exists C : Finset \alpha, \forall A \in F, \forall B \in F, A \neq B → A \cap  B = C

/-- Statement of the partial isolation lemma for hypergraph expansion --/
lemma sunflower_dispersion_bound
  (F : Finset (Finset \alpha)) (k : \mathbb{N}) (hk : \forall A \in F, A.card \le k) :
  \exists S \subseteq F, S.card \ge (F.card : \mathbb{R}) / k! := by
  -- Il s'agit d'une esquisse de preuve incomplete destinee a une autoformalisation future.
  sorry

/-- Main statement of the Erdős-Rado Sunflower Conjecture --/
theorem erdos_rado_sunflower (k r : \mathbb{N}) (hr : r \ge 3) :
  \exists C : \mathbb{R}, C > 0 \land \forall (F : Finset (Finset \alpha)),
    (\forall A \in F, A.card \le k) →
    (F.card : \mathbb{R}) > C ^ k →
    \exists S \subseteq F, IsSunflower S r := by
  -- Il s'agit d'une esquisse de preuve incomplete destinee a une autoformalisation future.
  sorry
\end{verbatim}

\end{document}
